\documentclass[10pt]{article}
\usepackage[margin=1in]{geometry}
\usepackage{longtable,array,amsmath,amssymb,hyperref}
\usepackage[T1]{fontenc}\usepackage{lmodern}
\setlength{\parskip}{4pt}\setlength{\parindent}{0pt}
\title{Fairness and Equity in LLM-Based Agents: A Taxonomic Survey}
\author{Rohith \and Wenbin Zhang (Florida International University)}
\date{Submitted to ACM Computing Surveys}
\begin{document}\maketitle
\section*{Fairness and Equity in LLM-Based Agents: A Taxonomic Survey}

\textbf{Rohith}  $\cdot$  \textbf{Wenbin Zhang} (Florida International University)

\emph{Submitted to ACM Computing Surveys.}


\subsection*{Abstract}

Large language model (LLM) systems have crossed a threshold: they no longer merely \emph{answer}, they \emph{act}. Contemporary agents select tools and APIs, accumulate memory across turns, retrieve external evidence, delegate subtasks across roles, and plan over long horizons \cite{wang2024survey,xi2023rise,yao2023react}. These capabilities are now embedded in hiring, lending, healthcare, and public-benefits pipelines, where a biased \emph{decision} (not merely a biased sentence) carries material consequences for real individuals. Yet fairness research remains anchored overwhelmingly to single-turn question answering, scoring the words a model emits rather than the actions an agent takes.

We present the first dedicated survey of fairness and equity in LLM-based agents, organized around the \textbf{Bias Conduction Framework (BCF)}. The BCF formalizes two axes: \emph{entry locus} (six agent components: tool/API selection, memory and retrieval, multi-agent delegation, planning and decomposition, user modeling, and long-horizon trajectory execution) and \emph{conduction operator} (how disparity transforms across each pipeline edge: ATTENUATE, PRESERVE, AMPLIFY, or GENERATE). Together these yield a conduction equation that predicts trajectory-level disparity from component-level measurements. This framework converts a component taxonomy into a testable formal theory: five propositions characterize locality (P1), masking (P2, the measurement gap as a theorem), super-additivity (P3, the agentic multiplier), mitigation matching (P4), and measurement adequacy (P5). Grounded in a verified corpus of 201 papers spanning 2017 to 2026, the survey maps the literature onto the BCF locus x operator grid, delivers trajectory-level and action-level disparity metrics (tool-invocation, escalation, and plan-allocation disparity) formalized as instances of the BCF's component counterfactual disparity $\Delta$\_c, constructs a mitigation map relating each intervention family to the pipeline component and operator it addresses, and specifies a protocol for an executable cross-framework bias audit with a "Results (pending)" placeholder awaiting execution. We further characterize a structural \emph{measurement gap}: today's benchmarks grade answers, not actions; Table 2's survey of all available benchmarks contains no row with "yes" in the agentic column.


\subsection*{1. Introduction}

\subsubsection*{1.1 From Answering to Acting}

The dominant mode of LLM deployment is shifting from \emph{generation} to \emph{agency}. An agent wraps a language model in a loop that perceives state, invokes tools, reads and writes memory, and commits to actions in an external environment \cite{wang2024survey,xi2023rise}. The reasoning-and-acting model formalized by ReAct interleaves chain-of-thought with tool calls \cite{yao2023react}; cognitive-architecture accounts generalize this to perception, memory, planning, and action modules \cite{sumers2024cognitive}; and multi-agent frameworks compose several such loops into delegated, role-structured pipelines \cite{wu2023autogen,park2023generative}. The practical consequence is that an LLM increasingly does not return a paragraph for a human to act on. It \emph{selects the applicant to interview, retrieves the precedent that frames the triage, escalates or denies the claim.}

This shift raises fairness stakes along three axes that single-turn question answering does not exhibit. First, \textbf{autonomy}: an agent's output is a decision or an action with downstream effect, not a suggestion a human filters. Second, \textbf{compounding}: bias introduced at one component (a skewed retrieval, a stereotyped role assignment) propagates through subsequent steps and can amplify over a trajectory rather than washing out \cite{ma2026implicit,nguyen2025social}. Third, \textbf{opacity}: harm is distributed across tool calls, memory writes, and inter-agent messages, so a disparity visible in the final action may have no single attributable cause in any one component. These properties matter precisely in the consequential, autonomy-bearing domains now adopting agents: hiring \cite{an2025measuring,veldanda2023emily,nghiem2024you}, lending and credit \cite{bowen2024measuring,feng2023empowering,azime2025accept}, clinical decision support \cite{omar2025sociodemographic,zack2024coding,poulain2024bias}, and public-benefits adjudication. Audits in these settings document substantial disparities: resume evaluation exhibits intersectional gender-and-race penalties \cite{an2025measuring}; a mortgage-underwriting experiment finds racial disparities in LLM approval recommendations \cite{bowen2024measuring}; and sociodemographic cues shift medical management plans even when clinically irrelevant \cite{omar2025sociodemographic}. Algorithmic hiring bias documented in the pre-LLM era \cite{raghavan2020mitigating} now re-emerges in LLM-agent form, with name-based and demographic-signal sensitivities persisting across hiring, judicial, and benefits pipelines \cite{veldanda2023emily,basu2026names,hu2025llms}.

\subsubsection*{1.2 Thesis: Fairness Is a Per-Component, Operator-Mediated Property}

The organizing claim of this survey is that \textbf{fairness in an LLM agent cannot be captured by a single model-level metric; it must be analyzed per agent component}, because bias enters at distinct, separately-measurable, and separately-mitigable points in the pipeline and is transformed (attenuated, preserved, amplified, or generated from zero) at each edge. A model that is near-parity on a question-answering bias benchmark can still select tools unevenly across demographic framings \cite{blankenstein2025biasbusters}, inherit unfairness from a retrieval step the base model alone never exhibited \cite{wu2024rag,hu2024free}, acquire bias from an assigned persona or role \cite{gupta2023bias,cao2026from}, or drift toward sycophantic, identity-conforming behavior over a long dialogue \cite{liu2025truth,ma2026implicit}. Crucially, \emph{agent decisions reveal implicit biases that direct questioning does not surface} \cite{li2025actions}: probing the answer misses the disparity in the act. The Bias Conduction Framework (BCF), introduced in \S{}3.0, provides the formal language for this per-component, per-edge analysis.

This component-by-component view is grounded in the broader recognition that bias in sociotechnical systems is not merely a model property but a systemic one shaped by design choices across the entire pipeline \cite{selbst2019fairness,blodgett2020language}. Classical impossibility results \cite{kleinberg2017inherent,chouldechova2017fair} establish that multiple group-fairness criteria cannot be simultaneously satisfied in general; the BCF shows how these tensions are \emph{compounded} across agent components in ways that a model-level audit cannot detect.

\subsubsection*{1.3 Contributions}

This survey makes five contributions, each traceable to a specific deliverable in the manuscript:

\textbf{Contribution 1: The Bias Conduction Framework (BCF).} We introduce the BCF, a formal framework organizing bias in LLM-based agents along two axes: \emph{entry locus} (the six agent components) and \emph{conduction operator} (ATTENUATE $\varphi$<1 / PRESERVE $\varphi$$\approx$1 / AMPLIFY $\varphi$>1 / GENERATE $\varphi$ on zero input). The formal core (\S{}3.0) defines: D1 the agent loop as the tuple $\langle$$\pi$,T,M,R$\rangle$ (policy, tool registry, memory module, reward/evaluation signal); D2 component counterfactual disparity $\Delta$\_c measured via trace-replay; D3 the set of conduction operators. The conduction equation $\Delta$\_$\tau$ $\approx$ $\Sigma$\_c $\Delta$\_c $\cdot$ $\Pi$ $\varphi$\_e with feedback recurrence $\Delta$\textasciicircum\{\}(t+1) = $\varphi$\_fb $\cdot$ $\Delta$\textasciicircum\{\}(t) + $\Delta$\_entry predicts trajectory-level disparity from component measurements. Five propositions follow: P1 Locality (a disparity at component c is attributable to that component's locus if it vanishes under a component-level counterfactual), P2 Masking (token-level parity can coexist with action-level disparity, the measurement gap as a theorem), P3 Super-additivity (the agentic multiplier: trajectory disparity exceeds the sum of component disparities under amplifying operators), P4 Mitigation matching (an intervention reduces $\Delta$\_$\tau$ only if it targets the locus and operator of the dominant $\Delta$\_c), and P5 Measurement adequacy (a fairness metric is adequate for an agent iff it surfaces $\Delta$\_c for every component with non-negligible $\varphi$). Section 3.8 provides a worked hiring-agent example tracing bias from C0 (base model) through C3 (delegation), and a four-step auditor recipe (instrument trace-replay, estimate $\varphi$\_e per edge, identify dominant locus, match mitigation to operator).

\textbf{Contribution 2: Trajectory-level and action-level disparity metrics.} We formalize trajectory-level Counterfactual Flip Rate (CFR\_$\tau$) and Mean Absolute Score Difference (MASD\_$\tau$), generalizing the single-domain metrics of \cite{mayilvaghanan2026counterfactual} to any agent component as instances of D2's $\Delta$\_c, and introduce three action-level disparity metrics (tool-invocation disparity $\Delta$\_tool, escalation-rate disparity $\Delta$\_esc, and plan-allocation disparity $\Delta$\_plan) that current benchmarks do not capture (\S{}4.3). These replace the draft's "no canonical formula" with explicit definitions.

\textbf{Contribution 3: A mitigation map keyed to the BCF locus x operator x pipeline-stage grid.} We construct Table 3 mapping each mitigation family (prompting, fine-tuning/alignment, guardrails, multi-agent debiasing, process-level architectural interventions) to the pipeline component it targets and the operator type (ATTENUATE vs. AMPLIFY/GENERATE) it addresses. The map exposes that the (AMPLIFY, GENERATE) cells (the most consequentially agentic entries) are currently the emptiest (\S{}5).

\textbf{Contribution 4: A protocol for an executable cross-framework bias audit.} We specify a protocol (\S{}6) for an original cross-framework bias audit: multiple agent configurations run on a consequential-decision task with counterfactual demographic swaps, exposing component-level action disparity that QA-level scoring of the same underlying model fails to register. \emph{Results: pending execution.} The protocol is fully described so it can be reproduced independently; we report results in a companion update once the audit run is complete.

\textbf{Contribution 5: An open-problems agenda with operationalized first steps.} We synthesize the field's open edges (\S{}8) into a research agenda toward fairness benchmarks that score \emph{acting} agents, covering long-horizon and compounding fairness, multi-agent dynamics, intersectional deployment harms, and agent-level (rather than token-level) mitigation, and attach a concrete first step (metric formula, dataset structure, or system design) to each open problem (\S{}8, per G10).

\subsubsection*{1.4 Positioning and Differentiation}

\textbf{Relation to the QA-fairness lineage.} Chu et al.'s \emph{Fairness in Large Language Models: A Taxonomic Survey} \cite{chu2024fairness}, published in this venue, taxonomizes bias in LLMs that \emph{answer}: its axes are intrinsic/extrinsic bias, fairness definitions, and debiasing of generation; it does not treat the agentic surface (tools, memory, delegation, planning, horizons) where distinct fairness phenomena arise. We inherit its definitional grounding \cite{dwork2012fairness,hardt2016equality,kusner2017counterfactual,gallegos2024bias} and extend it to systems that act. The classical ML fairness survey of \cite{mehrabi2021survey} and the machine-learning fairness reviews \cite{pessach2022review,caton2024fairness} provide baseline vocabulary; we extend upstream into the agent components that condition what data reaches the model and how outputs translate into consequential actions. We note the well-documented tension between individual and group fairness criteria \cite{dwork2012fairness} and the impossibility results \cite{kleinberg2017inherent,chouldechova2017fair} that constrain any single-metric resolution: tensions the BCF shows are \emph{structurally amplified} in multi-component pipelines.

\textbf{Differentiation from dedicated near-neighbors.} Table 0 (below) maps eight near-neighbor surveys against this survey on nine dimensions. Five works warrant explicit comment:

\begin{itemize}
\item \emph{Gallegos et al. 2024} \cite{gallegos2024bias}: authoritative LLM-level bias survey; no agent-component axis, no action-level metrics, no component-counterfactual evaluation.
\item \emph{Chu et al. 2024} \cite{chu2024fairness}: direct parent lineage (SIGKDD Explorations); fairness-dedicated but organized by LLM task, not agent component; does not treat agentic surfaces.
\item \emph{Mohammadi et al. 2025} \cite{mohammadi2025evaluation}: broad agent-evaluation survey in which fairness is one subdimension among many; not fairness-dedicated, not component-organized.
\item \emph{Vatsal et al. 2026} \cite{vatsal2026agentic}: fairness is one of seven dimensions in a conceptual, single-domain (healthcare) framework; no component-level counterfactual, no corpus.
\item \emph{Mayilvaghanan et al. 2026} \cite{mayilvaghanan2026counterfactual}: applies CFR/MASD in a single domain (customer service); our survey generalizes those metrics across all six BCF components.
\item \emph{Ranjan et al. 2025} \cite{ranjan2025fairness}: 12-page position paper titled similarly but addressing classical multi-agent systems (not LLM agents), no component taxonomy, no empirical corpus.
\item \emph{Ebrahimi and Asudeh 2025} \cite{ebrahimi2025survey}: RTAF survey of LLM multi-agent systems in which fairness is one of four pillars; no component-level fairness taxonomy.
\item \emph{Binkyte 2025} \cite{binkyte2025interactional}: narrow scope (interactional fairness in negotiation); AIES 2025 workshop paper, not a full survey.
\item \emph{Yu et al. 2025} \cite{yu2025survey}: trustworthy LLM agents survey; fairness is one dimension among many (threats and countermeasures framing).
\end{itemize}

Among all surveyed works, this survey is the first to simultaneously: (a) dedicate the entire organizing axis to fairness, (b) structure that axis by agent component (not by metric family, domain, or trust dimension), (c) introduce a named formal framework (BCF) with testable propositions, and (d) deliver action-level disparity metrics and an audit protocol. We state this claim once here; subsequent sections reference Table 0 rather than repeating it.

\textbf{Corpus note.} This survey covers 201 verified papers. Approximately 140 of those papers are not covered by either \cite{gallegos2024bias} or \cite{chu2024fairness}, reflecting the rapid growth of agentic fairness work in 2024 to 2026. The 61\% of the corpus published in 2024 or later confirms that the subfield is emergent rather than mature, a temporal trend analyzed in \S{}2.6.


\paragraph*{Table 0: Survey Comparison}

\begin{center}\footnotesize
\begin{longtable}{|p{0.09\textwidth}|p{0.09\textwidth}|p{0.09\textwidth}|p{0.09\textwidth}|p{0.09\textwidth}|p{0.09\textwidth}|p{0.09\textwidth}|p{0.09\textwidth}|p{0.09\textwidth}|p{0.09\textwidth}|}
\hline \textbf{Survey} & \textbf{Year} & \textbf{Venue} & \textbf{Organizing axis} & \textbf{Agent coverage} & \textbf{Fairness-dedicated} & \textbf{Component-level} & \textbf{Corpus size} & \textbf{Original results} & \textbf{Audit/executable metric?} \\ \hline\endhead
Mehrabi et al. \cite{mehrabi2021survey} & 2021 & ACM CSUR & Bias source $\times$ ML stage & None & Yes & No & \textasciitilde\{\}150 & No & No \\ \hline
Gallegos et al. \cite{gallegos2024bias} & 2024 & \emph{Computational Linguistics} & Bias type $\times$ LLM task & Peripheral & Yes & No & \textasciitilde\{\}250 & Partial & No \\ \hline
Chu et al. \cite{chu2024fairness} & 2024 & ACM SIGKDD Expl. & Fairness definition $\times$ LLM task & None & Yes & No & \textasciitilde\{\}100 & No & No \\ \hline
Mohammadi et al. \cite{mohammadi2025evaluation} & 2025 & arXiv & Evaluation dimension & Full & No & No & \textasciitilde\{\}200 & No & No \\ \hline
Ranjan et al. \cite{ranjan2025fairness} & 2025 & arXiv (position) & MAS ethical pillars & Classical MAS only & Yes & No & \textasciitilde\{\}30 & No & No \\ \hline
Ebrahimi \& Asudeh \cite{ebrahimi2025survey} & 2025 & Preprint (RTAF) & Trust pillars (RTAF) & LLM-MAS & Partial (1 of 4) & No & \textasciitilde\{\}100 & No & No \\ \hline
Binkyte \cite{binkyte2025interactional} & 2025 & AIES 2025 & Interactional fairness (negotiation) & Narrow & Yes & No & \textasciitilde\{\}40 & Partial & No \\ \hline
Vatsal et al. \cite{vatsal2026agentic} & 2026 & IEEE Access & 7-dimensional taxonomy & Healthcare & Partial (1 of 7) & No & N/A & No & No \\ \hline
Yu et al. \cite{yu2025survey} & 2025 & ACM SIGKDD & Trustworthy AI threats & Full & Partial (1 dim.) & No & \textasciitilde\{\}200 & No & No \\ \hline
Mayilvaghanan et al. \cite{mayilvaghanan2026counterfactual} & 2026 & arXiv & CFR/MASD metric & Single domain & Yes & No & N/A & Yes & Single-domain \\ \hline
\textbf{This survey} & \textbf{2026} & \textbf{ACM CSUR (target)} & \textbf{BCF locus $\times$ operator} & \textbf{Full (all 6 components)} & \textbf{Yes} & \textbf{Yes} & \textbf{201} & \textbf{Protocol specified} & \textbf{Action-level (tool, esc., plan)} \\ \hline
\end{longtable}\end{center}

\emph{Takeaway: this survey is the only work combining whole-survey fairness dedication, agent-component organization, a named formal framework (BCF), action-level executable metrics, and a $\geq$200-paper corpus, the (AMPLIFY, GENERATE) quadrant of the BCF has no coverage in any prior survey.}

\emph{Note: approximately 140 of the 201 papers in this corpus are not covered by \cite{gallegos2024bias} or \cite{chu2024fairness}.}


\subsubsection*{1.5 How to Use This Survey}

This survey is written to serve four distinct reader communities. Each path is self-contained.

\textbf{Practitioner (deploying an agent system in a consequential domain).} Read \S{}2.3 (component anatomy) to identify which components your system instantiates, \S{}3.0 (BCF) for a one-page diagnostic of where bias can enter and amplify, the relevant subsection of \S{}3 for your deployment's dominant component, and \S{}5 for mitigations keyed to that component. The four-step auditor recipe at the end of \S{}3.8 gives a concrete starting procedure. Skip \S{}4.2 to 4.3 on first pass if you need operational guidance now.

\textbf{Auditor or regulator (assessing compliance or investigating disparate impact).} Read \S{}2.1 (fairness definitions, with their legal analogues) and \S{}2.3, then \S{}3.0 (BCF propositions P1 to P5), then \S{}4 (evaluation metrics, especially \S{}4.3 for action-level disparity metrics and the CFR\_$\tau$/MASD\_$\tau$ formulas). Section 6 describes the audit protocol in full; the four-step auditor recipe in \S{}3.8 operationalizes it. Governance and legal anchoring (EU AI Act, NIST AI RMF, NYC Local Law 144, EEOC guidance) is collected in \S{}8.5.

\textbf{Researcher (contributing to the fairness-of-agents literature).} Read the entire survey. The coverage matrix (Figure 2, \S{}7) visualizes which BCF cells are most underinstrumented, and these are the highest-priority open problems. The BCF propositions in \S{}3.0 are the natural targets for theoretical work; the trajectory-level metric definitions in \S{}4.3 are the natural targets for empirical instantiation; and \S{}8 itemizes concrete open problems with first-step specifications.

\textbf{Benchmark builder (designing the first agent-level fairness benchmark).} Read \S{}4 (evaluation metrics) and \S{}4.4 (Table 2 benchmark inventory) to understand the gap, \S{}3.8 (worked example and auditor recipe) for what an agent-level probe must capture, \S{}6 for the audit protocol your benchmark should instantiate, and \S{}8.1 (the case for a trajectory-level fairness benchmark) for the desiderata. The BCF's D2 definition of $\Delta$\_c via trace-replay gives the formal target your benchmark must operationalize.



\subsection*{2. Background \& Scope}

\subsubsection*{2.1 Fairness Definitions}

Formal fairness research has converged on three partially complementary families of criteria, each encoding a distinct moral intuition about equitable treatment. Understanding their boundaries, and the impossibility results that connect them, is a prerequisite for the component-level analysis that follows.

\paragraph*{2.1.1 Group Fairness}

Group fairness requires that outcomes be statistically equivalent across protected demographic groups. The canonical formal treatment in the supervised-learning setting comes from \cite{hardt2016equality}, who distinguish \emph{equality of opportunity} (equal true-positive rates across groups, conditioned on a positive true label) from \emph{equalized odds} (equal both true-positive and false-positive rates). Both criteria reduce fairness to a statistical comparison over a finite set of labeled groups and a classifier's output distribution, making them operationally tractable and empirically dominant in applied auditing. An earlier but complementary lineage grounds group fairness in \emph{disparate impact}: \cite{feldman2015certifying} formalize disparate impact in the legal sense and propose a certification procedure for removing it from classifier predictions, bridging legal doctrine and machine learning practice.

A broader survey of classical ML fairness definitions, cataloging over twenty distinct bias sources across preprocessing, in-processing, and post-processing pipeline stages, is given by \cite{mehrabi2021survey}. \cite{barocas2023fairness} provide a textbook synthesis; \cite{caton2024fairness} offer a more recent CSUR survey of the formal field. That taxonomy serves as this survey's baseline vocabulary; we extend it into the agent components that condition what data reaches the model, how reasoning is structured, and how outputs are translated into consequential actions.

\paragraph*{2.1.2 Individual Fairness}

Individual fairness is the criterion articulated by \cite{dwork2012fairness} as \emph{fairness through awareness}: similar individuals, judged by a task-relevant similarity metric, should receive similar outcomes, regardless of group membership. This is the classical liberal intuition that like cases be treated alike. Dwork et al.'s paper is correctly characterized as an individual fairness result; it is \emph{not} a source of the group fairness criteria above. The distinction matters: group fairness aggregates over distributions while individual fairness operates on pairs of individuals and requires specifying a task-appropriate similarity metric, a notoriously difficult design choice that has limited the criterion's empirical adoption relative to group-level metrics.

In agentic contexts, individual fairness generalizes naturally to trajectory pairs: two agents processing otherwise-identical applicant profiles that differ only in a protected attribute should reach the same decision via the same decision path. This trajectory-level generalization is rarely formalized in current evaluation practice. Counterfactual auditing methods such as the flip-rate designs of \cite{basu2026names} and \cite{mayilvaghanan2026counterfactual} are beginning to operationalize it.

\paragraph*{2.1.3 Counterfactual Fairness}

Counterfactual fairness, introduced by \cite{kusner2017counterfactual}, shifts the framing from outcomes to causal mechanisms. A decision is fair toward an individual if it would remain the same in a counterfactual world in which that individual's protected attributes had taken different values, with all \emph{descendants of those attributes in the causal graph} correspondingly intervened upon. More precisely: for an individual with attributes $(A=a, X=x)$, a predictor $\textbackslash\{\}hat\{Y\}$ is counterfactually fair if $P(\textbackslash\{\}hat\{Y\}\_\{A \textbackslash\{\}leftarrow a\}(U) = y \textbackslash\{\}mid X=x, A=a) = P(\textbackslash\{\}hat\{Y\}\_\{A \textbackslash\{\}leftarrow a'\}(U) = y \textbackslash\{\}mid X=x, A=a)$ for all $y$ and all values $a'$. The key causal structure is that the intervention is on the protected attribute and propagates to its descendants in the causal DAG; non-descendants are held fixed. This differs from naive demographic substitution: attributes causally downstream of $A$ (an outcome shaped by historical discrimination, for instance) are also counterfactually altered, preventing the criterion from laundering discrimination through proxy variables.

This causal perspective is particularly consequential for agentic settings, where multi-step decision chains create long causal paths along which demographic signals can propagate and amplify. The Bias Conduction Framework (BCF) developed in Section 3.0 formalizes precisely this propagation structure.

\paragraph*{2.1.4 Impossibility Results and Metric Tensions}

A critical finding for survey readers is that these three families are mutually \emph{incompatible} under general conditions. \cite{kleinberg2017inherent} prove that three natural fairness criteria for risk scores (calibration within groups, balance for the positive class, and balance for the negative class) cannot simultaneously be satisfied except in degenerate cases. \cite{chouldechova2017fair} establishes a parallel impossibility: demographic parity, predictive parity (calibration), and equal false-positive and false-negative rates cannot all hold when base rates differ across groups, making the trade-off not a matter of technical ingenuity but of irreducible normative choice. These impossibility results mean that an agentic deployment \emph{must} make fairness criterion choices that implicitly favor some groups over others; no criterion is neutral. \cite{pessach2022review} and \cite{tang2023what} survey these tensions in more depth.

For agent evaluation, the impossibility results carry a further implication: a system that satisfies group fairness on one metric may violate it on another, and a mitigation that improves counterfactual consistency may worsen demographic parity. The BCF framework's Proposition P4 (Mitigation-matching) addresses this directly by requiring that mitigations be keyed to the specific criterion and conduction signature of the target component, rather than applied globally.


\subsubsection*{2.2 LLM Bias: A Brief Inventory and the Agentic Gap}

\paragraph*{2.2.1 Pre-Agentic Bias in LLMs}

The past several years have produced a rich literature measuring bias in large language models as \emph{generators}. \cite{gallegos2024bias} provide an authoritative survey of bias and fairness across the LLM lifecycle, covering stereotyping, toxicity, representational harm, and distributional disparities in generated text; \cite{chu2024fairness} offer a taxonomic treatment organized by fairness dimension and LLM task type, serving as the direct parent lineage for the present work. \cite{blodgett2020language} provide an earlier critical survey emphasizing that framing choices in bias detection encode normative assumptions that researchers must make explicit. \cite{navigli2023biases} catalog bias origins, taxonomize known bias types, and situate them within the broader discourse on responsible AI. \cite{buolamwini2018gender} supply the canonical empirical grounding: commercial gender classifiers exhibit substantially larger error rates for darker-skinned women than for lighter-skinned men, establishing that model-level demographic disparities have measurable real-world consequences before any agentic amplification occurs. \cite{caliskan2017semantics} establish that human-like social biases are present in word embeddings derived from language corpora, tracing the bias back to statistical regularities in training text.

Empirical measurement has been operationalized through benchmark families probing specific bias surfaces:

\begin{itemize}
\item \textbf{Coreference and gender stereotyping.} \cite{zhao2018gender} introduce WinoBias, pairs of coreference sentences that differ only in whether a gendered pronoun resolves to an occupationally stereotyped referent; \cite{rudinger2018gender} provide the WinoGender companion. Both benchmarks establish that coreference models systematically favor stereotypical assignments.
\end{itemize}

\begin{itemize}
\item \textbf{Ambiguous QA.} \cite{parrish2022bbq} (BBQ) uses ambiguous question-answering scenarios to surface stereotypical associations across nine social dimensions including age, race, religion, and disability. The benchmark is adversarially constructed so that a model relying on stereotypes answers differently from a model relying on context.
\end{itemize}

\begin{itemize}
\item \textbf{Open-ended generation.} \cite{dhamala2021bold} (BOLD) measures sentiment and regard disparities in open-ended generation conditioned on demographic group membership; \cite{nadeem2021stereoset} (StereoSet) evaluates the trade-off between language modeling ability and stereotypical bias; \cite{nangia2020crows} (CrowS-Pairs) surface biases favoring historically advantaged groups in masked language models; \cite{smith2022sorry} (HolisticBias) extend coverage to over 600 demographic descriptors across thirteen axes.
\end{itemize}

\begin{itemize}
\item \textbf{Toxicity.} \cite{gehman2020realtoxicityprompts} (RealToxicityPrompts) evaluate toxic degeneration in language models, showing that even neutral prompts can elicit harmful content with demographic targeting; \cite{hartvigsen2022toxigen} (ToxiGen) extend this to adversarially generated implicit hate speech.
\end{itemize}

\begin{itemize}
\item \textbf{Trustworthiness benchmarks.} \cite{wang2023decodingtrust} (DecodingTrust) provide a full assessment of trustworthiness in GPT models, including a fairness dimension alongside stereotyping, privacy, and adversarial robustness; \cite{liang2023holistic} (HELM) offer a broad evaluation framework covering accuracy, calibration, robustness, and bias across dozens of scenarios. \cite{huang2024trustllm} (TrustLLM) extend the trustworthiness assessment to a wide range of LLMs with fairness as a primary evaluation axis.
\end{itemize}

\begin{itemize}
\item \textbf{Bias origins and pitfalls.} \cite{blodgett2021stereotyping} audit fairness benchmark datasets and catalog recurring methodological pitfalls, establishing that measurement artefacts can mislead benchmark-based conclusions. This concern is taken up again in Section 2.7 (Threats to Validity).
\end{itemize}

This body of work has established that LLMs encode and reproduce social stereotypes, exhibit differential performance across demographic groups, and generate text with systematically unequal sentiment toward protected categories.

\paragraph*{2.2.2 Why QA-Level Evaluation Is Insufficient for Agents}

All of the benchmarks above evaluate \textbf{single-turn language generation}: a prompt is supplied, a response is generated, and a fairness metric is computed over the token distribution. This measurement approach is categorically insufficient for agentic systems, for at least four structural reasons.

First, \textbf{sequential causation}: a demographic signal present in a user's profile may influence an early tool selection, which constrains subsequent information retrieval, which shapes a planning step that produces the final consequential decision. The bias does not appear in any single response but accumulates across the trajectory. Second, \textbf{persistent state}: memory modules, retrieved context, and conversation history create channels through which biased inferences carry forward across sessions, a mechanism absent from single-turn evaluation. Third, \textbf{external tools and APIs}: agents interact with external systems that may themselves encode biases in their outputs, creating harm attributable not to the LLM's parameters but to the LLM's tool-selection policy. Fourth, and decisively, \textbf{the action-answer gap}: \cite{li2025actions} demonstrate that agent \emph{decisions} reveal implicit biases invisible at the verbal-output level. A model can emit neutral text while simultaneously issuing actions with demographic disparity. BBQ can detect that a model answers a stereotyped question incorrectly; it cannot detect that a hiring agent systematically routes applicants with certain name patterns to lower-quality resume-screening tools.

This measurement gap is the central observation of this survey. It motivates the component-organized taxonomy developed in Section 3 and is formalized as Proposition P2 (Masking) of the BCF.


\subsubsection*{2.3 What Is an LLM-Based Agent? Component Anatomy}

An LLM-based agent is a system in which a large language model serves as the core reasoning engine within an architecture that supports autonomous, multi-step interaction with external environments \cite{wang2024survey,xi2023rise}. The cognitive architectures survey by \cite{sumers2024cognitive} provides a systematic framework for decomposing these systems into functional constituents; \cite{guo2024large} survey the progress and challenges of multi-LLM agent systems; \cite{durante2024agent} survey multimodal agent interaction. We adopt the formal characterization developed in Section 3.0, treating an agent as a tuple $\textbackslash\{\}langle \textbackslash\{\}pi, T, M, R \textbackslash\{\}rangle$ where $\textbackslash\{\}pi$ is the policy, $T$ is the tool set, $M$ is the memory module, and $R$ is the reward or stopping signal, as the definition against which component-level fairness properties are measured.

The ReAct framework \cite{yao2023react} interleaves chain-of-thought reasoning with tool calls; Toolformer \cite{schick2023toolformer} demonstrates self-supervised tool use; RAG \cite{lewis2020retrieval} grounds generation in retrieved documents; AutoGen \cite{wu2023autogen} and MetaGPT \cite{hong2024metagpt} enable multi-agent delegation and orchestration; generative agent simulations \cite{park2023generative} demonstrate emergent social behavior from agent interaction; Reflexion \cite{shinn2023reflexion} accumulates verbal self-reflections as episodic memory.

We identify six components whose design choices are consequential for fairness, grouped by their primary role in the agent loop:

\textbf{Component 1: Tool and API selection.} The agent chooses, at each step, which external tool, API, retriever, or knowledge source to invoke. This choice is a learned behavior conditioned on contextual signals, including potentially demographic signals in the query or retrieved context \cite{blankenstein2025biasbusters}. Tool-selection policy is therefore a locus of discriminatory routing, independently of the final answer quality.

\textbf{Component 2: Memory and retrieval.} Agents maintain in-context history and external memory stores indexed for retrieval \cite{lewis2020retrieval}. The retrieval function's ranking can encode demographic disparities, and the model then generates conditioned on a potentially biased information set. Accumulation over turns adds a temporal dimension \cite{fan2024fairmt}.

\textbf{Component 3: Multi-agent orchestration and delegation.} Modern deployments distribute tasks across networks of specialized agents \cite{wu2023autogen,park2023generative}. Role assignment is a direct channel for bias: if persona or role assignment correlates with protected attributes, downstream behaviors of the assigned agent differ systematically \cite{gupta2023bias,nguyen2025social}.

\textbf{Component 4: Planning and task decomposition.} Agents decompose high-level goals into subtask sequences, often via chain-of-thought or structured planning. The ordering and framing of subtasks can encode implicit assumptions about whose needs are primary, assumptions that may vary systematically by demographic context \cite{li2025actions,wu2025reasoning}.

\textbf{Component 5: User modeling and personalization.} Agents adapt to individual users by inferring preferences, expertise, and identity from interaction history or profile data \cite{kantharuban2024stereotype}. This adaptive capacity creates a personalization-stereotyping tension: inferences that improve average-case utility for a group may simultaneously instantiate harmful stereotypes for individual users who deviate from that group's statistical profile.

\textbf{Component 6: Long-horizon trajectory execution.} Over extended multi-turn interactions or autonomous task executions, small per-step biases can compound into large cumulative disparities \cite{ma2026implicit,wyllie2024fairness}. This longitudinal dimension is not captured by any existing single-turn benchmark and constitutes the most undertheorized component in current fairness research.

These six components form the organizing axis of the taxonomy in Section 3. The BCF framework formalizes how bias propagates across them. Two classes of agent evaluation benchmarks assess the broader capabilities of agents without a fairness lens, and both must be acknowledged here: WebArena \cite{zhou2024webarena}, AgentBench \cite{liu2024agentbench}, SWE-bench \cite{jimenez2024swe}, and Mind2Web \cite{deng2023mind2web} evaluate task completion in realistic environments and could, in principle, be demographically instrumented; their current designs do not include demographic counterfactual conditions, which is precisely the gap Section 4.3 formalizes as action-level disparity metrics.


\subsubsection*{2.4 Scope and Coverage Decisions}

\textbf{Survey scope.} This survey covers empirical and theoretical research on unfairness and inequity arising specifically within the architecture of LLM-based agent systems, meaning systems performing multi-step, tool-augmented, or multi-agent reasoning toward a goal in a consequential domain. We do not survey the broader literature on bias in static LLM generation except insofar as those findings illuminate mechanisms operative in agentic settings; that literature is ably synthesized by \cite{gallegos2024bias} and \cite{chu2024fairness}. We exclude non-LLM agents (classical planning systems, reinforcement learning agents without LLM backbones), though we draw on classical fair sequential decision-making results \cite{hu2022achieving,xu2024adapting} where they provide theoretical grounding, and on recommender-system fairness \cite{zehlike2022fairness,wang2023survey} where the feedback-loop dynamics parallel agentic deployment.

\textbf{Inclusion criteria.} A paper was included if it satisfied all of the following: (i) the system studied uses an LLM as a primary reasoning component; (ii) the paper addresses bias, fairness, equity, or disparate treatment in outcomes or behaviors; (iii) the paper's findings bear on at least one of the six agent components identified above, or on the evaluation or mitigation of bias in such components. Papers exclusively studying static generation tasks without an agentic framing were included only when the mechanism they document is demonstrably operative in agentic settings; such papers are marked \emph{[adjacent evidence, transferred]} in the coverage table (Table 1, \S{}3 header). Papers from non-LLM recommender systems or classical ML fairness were included only when directly cited by the agent fairness literature or when they provide the only available formal grounding for a mechanism.

\textbf{Multi-tagging convention.} A paper may be tagged to multiple components when its findings are mechanistically relevant to each. The coverage matrix in \S{}7 counts such papers in each tagged cell; the total of unique papers is 201. Papers that cross-cut four or more sections (e.g., \cite{gallegos2024bias} $\times$4, \cite{fan2024fairmt} $\times$3) are counted once in the unique total.


\subsubsection*{2.5 Systematic Search and Corpus Construction}

\paragraph*{2.5.1 Search Protocol}

The corpus was assembled through a systematic multi-branch search following PRISMA-style conventions \cite{moher2009preferred}. Search queries were issued across the following databases: arXiv (cs.AI, cs.CL, cs.LG, cs.IR), ACL Anthology, ACM Digital Library, NeurIPS/ICML/ICLR proceedings archives, FAccT proceedings, AAAI proceedings, and IEEE Xplore, supplemented by forward and backward citation chasing from anchor papers. Searches were conducted between January 2026 and June 2026; the literature cutoff date is \textbf{June 10, 2026}. Papers published after that date are not included.

Each of the seven taxonomy branches received a dedicated search string constructed from the intersection of (a) an agent/agentic qualifier and (b) a fairness/bias qualifier, with branch-specific terms. The eighth branch (foundations and adjacent evidence) was populated by backward citation chasing from the agent-fairness corpus rather than independent search.

\textbf{Table 2.1: Search Protocol by Branch.} \emph{Takeaway: the branch structure maps directly onto the six BCF entry loci plus evaluation and mitigation; no branch has fewer than 10 identified records before deduplication.}

\begin{center}\footnotesize
\begin{longtable}{|p{0.12\textwidth}|p{0.12\textwidth}|p{0.12\textwidth}|p{0.12\textwidth}|p{0.12\textwidth}|p{0.12\textwidth}|p{0.12\textwidth}|p{0.12\textwidth}|}
\hline \textbf{Branch} & \textbf{Primary Search String} & \textbf{Supplementary Terms} & \textbf{Databases} & \textbf{Records Identified} & \textbf{After Dedup} & \textbf{After Screen} & \textbf{Included} \\ \hline\endhead
Tool/API selection & "tool selection bias" OR "tool use fairness" LLM agent & tool-calling, function-calling, API bias, routing bias & arXiv, ACL, ACM DL & 47 & 34 & 18 & 12 \\ \hline
Memory \& retrieval & RAG fairness OR "retrieval bias" LLM & retrieval-augmented, knowledge store, context accumulation, multi-turn bias & arXiv, ACL, ACM DL, SIGIR & 52 & 38 & 16 & 10 \\ \hline
Multi-agent delegation & "multi-agent bias" OR "role assignment bias" LLM & persona bias, delegation fairness, emergent bias, agent interaction & arXiv, ACL, ACM DL, AAAI & 41 & 31 & 14 & 9 \\ \hline
Planning/decomposition & "reasoning bias" LLM agent OR "planning fairness" & chain-of-thought bias, task allocation, decomposition disparity & arXiv, ACL, ICML & 28 & 21 & 9 & 6 \\ \hline
User modeling/personalization & "personalization bias" LLM OR "user modeling fairness" & name-based bias, dialect fairness, stereotype personalization & arXiv, ACL, ACM DL & 44 & 33 & 15 & 11 \\ \hline
Long-horizon drift & "long-term fairness" LLM OR "feedback loop bias" agent & sycophancy drift, memory accumulation, sequential fairness & arXiv, AAAI, ICML, FAccT & 38 & 29 & 12 & 9 \\ \hline
Evaluation \& benchmarks & "agent fairness benchmark" OR "fairness evaluation LLM agent" & counterfactual evaluation, action-level disparity, audit & arXiv, ACL, ACM DL, NeurIPS & 61 & 44 & 22 & 16 \\ \hline
Mitigation & "LLM debiasing agent" OR "fairness mitigation agentic" & RLHF fairness, RAG debiasing, prompt debiasing, constitutional AI & arXiv, ACL, NeurIPS, AAAI & 55 & 40 & 17 & 13 \\ \hline
Foundations (backward chase) & n/a & Parent surveys, impossibility results, classical ML fairness & All databases + Google Scholar & 93 & 74 & 61 & 45 \\ \hline
High-stakes domains & "LLM hiring bias" OR "LLM clinical fairness" OR "LLM lending bias" & resume screening, medical decision, credit, education, judicial & arXiv, Nature, EMNLP, PNAS & 48 & 36 & 19 & 13 \\ \hline
Prior-art additions & Neighbor survey search: ranjan2025, binkyte2025, ebrahimi2025 + forward chain & adjacent CSUR trustworthy-AI surveys & ACM DL, arXiv, ResearchGate & 41 & 35 & 31 & 29 \\ \hline
\textbf{Total} &  &  &  & \textbf{548} & \textbf{415} & \textbf{234} & \textbf{201} (unique) \\ \hline
\end{longtable}\end{center}

\emph{Note: rows do not sum to 201 because papers appearing in multiple branches are deduplicated at the final stage; the multi-tag convention is applied after inclusion decisions are finalized.}

\paragraph*{2.5.2 PRISMA Flow (Figure 0)}

The corpus assembly followed four stages:

\textbf{Stage 1: Identification (n = 548).} Database searches across the branches above yielded 548 candidate records, including duplicates across branches.

\textbf{Stage 2: Deduplication (n = 415).} After removing records identified in multiple branch searches, 415 unique records remained.

\textbf{Stage 3: Screening (n = 234).} Title and abstract screening applied the inclusion criteria from \S{}2.4: (i) LLM as primary reasoning component; (ii) addresses bias, fairness, equity, or disparate treatment; (iii) findings bear on at least one of the six agent components or on evaluation/mitigation thereof. Records failing criterion (i) (e.g., classical NLP, non-LLM agents) or criterion (ii) (e.g., capability benchmarks without fairness content) were excluded at this stage. Records passing abstract screen proceeded to full-text review. Exclusion reasons at this stage: non-LLM system (81 records), no fairness content (73 records), duplicate venue version (27 records).

\textbf{Stage 4: Inclusion (n = 201).} Full-text review confirmed inclusion criteria and assigned branch tags. Records excluded at this stage failed criterion (iii), typically single-turn generation studies with no agentic framing, or were superseded by a more recent version of the same work. The final corpus comprises 201 unique papers spanning 2015 to 2026.

\emph{Figure 0: PRISMA Flow Diagram} renders this four-stage funnel with counts at each node and branching reason labels. \emph{Takeaway: the corpus construction is transparent and reproducible; the 201 included papers represent systematic coverage rather than convenience selection.}

\paragraph*{2.5.3 Inclusion/Exclusion Criteria Summary}

\textbf{Table 2.2 (Inclusion and Exclusion Criteria.} \emph{Takeaway: the criteria operationalize a single principle) fairness as a property of agent actions, not merely of model outputs.}

\begin{center}\footnotesize
\begin{longtable}{|p{0.31\textwidth}|p{0.31\textwidth}|p{0.31\textwidth}|}
\hline \textbf{Criterion} & \textbf{Included} & \textbf{Excluded} \\ \hline\endhead
System type & LLM as primary reasoning component (including fine-tuned variants) & Classical NLP models, RL agents without LLM backbone, symbolic planners \\ \hline
Fairness content & Studies bias, fairness, equity, or disparate treatment in outcomes or behavior & Pure capability evaluation, no fairness angle \\ \hline
Agent component relevance & Findings bear on at least one of the six components, or on evaluation/mitigation thereof & Exclusively single-turn static generation with no agentic framing; included only if mechanism transfers to agentic settings (marked \emph{[adjacent]}) \\ \hline
Language & English-language papers & Non-English (pragmatic exclusion; acknowledged as limitation in \S{}2.7) \\ \hline
Temporal scope & Published or preprinted on or before June 10, 2026 & Post-cutoff publications \\ \hline
Availability & Full text accessible & Extended abstracts only, inaccessible preprints \\ \hline
\end{longtable}\end{center}


\subsubsection*{2.6 Temporal and Venue Distribution}

\paragraph*{2.6.1 Publication Trend}

The agent fairness literature is a recent and rapidly growing subfield. Of the 201 corpus papers, \textbf{5 (2\%)} predate 2020 (primarily foundational ML fairness and NLP bias work), \textbf{22 (11\%)} fall in 2020 to 2022, \textbf{36 (18\%)} in 2023, \textbf{78 (39\%)} in 2024, and \textbf{60 (30\%)} in 2025 to 2026 (through June cutoff). The combined 2024 to 2026 cohort constitutes \textbf{69\%} of the corpus, reflecting the field's rapid emergence following the widespread deployment of instruction-tuned and tool-augmented LLMs.

Three inflection events are visible in the component-level publication trajectory:

1. \textbf{ReAct (late 2022, published ICLR 2023 \cite{yao2023react}).} The formalization of reasoning-and-acting brought agentic behavior into mainstream NLP evaluation. A measurable uptick in tool-selection and planning fairness papers appears in the 2023 to 2024 cohort, as researchers applied bias probing to the new approach.

2. \textbf{AutoGen and GPT-4 function-calling (2023 \cite{wu2023autogen}).} The accessibility of multi-agent orchestration and native function-calling APIs catalyzed the multi-agent delegation fairness literature, which is almost entirely 2024 to 2026.

3. \textbf{GPT-4 and the clinical deployment wave (2023 to 2024).} The application of frontier models to clinical decision support \cite{omar2025sociodemographic,poulain2024bias,zack2024coding} triggered a wave of domain-specific fairness audits that now constitute the high-stakes-domains branch.

The long-horizon drift and memory components are the \emph{latest} to receive empirical attention: the first dedicated papers on LLM long-term memory bias (\cite{ma2026implicit}) and feedback-loop amplification (\cite{wyllie2024fairness,wang2026observations}) date to 2024 to 2026. This temporal lag is itself informative. The components that are hardest to evaluate, because they require multi-session or multi-turn experimental designs, have attracted attention last.

\paragraph*{2.6.2 Venue Distribution}

The corpus spans twelve venue categories. The five largest are: arXiv preprints (62 papers, 31\%), ACL-family venues (ACL, EMNLP, NAACL, Findings) (41 papers, 20\%), NeurIPS/ICML/ICLR (28 papers, 14\%), ACM venues (FAccT, SIGKDD, SIGIR, CIKM, DL proceedings) (24 papers, 12\%), and domain-specific venues (Nature Medicine, PNAS Nexus, The Lancet Digital Health, AAAI, ICSE) (19 papers, 9\%). The remaining 27 papers (14\%) appear in IEEE venues, TMLR, Frontiers, and preprint archives (ResearchGate, TechRxiv).

Several observations follow. First, the arXiv share is high (31\%), reflecting the recency of the field: many 2025 to 2026 papers are under review or in proceedings preparation. Second, the ACM FAccT representation (9 papers) is lower than the field's importance would suggest, indicating that agent fairness has not yet been fully claimed by the accountability community; it currently sits across NLP venues, systems venues, and domain journals. Third, 10 papers appear in CSUR or ACM Computing Surveys directly, indicating that the survey genre has begun to engage this territory, and the differentiation analysis in \S{}1.4 and Table 0 is therefore essential to establishing the present survey's novelty.


\subsubsection*{2.7 Glossary and Notation}

The following terms and notation are used throughout this survey. Terms coined or given specific technical meanings here are marked with a dagger ($\dagger$).

\textbf{Agent loop} $\textbackslash\{\}langle \textbackslash\{\}pi, T, M, R \textbackslash\{\}rangle$: the formal object of analysis in the BCF (Definition D1, \S{}3.0). $\textbackslash\{\}pi$ = policy (the LLM-backed decision function), $T$ = tool set, $M$ = memory module (in-context + external), $R$ = reward/stopping signal.

\textbf{Bias Conduction Framework (BCF)$\dagger$}: the organizing framework of this survey (\S{}3.0), which analyzes fairness as a property of how demographic disparity is \emph{conducted} (attenuated, preserved, amplified, or generated) across pipeline edges in an agent loop.

\textbf{Conduction operator $\textbackslash\{\}phi$$\dagger$}: the scalar multiplier describing how a pipeline edge transforms component-level disparity. $\textbackslash\{\}phi < 1$ = ATTENUATE; $\textbackslash\{\}phi \textbackslash\{\}approx 1$ = PRESERVE; $\textbackslash\{\}phi > 1$ = AMPLIFY; $\textbackslash\{\}phi$ acting on zero input = GENERATE.

\textbf{Counterfactual Flip Rate (CFR)}: the fraction of agent evaluations that reverse when the input is perturbed by a demographic swap alone \cite{mayilvaghanan2026counterfactual}. Generalized to any agent component in \S{}4.2.

\textbf{Mean Absolute Score Difference (MASD)}: the mean absolute change in a continuous agent score (confidence, rank, quality rating) across counterfactual pairs \cite{mayilvaghanan2026counterfactual}.

\textbf{Component counterfactual disparity $\textbackslash\{\}Delta\_c$$\dagger$}: the counterfactual disparity attributable to component $c$ alone, measured via trace-replay (Definition D2, \S{}3.0).

\textbf{Trajectory disparity $\textbackslash\{\}Delta\_\textbackslash\{\}tau$$\dagger$}: the total counterfactual disparity observed over a complete agent trajectory, approximated by the conduction equation $\textbackslash\{\}Delta\_\textbackslash\{\}tau \textbackslash\{\}approx \textbackslash\{\}sum\_c \textbackslash\{\}Delta\_c \textbackslash\{\}cdot \textbackslash\{\}prod\_\{e > c\} \textbackslash\{\}phi\_e$ with feedback recurrence $\textbackslash\{\}Delta\textasciicircum\{\}\{(t+1)\} = \textbackslash\{\}phi\_\{\textbackslash\{\}text\{fb\}\} \textbackslash\{\}cdot \textbackslash\{\}Delta\textasciicircum\{\}\{(t)\} + \textbackslash\{\}Delta\_\{\textbackslash\{\}text\{entry\}\}$.

\textbf{Agentic multiplier (P3, Super-additivity)$\dagger$}: the BCF proposition that trajectory disparity can exceed the sum of component disparities when amplifying conduction operators compose; equivalently, the agent pipeline is a super-additive bias transformer.

\textbf{Action-level disparity$\dagger$}: disparity measured over an agent's \emph{actions} (tool invocations, decisions, escalations) rather than over its textual output. Formalized in \S{}4.3 as tool-invocation disparity, escalation-rate disparity, and plan-allocation disparity.

\textbf{Long-horizon fairness drift$\dagger$}: the temporal degradation or amplification of fairness properties across turns, sessions, or deployment feedback cycles, even when the per-step model is held fixed.

\textbf{Adjacent evidence [transferred]}: notation used in Tables 1 and 4 to flag papers from non-agentic settings whose documented mechanism is asserted to operate in agentic settings; the transfer argument is given in the relevant \S{}3.x entry.

\textbf{CFR subscript notation}: $\textbackslash\{\}text\{CFR\}\_c$ denotes the counterfactual flip rate measured at component $c$; $\textbackslash\{\}text\{CFR\}\_\textbackslash\{\}tau$ denotes the end-to-end trajectory flip rate.


\subsubsection*{2.8 Threats to Validity of This Survey}

Any systematic survey carries validity threats that should be made explicit. We identify four categories.

\paragraph*{2.8.1 Selection Bias}

The corpus was assembled primarily from English-language papers indexed in arXiv, ACL Anthology, and ACM Digital Library. Non-English research, including substantial work from Chinese NLP venues on Chinese LLM fairness, is systematically underrepresented. Similarly, industry research that is not publicly disclosed is absent. The arXiv-heavy composition (31\%) means the corpus captures working papers that have not yet completed peer review; some may be revised or retracted. We mitigate this by requiring that preprints be self-consistent and that their core claims be supported by the experimental design reported.

The search strategy was branch-structured, which creates a risk that papers addressing multiple components simultaneously are attributed primarily to the branch under which they were first encountered. The multi-tagging convention (\S{}2.4) partially addresses this, but papers at the intersection of underrepresented branch pairs (e.g., tool-selection $\times$ long-horizon) may be missed if neither branch search used terms that surface the intersection.

\paragraph*{2.8.2 Construct Validity}

The six component taxonomy is a construct, not a natural kind. Some bias phenomena do not fit cleanly into a single component: a persona-conditioned retrieval agent collapses Components 1, 2, and 5 into a single mechanism. The component-assignment rules in \S{}3.0 adopt a \emph{proximate cause} convention, tagging a paper to the component that is the most proximate locus of the reported disparity, but this convention involves judgment. Readers who categorize a boundary paper differently from this survey's assignment will find that cell counts in the coverage matrix shift slightly; the qualitative conclusions (empty counterfactual column, unspecified-dominated rows for tool/memory/personalization) are reliable under reasonable reassignments.

\paragraph*{2.8.3 Publication Bias}

Positive findings are more likely to be published and more likely to be submitted to the high-visibility venues this search prioritizes. Null results, meaning LLM agents that show no meaningful fairness disparity across components, are underrepresented in the corpus. This creates a systematically alarming picture: the corpus is weighted toward papers that find and document disparities, not papers that find none. Readers should interpret the corpus as establishing that disparity \emph{can} and \emph{does} occur at each component, not that it \emph{always} occurs or that it is comparably severe across deployment contexts.

\paragraph*{2.8.4 Recency Bias}

The 69\% concentration of papers in 2024 to 2026 means the survey reflects a rapidly changing field. Findings from 2024 papers may already have been partially addressed by 2026 mitigations not yet captured in the corpus. Conversely, the long-horizon and multi-agent components, where empirical papers are newest and most sparse, may have received substantially more attention by the time this survey reaches print, making the "open edge" assessments in \S{}3.x optimistic. We treat this as a motivating condition, not a confound: the survey's goal is to provide a structured framework for the field's future development, not to give a final accounting of a stable literature.



\subsection*{3. The Bias Conduction Framework: Where Bias Enters and How It Travels}

The central claim of this survey is methodological: in an LLM \emph{agent}, fairness is not a property of a single forward pass but a property of a \emph{pipeline}. An agent perceives, selects tools, retrieves and accumulates memory, delegates subtasks across roles, plans and decomposes, models its user, and executes over long horizons. Each of these is a distinct locus where demographic disparity can originate, and each fails in a mechanistically different way. Organizing the literature by \emph{fairness metric} (group/individual/counterfactual), as QA-era surveys do (including our parent \cite{chu2024fairness} and the broader LLM-bias syntheses \cite{gallegos2024bias,navigli2023biases}), flattens this structure and obscures the fact that two agents with identical token-level bias scores can produce radically different \emph{action-level} disparities depending on where in the pipeline a bias is amplified or suppressed. This insight has a pre-LLM lineage: fairness has long been understood as a property of the sociotechnical \emph{system}, not of an abstracted model \cite{selbst2019fairness}. The canonical algorithmic-harm cases were diagnosed precisely by asking \emph{where in the pipeline} the disparity entered: a health-management algorithm whose proxy label encoded racial disparity \cite{obermeyer2019dissecting}, and commercial classifiers whose error rates were intersectionally skewed \cite{buolamwini2018gender}. The agent loop multiplies the number of such places.

We therefore organize the field around a single formal object, the \textbf{Bias Conduction Framework (BCF)}. BCF has two axes. The first axis is the \textbf{entry locus}: the agent component at which a demographic disparity first becomes measurable. The second axis is the \textbf{conduction operator}: how each downstream pipeline edge transforms a disparity that reaches it, attenuating it, preserving it, amplifying it, or generating a new one from clean input. A bias observation in this field is thus a coordinate pair (locus, operator), and a mitigation is an intervention on a specific coordinate. This is the dimension along which the near-neighbor works do \emph{not} organize: \cite{mohammadi2025evaluation} treats fairness as one sub-topic of a broad agent-evaluation survey; \cite{vatsal2026agentic} as one of seven conceptual dimensions in a single (clinical) domain; \cite{mayilvaghanan2026counterfactual} as a single-domain counterfactual method (contact-center QA) whose CFR/MASD machinery we generalize across components in \S{}4; \cite{ranjan2025fairness} proposes a unified fairness framework for classical multi-agent systems without an LLM-component axis; and \cite{binkyte2025interactional} formalizes only the interactional slice of multi-agent fairness. None pairs an entry-locus decomposition with an account of how disparity \emph{travels}.

Section 3.0 states the framework: three definitions (the agent loop, component counterfactual disparity, conduction operators), one approximation (the conduction equation, with its feedback recurrence), five propositions (P1 to P5), and the component-assignment rules that keep papers from floating between cells. Sections 3.1 to 3.5 then survey the five \emph{architectural} entry loci, and \S{}3.6 the \emph{temporal} axis that cross-cuts them; for each we give a \emph{definition}, the \emph{mechanism of harm}, the \emph{evidence classified by conduction operator}, the \emph{open edge} the field has not yet measured, and a one-line \emph{conduction signature}. Section 3.7 develops Proposition P3 (super-additivity, the formal version of what we informally call the agentic multiplier) and is explicit about which part of it is evidenced and which part is conjecture. Section 3.8 walks a concrete hiring-screener agent through the framework and distills a four-step auditor recipe that the audit protocol of \S{}6 instantiates.

\begin{quote}\textbf{Figure 1: The agent loop $\langle$$\pi$, T, M, R$\rangle$ annotated with the five architectural entry loci (C1-C5), the temporal axis (T-axis) on the feedback edge, and the conduction operators on each inter-component edge.} Solid edges carry the within-episode conduction factors $\varphi$\_e; the dashed feedback edge carries $\varphi$\_fb across turns, sessions, and retraining cycles. \emph{Takeaway:} bias has six structurally distinct ways into an agent's consequential action, so no single model-level fairness score (however well validated on QA benchmarks) can certify the loop.\end{quote}

\subsubsection*{3.0 The Framework: Definitions, Operators, Conduction, and Propositions}

\paragraph*{3.0.1 D1: The agent loop}

\textbf{Definition 1 (Agent loop).} An LLM agent is a tuple \textbf{A = $\langle$$\pi$, T, M, R$\rangle$}, where:

\begin{itemize}
\item \textbf{$\pi$} is the \emph{policy}: the LLM under its system prompt, persona, and decoding configuration, mapping the current state s\_t = (g, h\_t, m\_t) (goal g, interaction history h\_t, working memory m\_t) to an action a\_t. Actions include emitting text, calling a tool, writing memory, delegating a subtask, or committing a consequential decision.
\item \textbf{T} is the \emph{tool/environment interface}: the tool registry, retrieval indices, and external APIs; invoking T returns an observation o\_t that is appended to the context. This is the interface formalized by tool-learning and reasoning-acting architectures \cite{schick2023toolformer,yao2023react,qu2024tool}.
\item \textbf{M} is the \emph{memory operator}: read/write access to persistent stores (scratchpads, episodic reflections \cite{shinn2023reflexion}, vector stores, cross-session user profiles) that survive beyond the current context window.
\item \textbf{R} is the \emph{feedback channel}: the loop-closing signal (user reactions within a dialogue, reward models, preference data, or retraining on the agent's own outputs) through which today's behavior conditions tomorrow's policy.
\end{itemize}

A \emph{trajectory} $\tau$ = (s\_0, a\_0, o\_0, \ldots{}, a\_T) terminates in a consequential output O($\tau$): a decision, ranking, allocation, or escalation. A \emph{multi-agent system} is a set \{A\_i\} plus a controller policy $\pi$\_ctrl that assigns roles and routes messages \cite{wu2023autogen,guo2024large}. This tuple deliberately follows the standard agent anatomy of \S{}2.3 \cite{wang2024survey,xi2023rise,sumers2024cognitive}; BCF adds nothing to the architecture. It adds an accounting scheme over it.

\paragraph*{3.0.2 D2: Component counterfactual disparity}

\textbf{Definition 2 (Component counterfactual disparity, $\Delta$\_c).} Fix a protected attribute A with a counterfactual perturbation x $\rightarrow$ x$^{\prime}$ (a demographic swap of names, dialect, stated identity, or profile fields, in the lineage of counterfactual fairness \cite{kusner2017counterfactual}). The \emph{trajectory disparity} is

$\Delta$\_$\tau$ = d( O($\tau$(x)), O($\tau$(x$^{\prime}$)) ),

for a task-appropriate distance d (a flip indicator for discrete decisions, an absolute score difference for continuous ones; exactly the CFR/MASD pair of \cite{mayilvaghanan2026counterfactual}, generalized in \S{}4.2). The \emph{component disparity} $\Delta$\_c for component c $\in$ \{$\pi$, T, M, R, $\pi$\_ctrl\} is defined by \textbf{trace-replay}: log the full trajectory under x; replay it holding every other component's realized inputs and outputs fixed; perturb only the input that c consumes; and measure the divergence of c's output distribution,

$\Delta$\_c = E[ d( c(z\_c(x)), c(z\_c(x$^{\prime}$)) ) ] under replay of the logged trace.

$\Delta$\_c is the quantity that the clinical decomposition study of \cite{xu2026ducx} estimates when it attributes separable shares of a tool-using agent's total unfairness to the model, the data, and the tool-invocation stage. To our knowledge this is the only corpus paper that operationalizes the full decomposition, which is why \S{}4.3 elevates trace-replay to the canonical estimator and \S{}6 builds its protocol around it.

\paragraph*{3.0.3 D3: Conduction operators}

\textbf{Definition 3 (Conduction operator).} Every pipeline edge e from component c to component c$^{\prime}$ carries a \emph{conduction factor} $\varphi$\_e = $\Delta$\_out / $\Delta$\_in: the ratio of the disparity measurable at c$^{\prime}$'s output to the disparity delivered to c$^{\prime}$'s input. Edges fall into four modes:

\begin{center}\footnotesize
\begin{longtable}{|p{0.18\textwidth}|p{0.18\textwidth}|p{0.18\textwidth}|p{0.18\textwidth}|p{0.18\textwidth}|}
\hline \textbf{Operator} & \textbf{Factor} & \textbf{Meaning} & \textbf{Canonical corpus evidence} & \textbf{Typical edge} \\ \hline\endhead
\textbf{ATTENUATE} & $\varphi$ < 1 & The edge dampens incoming disparity & identity anonymization in debate \cite{choi2025identity}; embedder control in RAG \cite{kim2025mitigating}; multi-LLM debiasing \cite{owens2024multi} & guardrail, debate, fair-ranking edges \\ \hline
\textbf{PRESERVE} & $\varphi$ $\approx$ 1 & The edge conducts disparity unchanged & biased chatbots become biased agents under role assignment \cite{cao2026from}; reward models inherit dialect penalties \cite{mire2025rejected} & role hand-offs, alignment signals \\ \hline
\textbf{AMPLIFY} & $\varphi$ > 1 & The edge magnifies incoming disparity & inter-agent propagation \cite{nguyen2025social}; multi-turn degradation \cite{fan2024fairmt}; self-consuming training loops \cite{wyllie2024fairness} & interaction, accumulation, feedback edges \\ \hline
\textbf{GENERATE} & $\Delta$\_out > 0 on $\Delta$\_in = 0 & The edge creates disparity from demographically clean input & RAG undermines fairness even for vigilant users \cite{hu2024free}; persona assignment injects reasoning bias \cite{gupta2023bias}; reasoning itself introduces bias \cite{wu2025reasoning} & retrieval, persona, planning edges \\ \hline
\end{longtable}\end{center}

\emph{Takeaway:} AMPLIFY and GENERATE are the two genuinely agentic modes. Single-turn evaluation cannot exhibit either, and (as Table 1 and \S{}7 show) these are also the ones the field has instrumented least.

GENERATE is formally a new entry rather than a transformation ($\varphi$ is undefined on zero input); we treat a GENERATE edge as contributing its own $\Delta$\_c term. ATTENUATE is also the operator class of every mitigation in \S{}5: a mitigation \emph{is} an engineered $\varphi$ < 1 edge, which is what makes the mitigation map of Table 3 commensurable with the taxonomy.

\paragraph*{3.0.4 The conduction equation and the feedback recurrence}

Treating disparities as small and edges as locally multiplicative (a first-order approximation), the trajectory disparity decomposes as the \textbf{conduction equation}:

\textbf{$\Delta$\_$\tau$ $\approx$ $\Sigma$\_c $\Delta$\_c $\cdot$ $\Pi$\_\{e $\in$ path(c $\rightarrow$ O)\} $\varphi$\_e ,}

i.e., each component's entry disparity, conducted through the product of edge factors on its path to the consequential output, summed over entry loci. Across the loop-closing channel R, disparity obeys the \textbf{feedback recurrence}:

\textbf{$\Delta$\textasciicircum\{\}(t+1) = $\varphi$\_fb $\cdot$ $\Delta$\textasciicircum\{\}(t) + $\Delta$\_entry ,}

whose behavior is qualitatively bimodal. For $\varphi$\_fb < 1 the disparity converges to the steady state $\Delta$\_entry / (1 $-$ $\varphi$\_fb), persistent but bounded. For $\varphi$\_fb $\geq$ 1 it grows without bound: the regime documented for synthetic-data training loops \cite{wyllie2024fairness}, self-consuming performative loops \cite{wang2026observations}, and recommender feedback \cite{mansoury2020feedback}. Both expressions are an \emph{organizing approximation}, a bookkeeping identity under linearization and not a validated dynamical model. Edges in real agents are nonlinear, context-dependent, and interactive; the equation's role is to make explicit \emph{what an audit must estimate} (the $\Delta$\_c's and $\varphi$\_e's) and \emph{what can defeat naive auditing} (cancellation and amplification along paths), which is exactly how the \S{}6 protocol and the \S{}3.8 recipe use it.

\paragraph*{3.0.5 Propositions P1-P5}

Five propositions follow from D1 to D3 and the conduction equation. They are stated as claims the survey's evidence base supports or motivates, not as theorems about all possible agents; each carries its grounding.

\textbf{P1 (Locality).} \emph{Every measurable trajectory disparity decomposes into component entries: $\Delta$\_$\tau$ > 0 implies $\Delta$\_c > 0 for at least one locus c (counting GENERATE edges as entries).} Justification: this is the premise validated whenever decomposition succeeds. \cite{xu2026ducx} attributes distinct, separable shares of a clinical agent's disparity to the tool-invocation stage versus the model, and the component-specific audits surveyed in \S{}\S{}3.1 to 3.5 each isolate a nonzero $\Delta$\_c with other components held fixed.

\textbf{P2 (Masking).} \emph{Endpoint metrics can read $\approx$ 0 while internal disparities are large: whenever a downstream ATTENUATE edge (or sign-opposed entries) cancels an upstream $\Delta$\_c, the conduction equation yields $\Delta$\_$\tau$ $\approx$ 0 with $\Sigma$\_c |$\Delta$\_c| $\gg$ 0.} This is the survey's measurement gap stated as a structural fact rather than an empirical complaint. Justification: \cite{li2025actions} shows verbally neutral models issuing demographically disparate \emph{decisions}, with endpoint text-parity masking action disparity; \cite{xu2025quantifying} shows surface-level semantic parity masking statistical differentiation; and \cite{turpin2023language} shows the reverse mask, reasoning traces that misreport the biased features actually driving the output.

\textbf{P3 (Super-additivity, the agentic multiplier).} \emph{$\Delta$\_$\tau$ exceeds the naive sum $\Sigma$\_c $\Delta$\_c if and only if at least one AMPLIFY or GENERATE edge lies on a conduction path; with only PRESERVE/ATTENUATE edges, composition is sub-additive and component audits compose into a conservative system bound.} Justification: within the multi-agent locus this is directly evidenced. Stereotypical bias emerges, propagates, and is amplified across interacting agents, making the collective more biased than its members \cite{nguyen2025social}, and group-level fairness is emergent rather than inherited \cite{madigan2025emergent}. \S{}3.7 develops P3 and marks its cross-component generalization as conjecture.

\textbf{P4 (Mitigation-matching).} \emph{A mitigation reduces $\Delta$\_$\tau$ only if it installs an ATTENUATE edge on a path that carries a nonzero $\Delta$\_c; mitigation must match the entry locus and sit downstream of it, and an intervention at locus c neither bounds disparity generated downstream of c nor survives upstream amplification.} Justification: fairness-aware prompting fails against retrieval-injected bias because the GENERATE edge sits downstream of the prompt \cite{hu2024free}, while intervening \emph{at} the retrieval representation succeeds \cite{kim2025mitigating}; symmetrically, debate-level anonymization attenuates delegation bias precisely because it edits the identity cues at the delegation edge \cite{choi2025identity}. P4 is the organizing key of the mitigation map in \S{}5/Table 3.

\textbf{P5 (Measurement-adequacy).} \emph{An evaluation is adequate for an agent only if it can identify every $\Delta$\_c: it must observe component-level inputs and outputs (trace-level instrumentation with counterfactual replay), not endpoints alone; by P2, endpoint-only metrics are systematically unsound, and by P3, they are not even conservative.} Justification: models that pass single-turn fairness tests fail multi-turn ones \cite{fan2024fairmt}; the measured bias depends on which demographic cue channel is perturbed \cite{tonneau2026different}; and the major agent benchmarks (WebArena \cite{zhou2024webarena}, AgentBench \cite{liu2024agentbench}, Mind2Web \cite{deng2023mind2web}, SWE-bench \cite{jimenez2024swe}) log exactly the traces P5 requires but carry no demographic counterfactual instrumentation, a confrontation developed in \S{}4.4.

\paragraph*{3.0.6 Component-assignment rules and reading conventions}

A recurring defect of pipeline taxonomies is that papers float between cells, most acutely between \emph{tool selection} and \emph{retrieval}, since a retriever is both a tool one selects and a content source one conditions on. BCF resolves this with five assignment rules, keyed to what a component \emph{consumes and emits} in D1:

\begin{itemize}
\item \textbf{R1 (Selection vs. content).} If the disparity lives in \emph{which} external resource is invoked (a discrete choice over a registry, including the choice \emph{among} retrievers or sources), it is \textbf{C1 (tool/API selection)}. If it lives in \emph{what} an invoked resource returns and how that content conditions generation or persists (ranking, coverage, accumulation), it is \textbf{C2 (memory \& retrieval)}. Thus source-selection agents \cite{singh2025biasb} sit in C1; the effect of retrieved content on generation \cite{wu2024rag,hu2024free} sits in C2; an IR-systems treatment spanning both \cite{dai2024bias,dai2025trustworthy} is tagged to C2 as primary.
\item \textbf{R2 (Role vs. plan).} Disparity caused by \emph{who is assigned} (persona, role, or identity cues attached to an agent) is \textbf{C3 (delegation \& roles)}. Disparity in \emph{how the task is cut up}, ordered, or resourced, independent of identity assignment to agents, is \textbf{C4 (planning \& decomposition)}. Team-flavored allocation studies \cite{parziale2026once} are therefore C4: the skew is in the allocation, not in an assigned persona.
\item \textbf{R3 (Subject vs. user).} Disparity keyed to protected attributes of the \emph{decision subject} (the applicant, patient, defendant) flows through C1 to C4; disparity keyed to attributes of the \emph{interacting user} (especially attributes the agent \emph{infers} rather than receives \cite{staab2023memorization}) is \textbf{C5 (user modeling \& personalization)}.
\item \textbf{R4 (Temporal indexing).} Any of the above measured \emph{as a function of time} (turn, session, or deployment cycle) is the same locus indexed by t, tagged with the temporal axis of \S{}3.6, not a separate locus. Long-horizon drift is the $\varphi$\_fb regime of the feedback recurrence, not a sixth architectural component (\S{}3.6 justifies this at length).
\item \textbf{R5 (Multi-tagging).} A paper may evidence several cells and is counted in each (this is the $\times$N convention of the corpus and the multi-tag accounting of \S{}7), but it carries exactly one \emph{primary} locus: the component whose input/output the paper directly measures.
\end{itemize}

\textbf{Transferred-evidence convention.} Consistent with the inclusion criteria of \S{}2.4, some mechanism evidence below originates in single-turn or chatbot settings rather than instrumented agents. We mark such citations with \textbf{[$\dagger$]}: the mechanism transfers to the agentic edge by D1's construction (the same model occupies $\pi$), but the agentic \emph{measurement} has not been performed. These markers are themselves a map of open problems.

\subsubsection*{3.1 Locus C1: Tool / API Selection Bias}

\textbf{Definition.} A tool-using agent must, at each step, choose \emph{which} external tool, API, retriever, or knowledge source to invoke. Tool-selection bias is systematic, demographically- or content-correlated skew in that discrete choice, independent of and often invisible to the quality of the model's final text. By rule R1, everything about \emph{which} resource is invoked belongs here; what the resource \emph{returns} belongs to C2.

\textbf{Mechanism of harm.} Selection is a discrete decision over a tool registry, and $\pi$ makes such decisions with the same surface-form sensitivities that distort multiple-choice answering. Because the selected tool determines what evidence the rest of the trajectory conditions on, a skewed selection is conducted by every downstream edge: in BCF terms, C1 disparities enter early and travel far. The harm has a pre-agentic lineage in search: search engines have long been documented to encode racial and gendered skew in what they surface \cite{u2018algorithms}[$\dagger$] and in visual representation \cite{makhortykh2021detecting}[$\dagger$]. An agent that chooses among such engines inherits the choice-over-biased-infrastructure problem and adds its own selection pathology on top.

\textbf{Evidence, by operator.}
\emph{GENERATE.} The anchor result is that selection skew arises on demographically clean input: across tool registries, LLMs prefer particular tools regardless of fitness for the query, a reproducible \emph{tool-selection bias} residing in the selection policy rather than the answer \cite{blankenstein2025biasbusters}. The base pathology is positional and surface-form sensitivity in discrete choice \cite{pezeshkpour2023large}[$\dagger$], which carries directly into registries presented as lists. The mode is weaponizable: ToolTweak edits only a tool's name and description to systematically inflate its selection rate \cite{sneh2025tooltweak}, and recommender-agents are similarly steered by injected preference biases \cite{tang2026your}, a fairness problem whenever tools map to providers, vendors, or sources with unequal demographic footprints.
\emph{PRESERVE.} In a consequential clinical deployment, the decomposition of \cite{xu2026ducx} attributes a distinct share of a chest-X-ray agent's total disparity to the tool-invocation stage and shows it conducted into the final read. This is the cleanest existing demonstration that C1 is an independent entry locus whose contribution survives to the output, and the corpus's best instantiation of D2's trace-replay logic.
\emph{ATTENUATE.} Bias-aware retrieval agents that monitor and rebalance \emph{source selection} reduce the skew \cite{singh2025bias,singh2025biasb}, an engineered $\varphi$ < 1 edge installed exactly at the selection interface, consistent with P4.

\textbf{Open edge.} Existing work measures \emph{that} selection skews; almost none measures whether the skew correlates with protected attributes of the user or subject under a counterfactual swap. A tool-selection-disparity metric (does swapping a profile's demographics change which API the agent calls, holding the query fixed?) exists in no benchmark; this is the $\Delta$\_tool instance of D2 that \S{}4.3 formalizes.

\textbf{Conduction signature.} \emph{Entry:} the $\pi$$\rightarrow$T selection interface. \emph{Observable:} tool-choice distribution across counterfactual input pairs. \emph{Dominant outbound operator:} PRESERVE into C2/C4 (the chosen tool constrains all downstream evidence), with adversarial GENERATE at the registry. \emph{Metric:} tool-invocation disparity (CFR over tool choice; \S{}4.3).

\subsubsection*{3.2 Locus C2: Memory \& Retrieval: Bias Injection and Accumulation}

\textbf{Definition.} Agents persist state (conversation history, scratchpad memory, retrieved documents, cross-session stores) and condition each step on that accumulated context. Memory-and-retrieval bias is the introduction, \emph{accumulation}, or amplification of demographic skew through this conditioning, distinct from any bias in a single isolated response. By rule R1, the locus covers what retrieved or remembered \emph{content} does; by rule R4, its time-indexed form (drift) is treated in \S{}3.6.

\textbf{Mechanism of harm.} Two sub-mechanisms operate, mapping cleanly onto BCF's two non-trivial operators. \emph{Retrieval injection} is a GENERATE edge: grounding shifts the generator's distribution, so a disparity can appear that the closed-book model never exhibited. \emph{Temporal accumulation} is an AMPLIFY edge: each turn's context conditions the next, so per-turn skews compound. Critically, this locus inherits an entire prior discipline: fairness in ranking is a mature field with exposure-based and probability-based criteria and a decade of algorithms \cite{zehlike2022fairness}, and the LLM-era IR community has reframed those results as challenges for retrieval pipelines feeding generators \cite{dai2024bias,dai2025trustworthy}. BCF makes the inheritance precise: a fair-ranking guarantee is a bound on $\Delta$\_c at the retrieval interface, but by P4 it bounds $\Delta$\_$\tau$ only if no downstream GENERATE edge re-introduces disparity, which is exactly what the RAG results below show can happen.

\textbf{Evidence, by operator.}
\emph{GENERATE.} RAG can \emph{undermine} fairness even when retrieved documents are individually innocuous and even for users who actively prompt for fairness \cite{hu2024free}; retrieval augmentation measurably changes social-bias levels relative to the closed-book model \cite{zhang2025evaluating}; and the systematic evaluation of \cite{wu2024rag} confirms the retrieval step itself (which corpus, which ranking) injects unfairness with the generator held fixed.
\emph{AMPLIFY.} The adversarial extreme is knowledge-store poisoning that amplifies bias and steers downstream outputs \cite{wang2025bias}. Benignly, accumulation amplifies: fairness degrades over multi-turn dialogue in ways single-turn tests miss \cite{fan2024fairmt}; accumulating context literally shifts the model's beliefs as history grows \cite{geng2025accumulating}; and conversational history "geometrically traps" the model, making early framings self-reinforcing \cite{simhi2026old}. Position effects compound the problem: long-context models privilege the head and tail of the window \cite{liu2023lost}[$\dagger$], so \emph{which} remembered facts about a person get attended to is itself position-dependent and therefore skewable. Multi-agent conversation exhibits the same conversational-bias accumulation \cite{coppolillo2025unmasking}.
\emph{ATTENUATE.} The bias is partially controllable at the representation layer (tuning the embedder reduces RAG bias \cite{kim2025mitigating}) and at the ranking layer, where applying fair-exposure ranking inside RAG changes, and can improve, both fairness and quality of the generated output \cite{kim2024fair}. This is the most direct existing bridge between the fair-ranking tradition \cite{zehlike2022fairness} and agentic pipelines.

\textbf{Open edge.} The field measures accumulation in \emph{dialogue} (FairMT-Bench) but not in \emph{agentic memory}, meaning long-lived, structured, cross-session stores with explicit write policies. Whether bias compounds faster under tool-augmented memory than under plain conversation is unmeasured, and no work yet audits the memory \emph{write policy} (what gets persisted about whom) as a fairness object in its own right.

\textbf{Conduction signature.} \emph{Entry:} the T/M$\rightarrow$$\pi$ conditioning interface. \emph{Observable:} divergence of retrieved sets and memory states across counterfactual pairs. \emph{Dominant outbound operator:} AMPLIFY (accumulation across turns), with GENERATE at the grounding edge. \emph{Metric:} retrieval-set disparity plus the per-turn slope of $\Delta$\textasciicircum\{\}(t) (\S{}4.3).

\subsubsection*{3.3 Locus C3: Multi-Agent Delegation \& Role Assignment}

\textbf{Definition.} In multi-agent systems a controller $\pi$\_ctrl assigns subtasks, personas, or roles to constituent agents and routes information among them \cite{wu2023autogen,guo2024large,qian2024chatdev}. Delegation bias is disparity introduced by \emph{who is assigned what} and by the social dynamics that emerge among role-bearing agents, over and above any single agent's bias. By rule R2, identity-cue effects sit here; identity-free task-cutting sits in C4.

\textbf{Mechanism of harm.} The seed is persona conditioning, a GENERATE edge at the assignment interface; the distinctive agentic harm is that interaction then \emph{amplifies} what assignment generated, producing the one corpus-evidenced case of P3's super-additivity. Role assignment is thus doubly dangerous: it injects disparity and it wires the injected disparities into a network whose collective behavior is not predictable from its parts.

\textbf{Evidence, by operator.}
\emph{GENERATE.} Merely assigning a persona induces implicit reasoning biases that degrade performance on tasks unrelated to the persona; role assignment is active bias injection, not neutral scaffolding \cite{gupta2023bias}[$\dagger$]. Persona assignment also multiplies toxic degeneration in chat settings \cite{deshpande2023toxicity}[$\dagger$]. Benchmarked agentically, MALIBU uncovers implicit bias specifically in multi-agent role assignments \cite{mirza2025malibu}.
\emph{PRESERVE.} The chatbot-to-agent edge conducts: biased chatbots become biased agents through their assigned roles, with role assignment also degrading reliability \cite{cao2026from}. This is direct evidence that $\varphi$ $\approx$ 1 across the role hand-off, which is what licenses the [$\dagger$]-transfers above as mechanism evidence.
\emph{AMPLIFY.} The central result is that stereotypical bias \emph{emerges, propagates, and is amplified} across a multi-agent system, making the collective more biased than its members \cite{nguyen2025social}; emergent fairness in multi-agent decision systems is formally not inherited from components \cite{madigan2025emergent}; persona-induced bias scales from single agents to interacting societies \cite{li2025from}; implicit bias surfaces specifically \emph{in the interaction} between agents \cite{borah2024implicit}; and social hierarchy modulates the dynamics, with rank shaping persuasion and anti-social behavior \cite{campedelli2024want}.
\emph{ATTENUATE.} The mechanism is confirmed by its inverse: \emph{anonymizing} agent identities in debate reduces bias \cite{choi2025identity}. The disparity is caused by identity cues at the delegation edge, so deleting the cues installs $\varphi$ < 1 (P4 in action). Structured debate can also be steered toward equitable cultural alignment \cite{ki2025multiple}, and multi-LLM interaction can be engineered as a debiasing mechanism rather than an amplifier \cite{owens2024multi}: the same edge, opposite operator, depending on protocol design.

\textbf{Open edge.} Evidence is rich on \emph{whether} multi-agent setups amplify bias but thin on the \emph{delegation policy} as a controllable knob. There is no metric for assignment disparity (does $\pi$\_ctrl systematically route sensitive subtasks to lower-capability or more-biased sub-agents as a function of the subject's demographics?), and the amplification factor $\varphi$ in \cite{nguyen2025social} has been demonstrated but never \emph{bounded} as a function of network topology or protocol.

\textbf{Conduction signature.} \emph{Entry:} the $\pi$\_ctrl$\rightarrow$\{A\_i\} assignment interface. \emph{Observable:} the role/subtask assignment matrix across counterfactual pairs, plus inter-agent message content. \emph{Dominant outbound operators:} GENERATE (persona injection) followed by AMPLIFY (interaction). This is the only locus where both agentic operators are corpus-evidenced together. \emph{Metric:} assignment disparity over the delegation matrix (\S{}4.3).

\subsubsection*{3.4 Locus C4: Planning / Decomposition Disparities}

\textbf{Definition.} Agents plan: they decompose a goal into sub-goals, order steps, and reason explicitly (chain-of-thought, ReAct traces \cite{yao2023react}) before acting. Planning bias is demographic disparity that enters during this reasoning/decomposition stage, in the \emph{plan}, not merely the final answer. By rule R2, allocation skew belongs here even when the allocation is over agent-flavored "team members."

\textbf{Mechanism of harm.} The counterintuitive finding is that reasoning does not cleanly remove bias and can \emph{introduce} it, a GENERATE edge inside $\pi$ itself. Worse, this locus carries a distinctive P2 masking problem: the visible reasoning trace is not a faithful record of the features driving the decision, so auditing the \emph{stated} plan can miss the disparity the plan encodes. A scope caution is warranted: whether LLM "plans" constitute planning in the classical sense is contested \cite{kambhampati2024large}, but the fairness question is indifferent to that debate. Whatever the decomposition step is computationally, it emits orderings and allocations that can be counterfactually disparate.

\textbf{Evidence, by operator.}
\emph{GENERATE.} Adding explicit reasoning changes, and can worsen, social-bias outcomes relative to direct answering \cite{wu2025reasoning}; hidden veracity-linked beliefs surface \emph{during} problem-solving reasoning, encoding demographic assumptions a final-answer probe would not detect \cite{zhou2025veracity}; and in LLM-driven software-team composition, the planner allocates roles and tasks with demographic skew, a planning-stage disparity with direct hiring adjacency \cite{parziale2026once}.
\emph{Masking (P2).} Chain-of-thought explanations systematically misreport the biasing features that actually drive predictions \cite{turpin2023language}[$\dagger$], so a trace that \emph{reads} fair can implement a disparate decision procedure. Combined with the decisive measurement reframing of \cite{li2025actions} (agent \emph{decisions} reveal implicit biases that text-level evaluation misses), this locus is doubly hidden: neither the answer nor the stated reasoning is a sound observable. Only the counterfactual structure of the plan (D2 applied to plan objects) is.
\emph{ATTENUATE.} The same stage is a lever: reasoning-guided fine-tuning reduces bias \cite{kabra2025reasoning}, and fairness reward models can steer decision-making over the reasoning trace \cite{hall2025guiding}, interventions that act \emph{at} the generating edge, as P4 prescribes.

\textbf{Open edge.} No benchmark scores the \emph{plan itself} for disparity: are the sub-goals, their ordering, or the resource allocation a counterfactual function of the subject's demographics, independent of the final decision? Plan-allocation disparity is observed paper-by-paper but not operationalized; and by \cite{turpin2023language} any future plan-audit must be counterfactual-behavioral rather than trace-reading, a design constraint the field has not yet absorbed.

\textbf{Conduction signature.} \emph{Entry:} within $\pi$, at the decomposition step. \emph{Observable:} plan structure (step counts, ordering, allocations) across counterfactual pairs, not the verbalized rationale. \emph{Dominant outbound operator:} GENERATE, with strong P2 masking on the trace. \emph{Metric:} plan-allocation disparity and CFR over plan objects (\S{}4.3).

\subsubsection*{3.5 Locus C5: User Modeling \& Personalization Harms}

\textbf{Definition.} Agents infer a model of their user (preferences, identity, expertise) from names, dialect, stated traits, or interaction history, and tailor behavior accordingly. Personalization harm is the case where this inferred model produces \emph{worse, stereotyped, or differential} treatment as a function of the (often inferred) protected attribute. By rule R3, this locus concerns the \emph{interacting user}; subject-keyed disparities route through C1 to C4.

\textbf{Mechanism of harm.} The core tension is personalization versus stereotyping. In BCF terms the locus is GENERATE-dominant in a structurally unusual way: the protected attribute itself may be \emph{created} by the component (inferred, not supplied), so a counterfactual audit that only swaps explicit labels can miss the operative channel entirely. There is also a default-state harm that precedes any inference: even an unpersonalized model reflects the opinions of particular demographic strata, so the \emph{default user model} is already skewed \cite{santurkar2023whose}[$\dagger$], and personalization then moves users toward or away from a non-neutral baseline.

\textbf{Evidence, by operator.}
\emph{GENERATE.} Signaling user identity biases chatbot recommendations toward stereotype rather than genuine preference, collapsing adaptation into stereotyping \cite{kantharuban2024stereotype}[$\dagger$]. The trigger can be a mere name, which cues a presumed cultural identity that reshapes responses \cite{pawar2025presumed}; and the privacy precondition is established by the demonstration that LLMs \emph{infer} protected attributes (location, income, sex) from innocuous text \cite{staab2023memorization}, so an agent can model a user demographically without ever being told. The harm spans domains: personalized education delivers systematically different content quality by inferred student identity \cite{weissburg2024llms}; consumer-fairness evaluations of LLM recommenders find unequal recommendation utility across groups \cite{deldjoo2024cfairllm}, with a normative benchmarking framework following \cite{deldjoo2024normative}; and intrinsic socioeconomic bias supplies latent material for personalization to activate \cite{arzaghi2024understanding}.
\emph{PRESERVE (into the alignment signal).} Dialect is a second channel and shows conduction into training: reasoning performance drops on non-standard dialects \cite{lin2024assessing}[$\dagger$], and reward models penalize African American Language \cite{mire2025rejected}. The disparity is preserved \emph{into R}, the feedback channel, where \S{}3.6's recurrence takes over.
\emph{Method cautions.} Different demographic cue channels (name vs. explicit label vs. dialect) yield inconsistent conclusions about personalization bias, so single-cue audits can mislead \cite{tonneau2026different}, the empirical core of P5's adequacy demand; and debiasing personalized models trades off against utility \cite{vijjini2024exploring}, so ATTENUATE edges here are not free.

\textbf{Open edge.} Personalization harm has been measured in single responses, not over the \emph{persistent user model} an agent maintains across sessions, the object that bridges to \S{}3.6; and no agentic benchmark varies the cue \emph{channel} systematically despite \cite{tonneau2026different} showing conclusions are cue-fragile. The inference channel \cite{staab2023memorization} makes C5 the locus where D2's counterfactual swap is hardest to even define, an unresolved measurement-theory problem.

\textbf{Conduction signature.} \emph{Entry:} the inferred user model inside $\pi$/M. \emph{Observable:} behavior divergence under cue-channel-controlled counterfactual swaps (name, dialect, label, history). \emph{Dominant outbound operator:} GENERATE (the attribute is manufactured by inference), preserved into M and R. \emph{Metric:} cue-conditioned MASD across channels (\S{}4.2 to 4.3).

\subsubsection*{3.6 The Temporal Axis: Long-Horizon Fairness Drift Over Trajectories}

\textbf{Definition.} Long-horizon fairness drift is the \emph{temporal} degradation or self-amplification of fairness across turns, sessions, or deployment cycles, even when the per-step model is fixed.

\textbf{Why a temporal axis, not a sixth sibling.} Earlier drafts of this taxonomy, and adjacent surveys, treat long-horizon drift as a sixth component alongside the five architectural loci. BCF makes precise why that is a category error, and rule R4 codifies the repair. A component, in D1, is a \emph{place} in the architecture with its own inputs and outputs; drift has no such place. There is no "drift module" to instrument: drift is what the feedback recurrence $\Delta$\textasciicircum\{\}(t+1) = $\varphi$\_fb$\cdot$$\Delta$\textasciicircum\{\}(t) + $\Delta$\_entry \emph{does to any locus} when the loop closes: memory bias drifting (C2 $\times$ t), a user model ossifying (C5 $\times$ t), delegation patterns entrenching (C3 $\times$ t), or the policy itself shifting under retraining on its own outputs ($\pi$ $\times$ t through R). Treating drift as a sibling would double-count every time-indexed entry (violating P1's decomposition) and would imply the existence of a mitigation that targets "drift" as such, when by P4 every effective intervention must target a \emph{locus-and-edge}: either reducing some $\Delta$\_entry or forcing $\varphi$\_fb < 1. The taxonomy is therefore five architectural loci $\times$ one temporal axis: component $\times$ time made explicit. We retain the subsection at this rank because the temporal axis is where the field's least measured and most consequential phenomena live, not because it is a sixth place where bias enters.

\textbf{Mechanism of harm.} Three nested loops realize $\varphi$\_fb at increasing timescales, and all three are AMPLIFY-mode in the documented cases.

\textbf{Evidence, by operator.}
\emph{AMPLIFY, intra-trajectory (turns).} Fairness degrades over multi-turn dialogue \cite{fan2024fairmt}, and implicit bias accumulates and propagates specifically in LLM \emph{long-term memory}, making the agent more biased the longer it persists state \cite{ma2026implicit} (the C2 $\times$ t cell). Sycophancy is a named driver with a known origin: human-preference training itself induces models to conform to user views \cite{sharma2024understanding}[$\dagger$], so the conforming pressure is baked into $\pi$ by R before any dialogue begins. Multi-turn measurements then quantify the drift it produces: progressive conformity ("truth decay") across turns \cite{liu2025truth} and sycophancy growth over extended dialogues \cite{hong2025measuring}. An agent that incrementally adopts a user's biased framing is conducting that user's prejudice into its consequential actions.
\emph{AMPLIFY, deployment loops (retraining cycles).} Training on a model's own synthetic outputs amplifies bias across generations \cite{wyllie2024fairness}; the self-consuming performative loop has been characterized for LLMs specifically, with proposed remedies \cite{wang2026observations}; and the recommender-systems lineage supplies the canonical feedback-loop amplification result the LLM work inherits \cite{mansoury2020feedback}. Together these are the clearest empirical instances of the $\varphi$\_fb $\geq$ 1 regime.
\emph{Formal foundations (fair RL and sequential decision-making).} The sequential-fairness literature names the formal target: long-term fairness in sequential decision-making is definable and achievable under explicit dynamics \cite{hu2022achieving}; static fairness constraints provably can \emph{fail} to deliver fairness over a horizon, motivating criteria like the equal long-term benefit rate \cite{xu2024adapting}; and simulation-based long-term fairness analysis of ML-enabled systems supplies the systems-level tooling an agentic drift audit would need \cite{she2025fairsense}. In BCF terms, this literature is the theory of choosing interventions that force $\varphi$\_fb < 1. These are the only ATTENUATE results that exist for this axis, and none have yet been applied to an LLM agent.

\textbf{Open edge.} The sequential-fairness theory \cite{hu2022achieving,xu2024adapting} and the LLM-drift empirics \cite{fan2024fairmt,ma2026implicit} have not been joined: no work evaluates an \emph{LLM agent} against a formal long-term-fairness criterion over a realistic multi-session deployment. Equivalently, nobody has estimated $\varphi$\_fb for a deployed agent. The field's most consequential fairness regime is essentially unmeasured.

\textbf{Conduction signature.} \emph{Entry:} the R feedback edge (and time-indexed M). \emph{Observable:} the trajectory of $\Delta$\textasciicircum\{\}(t), its slope and curvature across turns, sessions, retraining cycles. \emph{Dominant operator:} AMPLIFY (the recurrence). \emph{Metric:} drift rate (slope of $\Delta$\textasciicircum\{\}(t)) and long-term benefit-rate gap \cite{xu2024adapting} (\S{}4.3).


\textbf{Table 1: Entry loci of the Bias Conduction Framework: harm mechanism, dominant operators, representative evidence, and measurement maturity.} Maturity: $\bullet$$\bullet$$\bullet$ = replicated findings \emph{and} standard benchmark; $\bullet$$\bullet$$\circ$ = replicated findings, no standard benchmark; $\bullet$$\circ$$\circ$ = isolated studies; $\circ$$\circ$$\circ$ = unmeasured. [$\dagger$] = transferred single-turn/chatbot evidence (\S{}3.0.6).

\begin{center}\footnotesize
\begin{longtable}{|p{0.18\textwidth}|p{0.18\textwidth}|p{0.18\textwidth}|p{0.18\textwidth}|p{0.18\textwidth}|}
\hline \textbf{Locus} & \textbf{Harm mechanism} & \textbf{Dominant operators} & \textbf{Representative evidence} & \textbf{Maturity} \\ \hline\endhead
\textbf{C1} Tool/API selection & Discriminatory routing over a registry; selection skew conducted into all downstream evidence & GENERATE, then PRESERVE & \cite{blankenstein2025biasbusters,sneh2025tooltweak,xu2026ducx,tang2026your} & $\bullet$$\circ$$\circ$ \\ \hline
\textbf{C2} Memory \& retrieval & Retrieval injection (grounding creates disparity); accumulation across turns & GENERATE + AMPLIFY & \cite{wu2024rag,hu2024free,wang2025bias,fan2024fairmt,kim2024fair} & $\bullet$$\bullet$$\circ$ \\ \hline
\textbf{C3} Delegation \& roles & Persona/role injection; emergent inter-agent amplification & GENERATE + AMPLIFY & \cite{gupta2023bias,cao2026from,nguyen2025social,madigan2025emergent,choi2025identity} & $\bullet$$\bullet$$\circ$ \\ \hline
\textbf{C4} Planning \& decomposition & Reasoning introduces disparity into plan structure; unfaithful traces mask it & GENERATE (P2-masked) & \cite{wu2025reasoning,zhou2025veracity,li2025actions,parziale2026once,turpin2023language}[$\dagger$] & $\bullet$$\circ$$\circ$ \\ \hline
\textbf{C5} User modeling \& personalization & Inferred-attribute stereotyping; skewed default user model; dialect penalties conducted into alignment & GENERATE, PRESERVE into R & \cite{kantharuban2024stereotype,staab2023memorization,pawar2025presumed,mire2025rejected,tonneau2026different} & $\bullet$$\bullet$$\circ$ \\ \hline
\textbf{T-axis} Long-horizon drift (any locus $\times$ time) & Feedback recurrence $\varphi$\_fb: sycophantic conformity, memory ossification, self-consuming retraining & AMPLIFY & \cite{ma2026implicit,liu2025truth,wyllie2024fairness,xu2024adapting,hu2022achieving} & $\bullet$$\circ$$\circ$ (theory mature; agent empirics nascent) \\ \hline
\end{longtable}\end{center}

\emph{Takeaway:} half the pipeline (tool selection, planning, and the entire temporal axis) sits at $\bullet$$\circ$$\circ$ despite documented real-world harm in hiring, lending, and clinical deployments; no locus reaches $\bullet$$\bullet$$\bullet$, because no standard benchmark instruments any component counterfactually.


\subsubsection*{3.7 Compounding Across the Pipeline: Proposition P3 and Its Limits}

The six entries are not additive, and Proposition P3 says exactly when they are not: \textbf{trajectory disparity is super-additive ($\Delta$\_$\tau$ > $\Sigma$\_c $\Delta$\_c) if and only if an AMPLIFY or GENERATE edge lies on a conduction path.} This is the formal content of what we call the \emph{agentic multiplier}, and it converts a rhetorical worry ("agents make bias worse") into a checkable property of a logged trace: estimate the $\varphi$\_e's; super-additivity is present precisely where some $\varphi$\_e > 1 or a fresh entry appears mid-path. Conversely (and this direction is just as consequential for auditors), if every edge in a pipeline is PRESERVE or ATTENUATE, component-level audits compose into a \emph{conservative} system-level bound, and certifying each component suffices. The entire practical question of whether per-component certification can ever underwrite system-level fairness claims therefore reduces to locating and bounding the amplifying edges.

\textbf{What is evidenced.} Within a single locus (C3) P3 is directly demonstrated. Stereotypical bias \emph{emerges} where none was salient, \emph{propagates} between agents, and is \emph{amplified} by their interaction, yielding a collective strictly more biased than its members \cite{nguyen2025social}; fairness in multi-agent decision systems is formally \emph{emergent}, not inherited, so component audits provably cannot be composed into a group-level guarantee in that setting \cite{madigan2025emergent}. Persona-induced bias scaling from single agents to societies \cite{li2025from} and interaction-borne implicit bias \cite{borah2024implicit} corroborate the same operator at the same locus. The temporal axis supplies a second evidenced amplifier, on the feedback edge rather than between components: self-consuming training loops with $\varphi$\_fb > 1 \cite{wyllie2024fairness,wang2026observations,mansoury2020feedback}.

\textbf{What is conjecture.} The cross-\emph{component} cascade is plausible under BCF, since every individual edge in that chain is separately evidenced as PRESERVE-or-worse. The intuitive story runs as follows: an attribute \emph{inferred} by user modeling (C5, \cite{staab2023memorization}) is \emph{persisted} by memory (C2, \cite{ma2026implicit}), conditions \emph{tool selection} (C1) and \emph{plan decomposition} (C4, \cite{li2025actions}), is \emph{delegated} along role lines (C3), and \emph{drifts} over the deployment horizon (T-axis). But we state plainly: \textbf{no paper in the corpus traces a single disparity across two or more *different* components of one agent.} We therefore record it as:

\begin{quote}\textbf{Conjecture 1 (Cross-component conduction).} In deployed tool-using agents, disparities entering at one locus are conducted with $\varphi$ $\geq$ 1 across component boundaries, so that multi-locus pipelines exhibit super-additive $\Delta$\_$\tau$ even when each component passes its own audit.\end{quote}

The conjecture is falsifiable with existing instruments: D2's trace-replay applied at two consecutive interfaces of the same logged trajectory estimates the cross-boundary $\varphi$ directly, and the audit protocol of \S{}6 is designed to produce exactly such paired estimates (results pending). Until such measurements exist, claims that "agentic bias compounds across the pipeline" (including informal versions in prior drafts of this survey) should be read as hypotheses licensed by P3's mechanism, not as established findings.

The compounding question is also the structural reason this survey is organized by component rather than by metric, and the conceptual core of the measurement gap of \S{}4: by P2, a token-level probe at any single hand-off can read clean while the consequential action is sharply disparate \cite{li2025actions,xu2025quantifying}; by P5, only trace-level counterfactual instrumentation can adjudicate Conjecture 1; and the coverage analysis of \S{}7 shows the cells P3 most implicates (action-level, counterfactual, long-horizon) are precisely where the field is emptiest.

\subsubsection*{3.8 Worked Example: A Hiring-Screener Agent Through the BCF, and a Four-Step Auditor Recipe}

To make the framework operational rather than ornamental, we walk one consequential system through it: a resume-screening agent, the deployment context with the deepest audit record \cite{raghavan2020mitigating,veldanda2023emily,an2025measuring,nghiem2024you,wilson2024gender} and the clearest regulatory exposure. We build the agent in four capability rungs, C0 to C3, and track what each rung adds to the conduction equation. (The rungs are capability levels of the worked example; the loci remain C1 to C5.)

\textbf{C0, bare model.} A single LLM call scores each resume. The loop degenerates to $\pi$ alone; the conduction equation collapses to $\Delta$\_$\tau$ = $\Delta$\_$\pi$. This rung is \emph{not} hypothetical parity: intersectional gender-and-race penalties in automated resume evaluation \cite{an2025measuring}, name-based callback disparities echoing the classic audit paradigm \cite{veldanda2023emily}, and occupation-skewed recommendations by name \cite{nghiem2024you} establish $\Delta$\_$\pi$ > 0 as the baseline entry. Everything QA-era fairness evaluation can see, it sees here.

\textbf{C1 rung, add retrieval.} The screener now retrieves similar past resumes and hiring precedents to ground its scores. Two new terms enter: $\Delta$\_ret at the C2 locus, conducted by the generation edge. This rung is also empirically grounded: race and gender bias arises \emph{in the retrieval step itself} when language models are used for resume screening \cite{wilson2024gender}, and by \cite{hu2024free}, instructing the generator to "score fairly" does not neutralize what retrieval injects (P4: the prompt sits upstream of the GENERATE edge). An endpoint audit that found C0 $\approx$ C1 parity could mask offsetting entries (P2); only comparing $\Delta$ at the retrieval interface against $\Delta$ at the output separates the terms.

\textbf{C2 rung, add tools and persistent memory.} The screener now chooses among background-check and skills-assessment APIs (C1 locus) and persists per-applicant notes across screening rounds (C2 $\times$ t). New terms: $\Delta$\_tool (does the demographic framing of a profile change \emph{which} verification API is invoked \cite{blankenstein2025biasbusters}, with vendor-controlled descriptions steerable adversarially \cite{sneh2025tooltweak}?) and a memory term that the recurrence now time-indexes. Early inferred attributes (possibly never stated, by \cite{staab2023memorization}) persist and accumulate \cite{ma2026implicit}, so an applicant re-screened in round two inherits round one's framing \cite{simhi2026old}.

\textbf{C3 rung, multi-agent committee.} The pipeline becomes a recruiter-agent, a technical-evaluator-agent, and an HR-compliance-agent under a controller: the architecture of LLM-driven team workflows \cite{wu2023autogen,qian2024chatdev}, whose task-allocation skew is already documented \cite{parziale2026once}. New terms: $\Delta$\_assign at C3, plus the only corpus-evidenced AMPLIFY-between-agents edge \cite{nguyen2025social,madigan2025emergent}. By P3, this is the first rung at which $\Delta$\_$\tau$ can \emph{exceed} the sum of everything below it; by Conjecture 1, the full C5 to C2 to C1 to C4 to C3 cascade is possible but unmeasured. Note what the ladder implies for governance: each rung is a routine engineering upgrade, invisible to a model-level fairness certificate issued at C0. The per-rung re-audit obligation this implies is taken up in \S{}8.5.

\textbf{A four-step auditor recipe.} The BCF reading of the ladder compresses into a procedure any team deploying an agent can execute, and which the \S{}6 protocol instantiates on a hiring-style task (results pending):

1. \textbf{Instrument the trace.} Log every component interface of D1 (tool calls and registry state, retrieved sets, memory reads/writes, inter-agent messages, plan objects, final action), not just the output. (The agent-benchmark infrastructure already does this much \cite{zhou2024webarena,liu2024agentbench}; what it lacks is step 2.)
2. \textbf{Counterfactually replay per component (D2).} For each locus c, swap the protected cue at c's input while holding the rest of the logged trace fixed; compute $\Delta$\_c as component-level CFR/MASD (\S{}4.2 to 4.3). Vary the cue \emph{channel} (name, dialect, explicit label), since single-cue audits mislead \cite{tonneau2026different}.
3. \textbf{Estimate the edge operators.} Compare disparities at consecutive interfaces to classify each edge as ATTENUATE / PRESERVE / AMPLIFY / GENERATE; on the feedback edge, regress $\Delta$\textasciicircum\{\}(t) across turns or sessions to estimate $\varphi$\_fb \cite{fan2024fairmt,she2025fairsense}. Any $\varphi$ > 1 or mid-path entry flags a P3 super-additivity risk.
4. \textbf{Match mitigation to the (locus, operator) cell and re-audit (P4).} Intervene at the entry locus or on the amplifying edge (embedder or ranking control for C2 \cite{kim2025mitigating,kim2024fair}, identity anonymization for C3 \cite{choi2025identity}, reasoning-stage steering for C4 \cite{kabra2025reasoning,hall2025guiding}, source-selection rebalancing for C1 \cite{singh2025biasb}), then repeat steps 2 to 3 to verify that the targeted $\Delta$\_c or $\varphi$\_e actually fell, and re-run over the horizon to confirm $\varphi$\_fb < 1.

Do not audit the verbalized rationale in place of step 2: stated reasoning is unfaithful to the operative features \cite{turpin2023language}, so only counterfactual behavior at component interfaces is a sound observable. Steps 1 to 2 are exactly the measurements Proposition P5 declares necessary, step 3 operationalizes P1 to P3, and step 4 operationalizes P4. The framework, run forward, \emph{is} the audit.



\subsection*{4. How to Evaluate Agent Fairness}

The translation from LLM-level fairness measurement to agent-level evaluation is neither direct nor complete. Single-turn benchmarks measure what a model \emph{says}; an agent must be evaluated on what the system \emph{does}: which tools it invokes, whose requests it escalates, how its plans allocate effort, and how disparities conduct across the loop. The Bias Conduction Framework of \S{}3 supplies the formal target: by Definition D2, fairness evaluation of an agent $\textbackslash\{\}mathcal\{A\} = \textbackslash\{\}langle \textbackslash\{\}pi, T, M, R \textbackslash\{\}rangle$ (Definition D1) is the estimation of component counterfactual disparities $\textbackslash\{\}Delta\_c$ via trace-replay, and by the conduction equation $\textbackslash\{\}Delta\_\textbackslash\{\}tau \textbackslash\{\}approx \textbackslash\{\}sum\_c \textbackslash\{\}Delta\_c \textbackslash\{\}cdot \textbackslash\{\}prod\_\{e\} \textbackslash\{\}varphi\_e$, an endpoint-only measurement estimates the left-hand side while leaving every term on the right unidentified. This section delivers the survey's second contribution in four steps. \S{}4.1 lifts the classical metric families (group, individual, counterfactual) onto the agent loop and asks what the impossibility results of \cite{kleinberg2017inherent} and \cite{chouldechova2017fair} become when the unit of evaluation is a trajectory rather than a prediction. \S{}4.2 formalizes trajectory-level CFR and MASD, generalizing \cite{mayilvaghanan2026counterfactual} into instances of D2's $\textbackslash\{\}Delta\_c$. \S{}4.3 supplies formal definitions for the three action-level disparity metrics the field currently lacks (tool-invocation, escalation, and plan-allocation disparity) as $\textbackslash\{\}Delta\_c$ specializations, discharging Proposition P5's requirement constructively. \S{}4.4 then inventories the benchmark field (Table\textasciitilde\{\}2) and shows, via Propositions P2 and P5, why no existing resource can certify an acting agent.

\subsubsection*{4.1 Metric Families Lifted to the Agent Loop}

Three classical fairness families dominate the measurement literature \cite{mehrabi2021survey,caton2024fairness,pessach2022review,barocas2023fairness}. Each imposes distinct demands and inherits distinct pathologies when lifted from single predictions to agent trajectories.

\paragraph*{4.1.1 Group fairness: from prediction distributions to action distributions}

Group fairness requires statistical equivalence of outcomes across protected groups. Its lineage runs from disparate-impact doctrine and the four-fifths rule, formalized for classifiers by \cite{feldman2015certifying}, through demographic parity to the error-rate criteria of equalized odds and equality of opportunity \cite{hardt2016equality}. The criteria are operationally attractive because they reduce fairness to a comparison of two conditional distributions.

Lifting group fairness to the agent loop changes \emph{what the outcome variable is}, not the comparison itself. In D1 terms, any measurable emission of the loop can serve as the outcome: the terminal decision $y$ (loan approval rates across applicant groups \cite{azime2025accept}, verdict severity across defendant groups \cite{hu2025llms}), but equally the intermediate emissions such as which tools in $T$ are invoked, what $M$ retrieves, which roles $R$ assigns, and whether the trajectory escalates to a human. Clinical audits already exploit this implicitly: \cite{omar2025sociodemographic} compares \emph{management plans} (an action distribution) rather than answer accuracy across sociodemographic variants, and \cite{adappanavar2025mfarm} scores harm-differentiated response patterns in decision support. What the QA-era surveys organize as a single "extrinsic disparity" number \cite{gallegos2024bias,chu2024fairness} is, under BCF, a \emph{family} of group-fairness statistics indexed by component: one $\textbackslash\{\}Delta\_c$-style comparison per locus. The practical failure of current group-fairness practice is therefore not the criterion but the denominator: existing evaluations aggregate over prompts and treat each response as atomic, so an agent that reaches equalized final answers through biased tool routing registers as fair. Proposition P1 (Locality) says the disparity has a component address; response-level group statistics discard the address.

\paragraph*{4.1.2 Individual fairness: from similar inputs to similar trajectories}

Individual fairness, in the metric formulation of \cite{dwork2012fairness}, demands that the treatment of two individuals differ by no more than their task-relevant dissimilarity, a Lipschitz condition $d\_\{\textbackslash\{\}mathrm\{out\}\}(h(x), h(x')) \textbackslash\{\}le L \textbackslash\{\}cdot d\_\{\textbackslash\{\}mathrm\{in\}\}(x, x')$. The notorious obstacle is specifying $d\_\{\textbackslash\{\}mathrm\{in\}\}$; the agentic lift adds a second: specifying $d\_\{\textbackslash\{\}mathrm\{out\}\}$ over \emph{trajectories}. Two near-identical applicants should receive the same verdict and be processed along comparably resourced paths (same depth of evidence retrieval, same tool quality, same escalation behavior). A natural candidate for $d\_\{\textbackslash\{\}mathrm\{out\}\}$ is an edit distance over action sequences, but no corpus paper yet defines or estimates one; the closest instantiation is the intervention-consistency audit of \cite{basu2026names}, which measures whether clinical and judicial recommendations change under name substitution while holding case features fixed, an individual-fairness probe at the terminal emission only. The trajectory version matters because of a failure mode group statistics cannot see: an agent can satisfy endpoint parity while routing counterfactual twins down systematically different paths (the \emph{path disparity} that Proposition P2 predicts can mask), so the like-cases-alike intuition is violated in the process even when it is satisfied in the verdict.

\paragraph*{4.1.3 Counterfactual fairness: from structural interventions to trace-replay}

Counterfactual fairness \cite{kusner2017counterfactual} requires that the outcome for an individual be invariant under an intervention on the protected attribute $A$: in structural-causal-model notation, $\textbackslash\{\}Pr(\textbackslash\{\}hat\{Y\}\_\{A \textbackslash\{\}leftarrow a\}(U) = y \textbackslash\{\}mid X = x, A = a) = \textbackslash\{\}Pr(\textbackslash\{\}hat\{Y\}\_\{A \textbackslash\{\}leftarrow a'\}(U) = y \textbackslash\{\}mid X = x, A = a)$, where the intervention propagates to the \emph{descendants of $A$} in the causal graph, not (as is sometimes loosely stated) to any variable that merely correlates with $A$. This family translates most naturally to agents, because the demographic swap can be applied to the full input and the full trajectory observed: this is exactly D2's trace-replay protocol, and \S{}4.2 develops it formally.

One honesty note is required, because it qualifies every counterfactual number reported in this literature. The swap operator $\textbackslash\{\}sigma$ used in practice (name substitution, demographic-cue editing) intervenes on the \emph{textual encoding} of $A$, not on the structural model: it alters the name but leaves other descendants of $A$ in the profile (dialect, address, school, employment gaps) untouched, which biases $\textbackslash\{\}Delta\_c$ estimates toward zero; conversely, an inconsistent swap that breaks profile coherence can inflate them. \cite{tonneau2026different} demonstrates the consequence empirically: different cue channels (names vs. explicit labels vs. dialect) yield inconsistent bias conclusions on the same models. Trace-replay counterfactuals are thus \emph{proxy interventions}, and rigorous protocols must vary the cue channel as a controlled factor, as the audit protocol of \S{}6 does.

\paragraph*{4.1.4 The upstream family: reward-model and alignment fairness}

A fourth family sits upstream of the loop. Agents are increasingly steered, scored, or filtered by reward models, and if the reward signal itself encodes demographic preferences, every downstream policy inherits them: \cite{song2025large} documents systematic inter-group reward gaps across demographic groups, and \cite{mire2025rejected} shows reward models penalize African American Language, a disparity \emph{generated} (operator GENERATE, on zero input disparity) at the alignment stage and then conducted through whatever the agent does. Evaluation must therefore propagate upstream from the action into the mechanism that rewards it; under BCF this is simply a $\textbackslash\{\}Delta\_c$ measurement at a locus outside the runtime loop, with conduction edges into $\textbackslash\{\}pi$.

\paragraph*{4.1.5 What the impossibility results become over trajectories}

The classical metric tensions do not dissolve at the agent level; they multiply. \cite{kleinberg2017inherent} prove that calibration and balance of error rates across groups cannot hold simultaneously except under equal base rates or perfect prediction; \cite{chouldechova2017fair} derives the same incompatibility between predictive parity and error-rate parity for recidivism instruments. These results bind any single decision point. Lifting them to trajectories surfaces three structurally new tensions, which we state as observations from the conduction equation rather than as theorems. None has yet been formally proved or empirically quantified for LLM agents, which is itself a finding about the field.

\textbf{Tension T1 (non-compositionality: per-step fairness does not compose).} An agent trajectory is a sequence of decision points, each formally subject to the classical impossibility. But satisfying a chosen criterion at every step does not deliver that criterion at the terminal decision, because each intermediate action reshapes the population reaching subsequent steps: the "base rates" at step $k$ are endogenous to the policy at steps $1, \textbackslash\{\}dots, k-1$. This is the sequential-screening analogue of the fair-sequential-decision results: \cite{hu2022achieving} show static per-step constraints fail to deliver long-term fairness, and \cite{xu2024adapting} show naive transfer of one-shot criteria can \emph{degrade} equal long-term benefit rates. The trajectory-level reading of \cite{kleinberg2017inherent} is therefore harsher than the original: the auditor must choose which criterion to satisfy and \emph{at which component} to satisfy it, knowing the choice perturbs the conditions under which every downstream criterion is evaluated.

\textbf{Tension T2 (outcome multiplicity under conduction coupling).} A single classifier has one outcome variable over which the impossibility binds. A trajectory has many: every component emission $z\_c$ plus the terminal $y$. The conduction equation couples them: $\textbackslash\{\}Delta\_\textbackslash\{\}tau \textbackslash\{\}approx \textbackslash\{\}sum\_c \textbackslash\{\}Delta\_c \textbackslash\{\}prod\_e \textbackslash\{\}varphi\_e$ means that driving one component's disparity to zero redistributes, rather than removes, the constraint when other edges amplify ($\textbackslash\{\}varphi\_e > 1$). Equalizing escalation rates while simultaneously equalizing terminal approval rates may be jointly infeasible when the escalation emission feeds an amplifying edge into the verdict. The label-choice problem of \cite{obermeyer2019dissecting} (where optimizing a proxy outcome (cost) produced large racial disparity on the target outcome (health need)) recurs here at every component boundary: each $z\_c$ is a proxy for the harm the auditor actually cares about, and the agent multiplies the number of proxy choices by the number of loci.

\textbf{Tension T3 (aggregation masking, P2 as a measurement tension).} Proposition P2 states that endpoint parity is consistent with large offsetting component disparities (attenuating edges, $\textbackslash\{\}varphi\_e < 1$, or sign-opposed $\textbackslash\{\}Delta\_c$ terms). As a metric tension, this means group fairness at the trajectory endpoint and individual fairness along the trajectory can disagree without bound: counterfactual twins reach identical verdicts through unequally resourced paths. The unfaithfulness results for chain-of-thought sharpen the warning. \cite{turpin2023language} show that verbalized reasoning does not reliably reflect the causes of model behavior, so $\textbackslash\{\}Delta\_c$ must be defined over \emph{executed actions} (tool calls issued, memory operations performed, delegations made), never over the agent's narrated rationale.

Finally, two construct-validity cautions cut across all families. \cite{blodgett2021stereotyping} catalogue pitfalls in fairness benchmark construction (incoherent contrasts, invalid stereotype operationalizations) that transfer wholesale to any agentic benchmark built by the same template, and \cite{selbst2019fairness} warn against the framing trap of evaluating an abstracted technical system apart from its deployment context: an agent loop scored in a sandbox omits exactly the institutional feedback (recourse, appeal, compounding access effects) through which the recurrence $\textbackslash\{\}Delta\textasciicircum\{\}\{(t+1)\} = \textbackslash\{\}varphi\_\{\textbackslash\{\}mathrm\{fb\}\} \textbackslash\{\}Delta\textasciicircum\{\}\{(t)\} + \textbackslash\{\}Delta\_\{\textbackslash\{\}mathrm\{entry\}\}$ operates in the world. Trajectory-level evaluation narrows the abstraction gap relative to QA scoring; it does not close it.

\subsubsection*{4.2 Trajectory-Level Counterfactual Metrics: CFR and MASD as Instances of $\textbackslash\{\}Delta\_c$}

We now make D2 operational. Fix an agent $\textbackslash\{\}mathcal\{A\} = \textbackslash\{\}langle \textbackslash\{\}pi, T, M, R \textbackslash\{\}rangle$, an input population $\textbackslash\{\}mathcal\{X\} = \textbackslash\{\}\{x\_1, \textbackslash\{\}dots, x\_N\textbackslash\{\}\}$ of task instances, and a demographic swap operator $\textbackslash\{\}sigma$ that rewrites the protected-attribute encoding of an input while holding task-relevant content fixed (\S{}4.1.3's proxy-intervention caveats apply, and $\textbackslash\{\}sigma$ should range over multiple cue channels \cite{tonneau2026different}). Executing $\textbackslash\{\}mathcal\{A\}$ on $x$ under the trace-replay protocol (pinned model versions, fixed decoding parameters, intercepted and demographically consistent tool/environment responses) yields a trajectory $\textbackslash\{\}tau\_\{\textbackslash\{\}mathcal\{A\}\}(x) = (s\_0, a\_1, o\_1, \textbackslash\{\}dots, a\_K, y)$ of actions $a\_k$, observations $o\_k$, and terminal decision $y$. For each component $c$ of the taxonomy, an extraction map $z\_c(\textbackslash\{\}cdot)$ reads off that component's emissions from the trace (the tool-invocation record for $T$, the retrieval set for $M$, the role assignments for $R$, the plan structure for $\textbackslash\{\}pi$'s decomposition, the terminal decision for the loop as a whole). Definition D2's component counterfactual disparity is then

$$
\textbackslash\{\}Delta\_c(\textbackslash\{\}mathcal\{A\}; \textbackslash\{\}mathcal\{X\}, \textbackslash\{\}sigma) \textbackslash\{\};=\textbackslash\{\}; \textbackslash\{\}mathbb\{E\}\_\{x \textbackslash\{\}sim \textbackslash\{\}mathcal\{X\}\} \textbackslash\{\}Big[ \textbackslash\{\}, d\_c\textbackslash\{\}big( z\_c(\textbackslash\{\}tau\_\{\textbackslash\{\}mathcal\{A\}\}(x)), \textbackslash\{\}; z\_c(\textbackslash\{\}tau\_\{\textbackslash\{\}mathcal\{A\}\}(\textbackslash\{\}sigma(x))) \textbackslash\{\}big) \textbackslash\{\}Big],
$$

where $d\_c$ is a distance on component $c$'s emission space. Every counterfactual metric in this literature is a choice of $(c, z\_c, d\_c)$. The two metrics introduced by \cite{mayilvaghanan2026counterfactual} for a contact-center quality-assurance agent are the canonical discrete and continuous choices, which we state in their trajectory-level generality.

\textbf{Definition 4.1 (Trajectory-level Counterfactual Flip Rate).} For a discrete emission $z\_c$ (a verdict, a tool choice, an escalation decision, a role assignment), the component CFR is $\textbackslash\{\}Delta\_c$ under the 0 to 1 metric:

$$
\textbackslash\{\}mathrm\{CFR\}\_c(\textbackslash\{\}mathcal\{A\}) \textbackslash\{\};=\textbackslash\{\}; \textbackslash\{\}frac\{1\}\{N\} \textbackslash\{\}sum\_\{i=1\}\textasciicircum\{\}\{N\} \textbackslash\{\}mathbb\{1\}\textbackslash\{\}big[\textbackslash\{\}, z\_c(\textbackslash\{\}tau\_\{\textbackslash\{\}mathcal\{A\}\}(x\_i)) \textbackslash\{\}neq z\_c(\textbackslash\{\}tau\_\{\textbackslash\{\}mathcal\{A\}\}(\textbackslash\{\}sigma(x\_i))) \textbackslash\{\},\textbackslash\{\}big].
$$

The terminal case $\textbackslash\{\}mathrm\{CFR\}\_y$ (with $z\_c = y$) recovers the original flip rate of \cite{mayilvaghanan2026counterfactual}; the intervention-consistency rate of \cite{basu2026names} is $\textbackslash\{\}mathrm\{CFR\}\_y$ with $\textbackslash\{\}sigma$ instantiated as name substitution.

\textbf{Definition 4.2 (Trajectory-level Mean Absolute Score Difference).} For a continuous emission $s\_c$ (a recommendation score, a rank, a confidence, a quality rating), the component MASD is $\textbackslash\{\}Delta\_c$ under the absolute-difference metric:

$$
\textbackslash\{\}mathrm\{MASD\}\_c(\textbackslash\{\}mathcal\{A\}) \textbackslash\{\};=\textbackslash\{\}; \textbackslash\{\}frac\{1\}\{N\} \textbackslash\{\}sum\_\{i=1\}\textasciicircum\{\}\{N\} \textbackslash\{\}big| \textbackslash\{\}, s\_c(\textbackslash\{\}tau\_\{\textbackslash\{\}mathcal\{A\}\}(x\_i)) - s\_c(\textbackslash\{\}tau\_\{\textbackslash\{\}mathcal\{A\}\}(\textbackslash\{\}sigma(x\_i))) \textbackslash\{\}, \textbackslash\{\}big|.
$$

CFR and MASD decompose the counterfactual question cleanly: \emph{does} the emission change under the swap, and \emph{by how much}? Reporting both prevents a benchmark from hiding many small shifts behind a low flip count or a few large flips behind a small mean.

\textbf{Estimating the conduction operators.} Because Definitions 4.1 to 4.2 are component-indexed, they also yield the first estimator for D3's operators. For a pipeline edge $e: c \textbackslash\{\}to c'$, the empirical conduction coefficient is the disparity ratio across the edge,

$$
\textbackslash\{\}hat\{\textbackslash\{\}varphi\}\_e \textbackslash\{\};=\textbackslash\{\}; \textbackslash\{\}frac\{\textbackslash\{\}Delta\_\{c'\}(\textbackslash\{\}mathcal\{A\}; \textbackslash\{\}mathcal\{X\}, \textbackslash\{\}sigma)\}\{\textbackslash\{\}Delta\_c(\textbackslash\{\}mathcal\{A\}; \textbackslash\{\}mathcal\{X\}, \textbackslash\{\}sigma)\} \textbackslash\{\}qquad (\textbackslash\{\}Delta\_c > 0),
$$

classifying the edge as attenuating ($\textbackslash\{\}hat\{\textbackslash\{\}varphi\}\_e < 1$), preserving ($\textbackslash\{\}hat\{\textbackslash\{\}varphi\}\_e \textbackslash\{\}approx 1$), or amplifying ($\textbackslash\{\}hat\{\textbackslash\{\}varphi\}\_e > 1$); the GENERATE case is diagnosed when $\textbackslash\{\}Delta\_c \textbackslash\{\}approx 0$ at every upstream locus yet $\textbackslash\{\}Delta\_\{c'\} > 0$. This turns the BCF's operator taxonomy from a descriptive vocabulary into a measurable quantity, and it is the estimator the audit protocol of \S{}6 is designed to compute (results pending).

\textbf{The existing counterfactual evaluation frameworks, read against this scaffold.} The 2025 to 2026 wave of counterfactual evaluation frameworks shares a common three-step scaffold (construct demographically paired inputs, observe outputs under each condition, compute a distance statistic), and the frameworks differ chiefly in which corner of the $(c, z\_c, d\_c, \textbackslash\{\}sigma)$ design space they develop. CAFFE \cite{parziale2025systematic} and its companion SCOPE dataset \cite{parziale2026scope} develop the $\textbackslash\{\}sigma$ axis: systematic, multi-attribute perturbation generation with coverage of attributes existing benchmarks leave sparse. FiSCo \cite{xu2025quantifying} develops the $d\_c$ axis: a semantic, statistically-tested group-difference measure that detects disparities surface-level token comparison misses, important for agents whose verbally equalized answers can sit atop differentiated reasoning. The Equal-Access audit \cite{amirimargavi2026equal} develops the $z\_c$ axis toward service quality: breadth and actionability of responses as the emission, mapping in agentic settings onto retrieval depth and tool-complement disparities. \cite{basu2026names} develops the deployment-realism axis, applying flip-rate logic in clinical and judicial scenarios where the emission is a consequential recommendation. What none of them yet provides, including the originating CFR/MASD deployment of \cite{mayilvaghanan2026counterfactual} (which scores a single QA-scoring emission of a deployed agent), is coverage of \emph{intermediate} trajectory state: the tool selected at step 3, the memory written at step 5, the subagent delegated to at step 2. In BCF terms, the field has built increasingly refined estimators of $\textbackslash\{\}Delta\_\textbackslash\{\}tau$ while leaving every interior $\textbackslash\{\}Delta\_c$ and every $\textbackslash\{\}varphi\_e$ unestimated. By Proposition P1 those interior quantities are where the harm has an address; by Proposition P4, they are what mitigation must be matched against.

\subsubsection*{4.3 Action-Level Disparity Metrics: Formal Definitions}

The sharpest edge of the measurement gap has been the absence of \emph{defined} metrics for action-level disparity: prior work demonstrates that actions carry bias that answers conceal \cite{li2025actions,blankenstein2025biasbusters,xu2026ducx} but stops short of canonical formulas. We supply them here as $\textbackslash\{\}Delta\_c$ specializations, in both counterfactual (paired) and group (population) forms. Throughout, $a \textbackslash\{\}in \textbackslash\{\}\{0, 1\textbackslash\{\}\}$ indexes the protected-group contrast for group forms, and the counterfactual forms are paired over $(x, \textbackslash\{\}sigma(x))$ as in \S{}4.2.

\textbf{Definition 4.3 (Tool-Invocation Disparity, TID).} Let $n\_t(\textbackslash\{\}tau)$ be the number of invocations of tool $t \textbackslash\{\}in T$ in trajectory $\textbackslash\{\}tau$, and $p\_\textbackslash\{\}tau(t) = n\_t(\textbackslash\{\}tau) / \textbackslash\{\}sum\_\{t' \textbackslash\{\}in T\} n\_\{t'\}(\textbackslash\{\}tau)$ the empirical invocation distribution. The counterfactual form is $\textbackslash\{\}Delta\_c$ at the tool-selection locus with $z\_c = p\_\textbackslash\{\}tau$ and $d\_c$ the total-variation distance:

$$
\textbackslash\{\}mathrm\{TID\}\_\{\textbackslash\{\}mathrm\{cf\}\}(\textbackslash\{\}mathcal\{A\}) \textbackslash\{\};=\textbackslash\{\}; \textbackslash\{\}mathbb\{E\}\_\{x \textbackslash\{\}sim \textbackslash\{\}mathcal\{X\}\} \textbackslash\{\}left[ \textbackslash\{\}tfrac\{1\}\{2\} \textbackslash\{\}sum\_\{t \textbackslash\{\}in T\} \textbackslash\{\}big| \textbackslash\{\}, p\_\{\textbackslash\{\}tau\_\{\textbackslash\{\}mathcal\{A\}\}(x)\}(t) - p\_\{\textbackslash\{\}tau\_\{\textbackslash\{\}mathcal\{A\}\}(\textbackslash\{\}sigma(x))\}(t) \textbackslash\{\}, \textbackslash\{\}big| \textbackslash\{\}right],
$$

and the group form compares group-conditional mean invocation distributions, $\textbackslash\{\}mathrm\{TID\}\_\{\textbackslash\{\}mathrm\{grp\}\} = \textbackslash\{\}tfrac\{1\}\{2\} \textbackslash\{\}sum\_\{t \textbackslash\{\}in T\} \textbackslash\{\}big| \textbackslash\{\}mathbb\{E\}[\textbackslash\{\}, p\_\textbackslash\{\}tau(t) \textbackslash\{\}mid a\{=\}1 \textbackslash\{\},] - \textbackslash\{\}mathbb\{E\}[\textbackslash\{\}, p\_\textbackslash\{\}tau(t) \textbackslash\{\}mid a\{=\}0 \textbackslash\{\},] \textbackslash\{\}big|$. TID answers the question \S{}3.1 left open: does swapping a subject's demographics change \emph{which} API the agent calls, holding the query fixed? The existence evidence is already in place: tool-selection skew is measurable and manipulable \cite{blankenstein2025biasbusters,sneh2025tooltweak}, rides on the same surface-form sensitivities as option-order bias \cite{pezeshkpour2023large}, and accounts for an attributable share of total disparity in the one clinical agent where the decomposition has been run \cite{xu2026ducx}. No benchmark computes TID under demographic counterfactuals.

\textbf{Definition 4.4 (Escalation Disparity, ESD).} Let $e(\textbackslash\{\}tau) \textbackslash\{\}in \textbackslash\{\}\{0, 1\textbackslash\{\}\}$ indicate that the trajectory escalates: defers to a human reviewer, requests additional documentation, refuses, or routes to a higher-scrutiny queue. The group form is a parity gap, with a companion ratio form aligned with the four-fifths convention of disparate-impact doctrine \cite{feldman2015certifying}:

$$
\textbackslash\{\}mathrm\{ESD\}\_\{\textbackslash\{\}mathrm\{grp\}\} \textbackslash\{\};=\textbackslash\{\}; \textbackslash\{\}big| \textbackslash\{\}Pr(e(\textbackslash\{\}tau)\{=\}1 \textbackslash\{\}mid a\{=\}1) - \textbackslash\{\}Pr(e(\textbackslash\{\}tau)\{=\}1 \textbackslash\{\}mid a\{=\}0) \textbackslash\{\}big|,
\textbackslash\{\}qquad
\textbackslash\{\}mathrm\{ESD\}\_\{\textbackslash\{\}mathrm\{ratio\}\} \textbackslash\{\};=\textbackslash\{\}; \textbackslash\{\}frac\{\textbackslash\{\}min\_a \textbackslash\{\}Pr(e(\textbackslash\{\}tau)\{=\}1 \textbackslash\{\}mid a)\}\{\textbackslash\{\}max\_a \textbackslash\{\}Pr(e(\textbackslash\{\}tau)\{=\}1 \textbackslash\{\}mid a)\},
$$

and the counterfactual form is the escalation-specific flip rate, $\textbackslash\{\}mathrm\{ESD\}\_\{\textbackslash\{\}mathrm\{cf\}\} = \textbackslash\{\}mathbb\{E\}\_x \textbackslash\{\}big[ \textbackslash\{\}, | e(\textbackslash\{\}tau\_\{\textbackslash\{\}mathcal\{A\}\}(x)) - e(\textbackslash\{\}tau\_\{\textbackslash\{\}mathcal\{A\}\}(\textbackslash\{\}sigma(x))) | \textbackslash\{\}, \textbackslash\{\}big]$, i.e., $\textbackslash\{\}mathrm\{CFR\}\_c$ with $z\_c = e$. Escalation is the margin where differential treatment hides behind identical verdicts: two counterfactual twins may both be approved, but one only after demands for extra documentation that impose cost and attrition. Clinical evidence that management and referral behavior shifts with sociodemographic cues \cite{omar2025sociodemographic}, and the safety-utility asymmetries of personalized refusal behavior \cite{vijjini2024exploring}, indicate this margin is live; no benchmark instruments it.

\textbf{Definition 4.5 (Plan-Allocation Disparity, PAD).} Let $\textbackslash\{\}mathbf\{r\}(\textbackslash\{\}tau) \textbackslash\{\}in \textbackslash\{\}mathbb\{R\}\textasciicircum\{\}m$ be the plan's resource vector, e.g., number of plan steps, tool-call budget consumed, retrieval depth, deliberation tokens, count of verification actions. The counterfactual form is $\textbackslash\{\}Delta\_c$ at the planning locus under a weighted $\textbackslash\{\}ell\_1$ metric:

$$
\textbackslash\{\}mathrm\{PAD\}\_\{\textbackslash\{\}mathrm\{cf\}\}(\textbackslash\{\}mathcal\{A\}) \textbackslash\{\};=\textbackslash\{\}; \textbackslash\{\}mathbb\{E\}\_\{x \textbackslash\{\}sim \textbackslash\{\}mathcal\{X\}\} \textbackslash\{\}Big[ \textbackslash\{\}, \textbackslash\{\}big\textbackslash\{\}| \textbackslash\{\}mathbf\{r\}(\textbackslash\{\}tau\_\{\textbackslash\{\}mathcal\{A\}\}(x)) - \textbackslash\{\}mathbf\{r\}(\textbackslash\{\}tau\_\{\textbackslash\{\}mathcal\{A\}\}(\textbackslash\{\}sigma(x))) \textbackslash\{\}big\textbackslash\{\}|\_\{1, w\} \textbackslash\{\}Big],
\textbackslash\{\}qquad \textbackslash\{\}|\textbackslash\{\}mathbf\{v\}\textbackslash\{\}|\_\{1, w\} = \textbackslash\{\}sum\_\{j=1\}\textasciicircum\{\}\{m\} w\_j |v\_j|,
$$

with weights $w\_j$ fixed in advance to reflect the deployment's resource stakes. When the plan allocates work \emph{to people} rather than effort to a case (team composition, task assignment), the relevant form is the maximum per-role assignment gap, $\textbackslash\{\}mathrm\{PAD\}\_\{\textbackslash\{\}mathrm\{alloc\}\} = \textbackslash\{\}max\_\{\textbackslash\{\}rho \textbackslash\{\}in \textbackslash\{\}mathcal\{R\}\} \textbackslash\{\}big| \textbackslash\{\}Pr(\textbackslash\{\}mathrm\{assign\}(\textbackslash\{\}rho) \textbackslash\{\}mid a\{=\}1) - \textbackslash\{\}Pr(\textbackslash\{\}mathrm\{assign\}(\textbackslash\{\}rho) \textbackslash\{\}mid a\{=\}0) \textbackslash\{\}big|$, which is the metric implicit in the demographically skewed task allocations documented by \cite{parziale2026once} and the agent-decision disparities of \cite{li2025actions}, neither of which states it as a reusable formula.

Three remarks complete the definitions. \emph{Estimation.} All three metrics are paired statistics under trace-replay: variance is computed over counterfactual pairs, confidence intervals attach to each $\textbackslash\{\}Delta\_c$, and the configuration-by-cue-channel grid requires multiple-comparison correction. The audit protocol of \S{}6 instantiates exactly these estimators for a hiring task across agent configurations C0 to C3 (results pending). Determinism matters: without pinned versions, fixed decoding, and intercepted environment responses, flip rates confound demographic sensitivity with sampling noise. \emph{Scope.} The definitions deliberately score executed actions, not narrated rationales, per the unfaithfulness caveat of \cite{turpin2023language}. \emph{Adequacy.} Proposition P5 states that an evaluation suite is adequate for an agent only if it bounds $\textbackslash\{\}Delta\_c$ for every component with nonzero conduction weight in $\textbackslash\{\}Delta\_\textbackslash\{\}tau \textbackslash\{\}approx \textbackslash\{\}sum\_c \textbackslash\{\}Delta\_c \textbackslash\{\}prod\_e \textbackslash\{\}varphi\_e$. Definitions 4.3 to 4.5, together with the component CFR/MASD of \S{}4.2, constitute a minimal instrument set meeting P5 for the $T$, $R$, and planning surfaces of D1. By Proposition P4 (mitigation-matching), they are also the precondition for crediting any mitigation with changing what an agent \emph{does} rather than what it says. What was previously described in this literature as having "no canonical formula" now has one; what remains missing is the benchmark infrastructure to run it against, which is \S{}4.4's subject.

\subsubsection*{4.4 Datasets and Benchmarks: An Inventory of the Gap}

Table\textasciitilde\{\}2 inventories the evaluation resources in the corpus alongside the agent-capability tools against which any claim of a "missing benchmark" must be tested. We classify each resource by \emph{fairness level} (QA (scored answers to fixed questions), generation (scored open-ended text), or action (scored executed behavior)) and by whether it is \emph{agentic}, which we define strictly: a resource is agentic if and only if it (i) executes multi-step, tool-using trajectories and (ii) scores a fairness quantity (some $\textbackslash\{\}Delta\_c$) over them under demographic instrumentation. Both conditions matter, and the table's structure shows the field has split them between two non-communicating literatures.

\textbf{Table 2. Fairness evaluation resources and agent capability harnesses.}

\begin{center}\footnotesize
\begin{longtable}{|p{0.15\textwidth}|p{0.15\textwidth}|p{0.15\textwidth}|p{0.15\textwidth}|p{0.15\textwidth}|p{0.15\textwidth}|}
\hline \textbf{Benchmark / resource} & \textbf{Fairness level} & \textbf{Domain} & \textbf{Primary metric} & \textbf{Agentic?} & \textbf{Code/Data released?} \\ \hline\endhead
BBQ \cite{parrish2022bbq} & QA & General social bias & Accuracy/bias score gap & No & Yes \\ \hline
BOLD \cite{dhamala2021bold} & Generation & Open-ended generation & Sentiment/regard disparity & No & Yes \\ \hline
StereoSet \cite{nadeem2021stereoset} & QA/Gen & General & Stereotype score vs. LM score & No & Yes \\ \hline
CrowS-Pairs \cite{nangia2020crows} & Generation & General & Paired pseudo-likelihood gap & No & Yes \\ \hline
HolisticBias \cite{smith2022sorry} & Generation & General (600+ descriptors) & Perplexity/regard disparity & No & Yes \\ \hline
DecodingTrust \cite{wang2023decodingtrust} & QA/Gen & Trustworthiness (fairness subset) & Stereotype agreement; DP/EO gaps & No & Yes \\ \hline
TrustLLM \cite{huang2024trustllm} & QA/Gen & Trustworthiness & Fairness subscores & No & Yes \\ \hline
TrustGPT \cite{huang2023trustgpt} & Generation & Toxicity/bias/value & Toxicity and bias gaps & No & n.r. \\ \hline
HELM \cite{liang2023holistic} & QA/Gen & Broad capability & Bias/fairness slice metrics & No & Yes \\ \hline
FairMedQA \cite{xiao2025fairmedqa} & QA & Clinical & Group response disparity & No & n.r. \\ \hline
MedEqualQA \cite{ghosh2025medequalqa} & QA & Clinical & Counterfactual consistency & No & n.r. \\ \hline
mFARM \cite{adappanavar2025mfarm} & QA & Clinical decision support & Multi-faceted harm scores & No & n.r. \\ \hline
FMBench \cite{wu2024fmbench} & QA & Medical multimodal & Group performance parity & No & n.r. \\ \hline
Health-equity toolbox \cite{pfohl2024toolbox} & QA/Gen & Clinical & Rubric-rated harm rates & No & Yes \\ \hline
Accept-or-Deny \cite{azime2025accept} & QA (decision) & Lending & Approval-parity gap & No & n.r. \\ \hline
LLMs on Trial \cite{hu2025llms} & QA (decision) & Judicial & Verdict disparity & No & n.r. \\ \hline
Name-intervention audit \cite{basu2026names} & QA (decision) & Clinical/judicial & Intervention consistency (flip rate) & No & n.r. \\ \hline
CEB \cite{wang2024ceb} & QA/Gen & General (compositional) & Composite stereotype/toxicity/performance & No & Yes \\ \hline
CAFFE \cite{parziale2025systematic} & QA/Gen & General & Multi-attribute flip metrics & No & n.r. \\ \hline
SCOPE \cite{parziale2026scope} & QA/Gen & General & Counterfactual flip rate & No & n.r. \\ \hline
LangFair \cite{bouchard2024bring} & QA/Gen & Use-case-specific toolkit & Configurable group + counterfactual & No & Yes \\ \hline
FiSCo \cite{xu2025quantifying} & Generation & General & Semantic group-difference statistic & No & n.r. \\ \hline
Equal-Access audit \cite{amirimargavi2026equal} & Generation & Service quality & Response-quality disparity & No & n.r. \\ \hline
Reward-model fairness \cite{song2025large} & Upstream (reward) & Alignment & Inter-group reward gap & No & n.r. \\ \hline
CFR/MASD deployment audit \cite{mayilvaghanan2026counterfactual} & QA (agent output) & Contact center & CFR, MASD & Partial$\textasciicircum\{\}\{\textbackslash\{\}ast\}$ & No (proprietary) \\ \hline
FairMT-Bench \cite{fan2024fairmt} & Generation (multi-turn) & Dialogue & Multi-turn bias rate & Partial$\textasciicircum\{\}\{\textbackslash\{\}ast\}$ & Yes \\ \hline
MALIBU \cite{mirza2025malibu} & Generation (multi-agent) & Multi-agent personas & Persona bias scores & Partial$\textasciicircum\{\}\{\textbackslash\{\}ast\}$ & n.r. \\ \hline
WebArena \cite{zhou2024webarena} & None & Web tasks & Task success rate & No$\textasciicircum\{\}\{\textbackslash\{\}dagger\}$ & Yes \\ \hline
Mind2Web \cite{deng2023mind2web} & None & Web tasks & Step/element accuracy & No$\textasciicircum\{\}\{\textbackslash\{\}dagger\}$ & Yes \\ \hline
AgentBench \cite{liu2024agentbench} & None & Eight interactive environments & Task success rate & No$\textasciicircum\{\}\{\textbackslash\{\}dagger\}$ & Yes \\ \hline
SWE-bench \cite{jimenez2024swe} & None & Software engineering & Issue resolution rate & No$\textasciicircum\{\}\{\textbackslash\{\}dagger\}$ & Yes \\ \hline
\end{longtable}\end{center}

$\textasciicircum\{\}\{\textbackslash\{\}ast\}$ \emph{Partial}: fairness-instrumented beyond a single turn (multi-turn dialogue, multi-agent output, or a deployed agent's scoring emission) but does not instrument tool calls, memory operations, plans, or delegations, i.e., estimates $\textbackslash\{\}Delta\_\textbackslash\{\}tau$ or a single $\textbackslash\{\}Delta\_c$, never the interior of the conduction equation. $\textasciicircum\{\}\{\textbackslash\{\}dagger\}$ Executes genuine multi-step, tool-using trajectories but contains no demographic instrumentation, so it cannot score fairness at any level. n.r. = public code/data release not verified at the time of writing. \emph{Takeaway: no row pairs fairness instrumentation with agentic execution (the Agentic? column contains no "yes") which is the structured, falsifiable form of the measurement gap predicted by Propositions P2 and P5.}

\textbf{Three lineages, one boundary.} Read as a whole rather than row by row, the table contains three lineages that all halt at the same boundary. The \emph{QA-bias lineage} (BBQ through HolisticBias) fixed the field's unit of analysis (a prompt, a response, a disparity statistic) before agents existed, and the \emph{trustworthiness suites} (DecodingTrust, TrustLLM, TrustGPT, HELM) scaled that unit up in breadth without changing it: fairness is one slice of a prompt-response evaluation, never a property of executed behavior \cite{wang2023decodingtrust,huang2024trustllm,liang2023holistic}. The \emph{domain-decision tier} (lending, judicial, clinical) moved the content to consequential decisions but kept the form: a vignette in, a recommendation out, a parity statistic over recommendations \cite{azime2025accept,hu2025llms}. The \emph{counterfactual tools} (CAFFE, SCOPE, LangFair, FiSCo, the audits of \S{}4.2) refined the statistics (they are precisely the $(\textbackslash\{\}sigma, d\_c)$ machinery \S{}4.2 formalizes) but apply it at natural-language interfaces, not at trajectory interiors. The dataset survey of \cite{zhang2025datasets} confirms the pattern across a much larger inventory than ours, and the agent-evaluation surveys concede the complement: fairness appears as an aspirational dimension without trajectory-level metrics in \cite{mohammadi2025evaluation}, \cite{yehudai2025survey}, and \cite{chang2024survey}.

\textbf{Why clinical benchmarks stop at vignettes.} The clinical tier deserves a structural explanation rather than a complaint, because its authors are methodologically careful and the vignette boundary is not an oversight. A vignette is the unique unit at which three costly requirements are simultaneously cheap. First, \emph{ground truth}: a static case has an expert-adjudicable answer, whereas a trajectory's quality is path-dependent. Judging whether an agent's third tool call was clinically appropriate requires re-adjudicating the state the agent itself created, at expert cost per step rather than per case (the rubric-and-adversarial design of \cite{pfohl2024toolbox} already strains this budget at the answer level). Second, \emph{privacy and provenance}: real clinical trajectories are protected health information end to end; vignettes are shareable artifacts, which is why every released clinical resource in the table is vignette-formed. Third, \emph{counterfactual control}: in a static case, the demographic swap $\textbackslash\{\}sigma$ is a text edit; in an interactive environment, the swap must propagate consistently through every tool response, retrieved record, and intermediate observation, or the counterfactual pair silently decoheres. That is exactly the instrumented-environment engineering that no fairness benchmark has built. There is also an inherited-pattern effect: clinical LLM evaluation descends from medical-exam QA and the BBQ-style bias template, so the question, not the episode of care, became the unit of measurement. The cost of the boundary is stated by Proposition P2: a model certified near-parity on vignette QA can still triage, escalate, and order differently across demographic counterfactuals once it acts, because the certified emission and the harmful emissions are different $z\_c$'s. The label-choice lesson of \cite{obermeyer2019dissecting} indicates how this ends if uncorrected: the proxy that is cheapest to measure becomes the de facto fairness standard, while the disparity migrates to the unmeasured variable.

\textbf{Why the agent harnesses are not demographically instrumented.} The four capability harnesses at the foot of Table\textasciitilde\{\}2 prove that trajectory-executing evaluation infrastructure exists (reproducible environments, deterministic resets, full action logging \cite{zhou2024webarena,liu2024agentbench,deng2023mind2web,jimenez2024swe}), which makes their fairness silence diagnostic rather than incidental. Three design commitments explain it. Their environments are \emph{demographically sparse by construction}: self-hosted websites, shells, databases, and GitHub issues contain few persons at all, so there is no protected attribute to swap. Their success metrics are \emph{functional and scalar}: task completion and issue resolution are demographic-blind by definition, and nothing in the harness records who was affected by an action, only whether the action achieved the goal. Their tasks are \emph{non-allocative}: even a perfectly instrumented demographic swap would find no consequential decision to flip in "book the cheapest flight" or "fix this failing test." Retrofitting them is therefore not a matter of adding a metric but of adding \emph{paired worlds}: environment states that are identical under $\textbackslash\{\}sigma$ in every record, page, and tool response, plus seeded determinism so that trace-replay attributes divergence to the swap rather than to sampling. Those are the engineering requirements of \S{}4.2's protocol. In P5's terms, the capability harnesses achieve execution adequacy without measurement adequacy, while every fairness resource above them achieves the reverse. The two halves of the needed benchmark exist, fully built, in literatures that do not cite each other.

\textbf{The conclusion, stated through the framework.} The Agentic? column of Table\textasciitilde\{\}2 contains no "yes." This is not an accident of sampling. The inventory spans the QA-bias lineage, the trustworthiness suites, the clinical and decision tiers, the counterfactual tools, the multi-turn and multi-agent partials, and the capability harnesses: it is the structural consequence of Propositions P2 and P5. By P2 (Masking), endpoint parity on any QA- or generation-level resource is consistent with arbitrarily large offsetting component disparities, so every "No" row can certify as fair an agent whose interior $\textbackslash\{\}Delta\_c$ terms are large. By P5 (Measurement adequacy), certification of an acting agent requires bounding each $\textbackslash\{\}Delta\_c$ with nonzero conduction weight, which demands exactly the action-level instruments of \S{}4.3 running inside exactly the paired-world environments the capability harnesses almost provide. Read jointly with the coverage matrix of \S{}7 (Figure\textasciitilde\{\}2) (where counterfactual fairness is the named framing in exactly one component-tagged corpus paper), the gap is two-sided: the metric dimension agents most need is the least applied, and the benchmark layer that would apply it does not exist. The audit protocol of \S{}6 is specified to demonstrate, on a small scale, that this instrumentation is buildable with current tooling; the benchmark agenda of \S{}8.1 specifies what a full-scale version requires.



\subsection*{5. How to Mitigate: Proposition P4 and the Matching Principle}

The Bias Conduction Framework articulates \textbf{Proposition P4 (Mitigation Matching)}: an intervention reduces trajectory-level disparity $\Delta$\textbackslash\{\}\_$\tau$ only if it acts \emph{at or upstream of} the entry locus producing $\Delta$\textbackslash\{\}\_c \textbf{and} on the active conduction operator $\varphi$ at that locus. When an intervention operates downstream of the entry locus, it may suppress the observable symptom (the final text token) while the operator continues to conduct disparity forward into subsequent pipeline stages. When it acts at the wrong operator type (e.g., targeting ATTENUATE dynamics with an intervention calibrated for AMPLIFY), it may be ineffective or, in the case of feedback-recurrence paths, counter-productive. This section surveys the mitigation literature through the P4 lens, asks at which component and operator type each intervention acts, and identifies the gaps that remain. Table 3 formalizes the mapping. Table 4 closes the \S{}3$\rightarrow$\S{}5$\rightarrow$\S{}8 loop by specifying, for each component, the coverage of existing mitigation evidence and the critical open gap.

Three preliminary observations structure the review. First, the overwhelming majority of published interventions act on the base language model's token distribution (prompting, alignment, guardrails, inference-time steering) and are therefore upstream of the agentic components (tool selection, memory, delegation, planning, personalization, drift) where Section 3 locates most of the harm. P4 predicts this mismatch systematically: a model-level debiasing that reduces group disparity in a QA setting neither reaches the tool-selector's discrete routing policy nor the retrieval ranker's document distribution. Second, the few interventions that do act at the agent level (multi-agent debiasing, RAG-layer controls, memory hygiene, identity anonymization) are recent, sparse, and evaluated in narrow settings; the (AMPLIFY, GENERATE) operator cells of Table 3 are empty by evidence, not by design. Third, compound pipeline architectures create a constraint that is easy to underestimate: even a component-matched intervention at locus \emph{c} can fail if an upstream GENERATE operator has already inserted disparity on zero input, since the mitigation target did not exist at the pre-entry stage. P4 therefore calls for matching locus and operator and for auditing the full conduction chain to locate the \emph{first} non-zero $\Delta$\textbackslash\{\}\_c. The intervention must reach that earliest point.

The token-level debiasing literature (intrinsic LLM debiasing at the embedding, pre-training, or decoding level) is the scope of the parent survey \cite{chu2024fairness} and is treated here in a single compressed paragraph. The agentic contribution of this section is the mapping to operator type, component locus, and pipeline stage, not a rehearsal of the QA-era literature.

\subsubsection*{5.1 Prompting and In-Context Debiasing}

Prompting interventions work by restructuring the information presented to the model at inference time, without modifying weights. Under the BCF, they act at the \emph{input pre-processing stage}, prior to all six agent components, and are best characterized as ATTENUATE operators: they reduce the salience of protected-attribute signals in the model's context before the first agentic decision is made. Two causal prompting approaches from the corpus share this structure and permit useful comparative analysis; a third extends the approach to example selection.

\cite{li2024prompting} and \cite{zhang2024causal} both operationalize front-door adjustment at the prompt level, structuring the chain of reasoning to block spurious demographic pathways from reaching the model's causal graph. \cite{furniturewala2024thinking} achieves ATTENUATion through a different mechanism: multi-step deliberative templates that force explicit reasoning before a demographic-adjacent judgment. The key empirical difference between these approaches is where on the reasoning trace the attenuation is applied. Causal prompting blocks the path at the \emph{input encoding} step, before chain-of-thought begins; structured deliberation ("thinking fair and slow") blocks it at an \emph{intermediate reasoning gate}, tolerating early demographic activation but redirecting it before commitment to an action. The \cite{furniturewala2024thinking} finding that structured prompts outperform zero-shot fairness instructions is consistent with the BCF prediction: a flat instruction competes with the model's learned associations, whereas a multi-step structure enforces ATTENUATion at a reasoning node that precedes the action commit. \cite{du2024causal} extends the causal approach into few-shot example selection, choosing in-context demonstrations by causal influence score to block demographic confounds, which attenuates bias \emph{in the information selected} rather than in the prompt structure itself.

Despite their accessibility, prompting interventions have a structural ceiling defined by P4. They act upstream of every agent component, but only on the PRESERVE/ATTENUATE operators of the base model's reasoning; they cannot suppress the AMPLIFY operator introduced when a tool registry presents demographically unequal options (\S{}3.1) or the GENERATE operator active when a retrieval index returns a biased document set for a demographically marked query (\S{}3.2). An agent that reformulates the system prompt between turns (a common architecture pattern in ReAct-style loops \cite{yao2023react}) further degrades prompt-level guarantees because each reformulation resets the context structure, potentially re-introducing the blocked demographic pathway. Prompting is therefore a viable first layer, deployable with no infrastructure change, but P4 marks it as insufficient at any component locus downstream of the model's initial token prediction.

\subsubsection*{5.2 Fine-Tuning and Alignment}

Training-time interventions shape the model's weight distribution to reduce the propensity to generate demographically differentiated outputs, targeting the PRESERVE$\rightarrow$ATTENUATE transformation at the model level. A comparative reading of the alignment literature reveals a structural progression. Standard RLHF \cite{ouyang2022training} is the enabling foundation, but it does not target fairness; it targets human-preference alignment, which \cite{xiao2024algorithmic} shows can itself \emph{generate} minority-group disparity through preference collapse, a mechanism where majority-group preferences dominate the reward signal and the model drifts toward AMPLIFY behavior for minority demographics. MaxMin-RLHF \cite{chakraborty2024maxmin} directly corrects this by replacing the maximization objective with a maximin criterion, ensuring the worst-off preference subgroup is optimized. That fairness-at-training-time intervention acts on the same AMPLIFY-by-preference-collapse pathway that standard RLHF inadvertently activates. BiasDPO \cite{allam2024biasdpo} applies the Direct Preference Optimization alternative to RLHF, constructing biased/debiased preference pairs to steer away from discriminatory completions. Reasoning-guided fine-tuning \cite{kabra2025reasoning} incorporates fairness-aware chain-of-thought supervision so that the model learns to reason through demographic parity as part of task execution, an approach particularly relevant for planning-stage (\S{}3.4) decomposition tasks where an agent allocates subtasks with differential treatment.

Constitutional AI \cite{bai2022constitutional} and its collective variant \cite{huang2024collective} occupy a distinct tier within alignment: rather than optimizing against a reward model that encodes human preferences, they install a set of normative principles against which the model self-critiques. The fairness implication is significant. Standard RLHF can encode the biases of its annotator pool; a constitution-governed model can articulate fairness constraints explicitly, providing a mechanism for ATTENUATION of group-differential behavior that does not depend on the demographic representativeness of preference data. \cite{huang2024collective} strengthens this by deriving the constitution through public deliberation, which reduces the probability that the normative constraints themselves encode majority-group values.

The shared limitation of all training-time methods under P4 is that they act on the model's weight distribution, not on the agent's action-selection architecture. Once the aligned base model is embedded in a tool-using or multi-agent pipeline, the alignment signal does not propagate to: (a) the tool-selection policy, which is a discrete routing choice conditioning on tool metadata that the reward model never observed; (b) the retrieval-stage embedding that ranks candidate documents; or (c) the role-assignment protocol in a multi-agent orchestrator. The aggregate reward signal, summed across a full trajectory, cannot distinguish between a bias introduced at the planning step and one introduced at tool invocation; targeted correction of specific pipeline loci is invisible at training time. This is the BCF's component-locality principle (P1): a weight-level intervention has $\Delta$\textbackslash\{\}\_c = 0 at every non-model locus.

\subsubsection*{5.3 Guardrails and Filters}

Guardrails interpose a protective layer between agent outputs and their consumers, operating at the inference-time boundary between components. \cite{dong2024building} provides the systematic taxonomy of this design space, distinguishing input filters (blocking demographic-sensitive inputs before they reach the model), output filters (screening completions post-generation), and architectural wrappers (modular safety classifiers surrounding the core model). For agent deployment, the natural extension is \emph{action-level guardrails}: before an agent executes a tool call or commits a delegation decision, a classifier audits the proposed action for discriminatory routing.

Under the BCF, guardrails act as post-hoc ATTENUATE operators: they do not reduce the upstream disparity generation but intercept it before the output commits. Their P4 status depends critically on \emph{where in the pipeline} they are inserted. A guardrail at the final output stage satisfies neither the "at or upstream of the entry locus" requirement nor the operator-matching requirement for components like retrieval (where the AMPLIFY/GENERATE operator is active \emph{inside} the retrieval step, not at its output). Guardrails become more P4-compliant when inserted between pipeline components (at the tool-invocation boundary, at the delegation boundary, at memory write) rather than only at the response boundary. In principle, this component-interposed architecture is achievable: the same classifier logic that \cite{dong2024building} applies to conversational outputs can be applied to action proposals. In practice, no corpus paper reports a guardrail specifically trained on action-level fairness signals; all guardrail deployments in the agent fairness literature re-use classifiers trained on natural-language toxicity or refusal tasks, which are neither calibrated to demographic disparity in tool selection nor able to detect GENERATE-type operator effects (disparity created from zero input when a tool produces a biased output).

The measurement gap documented in \S{}4 is directly implicated here: without action-level fairness training data, training action-level guardrails is not tractable. The guardrail tier is thus a plausible P4-compliant architecture whose evidence base is at present thin.

\subsubsection*{5.4 Multi-Agent Debiasing}

Multi-agent architectures introduce distinct bias pathways (\S{}3.3), but they also create structural opportunities for peer-level bias correction that single-agent systems cannot exploit. The multi-agent debiasing literature is the only mitigation family in Table 3 that acts natively at the agent-architecture level, operating on the delegation protocol and identity-signaling structure rather than on the base model's token distribution.

Two parallel design philosophies emerge from comparative analysis. The first, \emph{diversity ensemble}, assigns multiple agents to generate responses independently, then suppresses idiosyncratic biases through aggregation. \cite{owens2024multi} instantiates this as a general multi-LLM debiasing framework; \cite{xu2024mitigating} instantiates it as a multi-objective optimization in which specialized agents maximize distinct fairness criteria and a coordinator agent reconciles them. Both approaches convert the AMPLIFY operator in multi-agent interaction into an ATTENUATE operator through deliberate heterogeneity: the collective is less biased than any member because divergent biases cancel in the aggregation step. The second philosophy, \emph{identity neutralization}, acts directly on the GENERATE pathway that produces bias from role and identity signals. \cite{choi2025identity} demonstrates this cleanly: anonymizing agent identity cues in a multi-agent debate context reduces identity-driven bias without eliminating the epistemic value of role differentiation. This is a direct P4-compliant intervention: the GENERATE operator that creates disparity from demographic role signals is deactivated at the locus where it acts (the identity-conditioning step of the debate protocol), rather than downstream. \cite{ki2025multiple} extends the debate approach to cultural fairness: structured adversarial debate surfaces and corrects cultural stereotypes more effectively than single-agent self-critique, with the multi-agent argumentative dynamic functioning as an ATTENUATION operator on culturally skewed priors.

At the implicit interaction level, \cite{borah2024implicit} addresses AMPLIFY dynamics that arise through accumulated interaction history. A detection-and-mitigation pipeline analyzes multi-agent interaction logs to identify bias propagation patterns before they crystallize into persistent behavioral tendencies. This temporal framing is P4-relevant because it targets the feedback-recurrence path of the BCF equation (the loop $\Delta$\textasciicircum\{\}(t+1) = $\varphi$\textbackslash\{\}\_fb $\cdot$ $\Delta$\textasciicircum\{\}(t) + $\Delta$\textbackslash\{\}\_entry) by interrupting the recurrence before the accumulated disparity exceeds a threshold.

The collective limitation of multi-agent debiasing is architectural: these interventions require a multi-agent infrastructure to function. They are not portable to single-agent settings, and their efficacy is contingent on the specific topology and communication protocol of the multi-agent system. Crucially, none of the corpus papers provides a general mapping from the multi-agent communication graph structure to the magnitude of the resulting AMPLIFY operator. The conduction factor $\varphi$ at the delegation locus remains empirically unbounded.

\subsubsection*{5.5 Inference-Time Steering}

Inference-time steering operates on the model's internal representations during the forward pass, without modifying weights and without requiring multi-agent infrastructure. This makes it structurally distinct from alignment (which modifies weights) and from prompting (which acts on the input). \cite{li2025fairsteer} operationalizes this as FairSteer, which identifies fairness-relevant directions in the model's activation space and applies dynamic activation steering to suppress those directions at inference time.

Under the BCF, FairSteer is a PRESERVE$\rightarrow$ATTENUATE operator acting on the model's internal representation layer. Its key architectural advantage over prompting is that it operates inside the forward pass rather than on the input conditioning, making it more reliable against prompt reformulation and more precisely targeted at the representation subspace encoding demographic associations. For agent deployment, activation steering can in principle be conditioned on agent-state signals (the current tool call context, the retrieved document set, the assigned role), which would enable P4-compliant component-level targeting: suppress fairness-relevant activations specifically when the agent is in a tool-selection or delegation state, rather than uniformly throughout the trajectory.

This state-conditional extension has not been demonstrated in the corpus, but it represents a plausible path toward inference-time interventions that match the component architecture P4 demands. The current limitation of FairSteer is shared with other representation-level methods: it acts on the base model's encoding and does not reach the discrete routing logic of tool selection, the ranking function of retrieval, or the delegation protocol of a multi-agent orchestrator.

\subsubsection*{5.6 Process-Level, Architectural, and RAG-Level Interventions}

The most architecturally integrated tier of mitigation embeds fairness enforcement directly into the agent's operational components (the retrieval layer, the memory module, the planning trace) rather than applying it as an external correction over the model's outputs. These interventions are the most P4-compliant in the corpus, and they are also the most sparsely evidenced.

\textbf{RAG-layer interventions} address the retrieval component (\S{}3.2), where the AMPLIFY and GENERATE operators are active on the information that conditions all downstream agentic reasoning. \cite{kim2024fair} establishes the theoretical and empirical case that ranking fairness in the retrieval stage propagates directly into downstream generation quality across demographic groups, a component-matched diagnosis that names the entry locus precisely. \cite{kim2025mitigating} translates this into a practical P4-compliant intervention: by acting on the \emph{embedding model} rather than on the ranker or the generator, the approach corrects demographic skew at the source representation layer, reducing disparity before it enters the ranking pipeline. This is the BCF's strongest existing example of locus-matching in the retrieval component. The intervention acts at the GENERATE locus (where the embedding model encodes the demographic signal into the retrieval key) rather than at a downstream symptom. The limitation is that the embedding-level correction addresses demographic skew in the representation of queries and documents but does not prevent adversarial amplification through knowledge-store poisoning \cite{wang2025bias}, which operates at a different point in the pipeline.

\textbf{Memory-level hygiene} addresses \S{}3.6's long-horizon drift by intervening on the agent's persistent state. \cite{ma2026implicit} characterizes how implicit biases accumulate across LLM long-term memory and proposes memory audit and selective forgetting mechanisms. Under the BCF, accumulated memory disparity is a feedback-recurrence path: each turn's $\Delta$\textbackslash\{\}\_c is written into memory and re-read as $\Delta$\textbackslash\{\}\_entry for the next turn, creating the compounding effect P3 (Super-additivity) names. Memory hygiene (auditing which inferences are stored, which are forgotten) acts to reduce $\varphi$\textbackslash\{\}\_fb in the recurrence equation, lowering the amplification factor over the trajectory. No corpus paper has demonstrated memory hygiene at scale in an agentic deployment, making this an architecturally motivated but empirically nascent intervention.

\textbf{Bias-aware retrieval and source selection} are addressed by \cite{singh2025bias} and \cite{singh2025biasb}, which frame the source-selection decision (which knowledge sources to query) as a fairness-critical routing choice analogous to tool selection in \S{}3.1. \cite{singh2025biasb} proposes a bias mitigation \emph{agent} dedicated to optimizing source selection for balanced knowledge retrieval across demographic groups, a design pattern that inserts a specialized fairness-enforcement agent into the tool-use pathway as a gatekeeper. This is architecturally P4-compliant: the intervention acts at the entry locus (retrieval source selection) on the AMPLIFY operator (the source routing that conditions what the agent retrieves).

\textbf{Planning and decision-level fairness rewards} address \S{}3.4 directly. \cite{hall2025guiding} inserts fairness reward models into the planning trace, scoring intermediate reasoning steps against demographic parity criteria, which functions as a component-level guardrail on the planning stage's PRESERVE operator, converting it toward ATTENUATE. This approach is particularly significant because it is the only corpus paper that \emph{scores the plan itself} for fairness rather than scoring the plan's output, which corresponds to D2's $\Delta$\textbackslash\{\}\_c measurement at the planning locus rather than the trajectory endpoint.

\textbf{Summary assessment under P4.} Across all five subsections, the dominant empirical picture is asymmetric: the ATTENUATE and PRESERVE operators at the model level are densely populated with mitigation evidence, while the AMPLIFY and GENERATE operator cells (which correspond to the agentic failure modes: retrieval poisoning, role-induced bias generation, delegation amplification, feedback drift) are either thinly covered or empty. No published intervention jointly addresses tool selection bias and long-horizon drift within a unified architectural framework. Building such a framework, one that treats fairness as an agent-level property enforced across the full component stack, is the central open problem of \S{}8.6.


\textbf{Table 3: Mitigation Map: Entry Locus $\times$ Conduction Operator $\times$ Pipeline Stage $\times$ Evidence}

\emph{Takeaway: Mitigation evidence is concentrated in the model-level ATTENUATE column; the (AMPLIFY, GENERATE) $\times$ agent-component cells are empty, marking the gap between where interventions act and where agentic harm originates.}

\begin{center}\footnotesize
\begin{longtable}{|p{0.15\textwidth}|p{0.15\textwidth}|p{0.15\textwidth}|p{0.15\textwidth}|p{0.15\textwidth}|p{0.15\textwidth}|}
\hline \textbf{Entry Locus} & \textbf{Conduction Operator Targeted} & \textbf{Pipeline Stage} & \textbf{Intervention Class} & \textbf{Representative Evidence} & \textbf{Evidence Density} \\ \hline\endhead
Model / input pre-processing & PRESERVE $\rightarrow$ ATTENUATE & Pre-inference (prompt) & Causal prompting & \cite{li2024prompting}, \cite{zhang2024causal}, \cite{furniturewala2024thinking} & High \\ \hline
Model / input pre-processing & PRESERVE $\rightarrow$ ATTENUATE & Pre-inference (few-shot) & Causal active learning & \cite{du2024causal} & Medium \\ \hline
Model (weights) & PRESERVE $\rightarrow$ ATTENUATE & Training (RLHF) & Alignment / MaxMin-RLHF & \cite{ouyang2022training}, \cite{chakraborty2024maxmin}, \cite{xiao2024algorithmic} & High \\ \hline
Model (weights) & PRESERVE $\rightarrow$ ATTENUATE & Training (DPO) & BiasDPO / preference optimization & \cite{allam2024biasdpo} & Medium \\ \hline
Model (weights) & PRESERVE $\rightarrow$ ATTENUATE & Training (SFT + CoT) & Reasoning-guided fine-tuning & \cite{kabra2025reasoning} & Medium \\ \hline
Model (weights) & PRESERVE $\rightarrow$ ATTENUATE & Training (constitutional) & Constitutional / Collective AI & \cite{bai2022constitutional}, \cite{huang2024collective} & Medium \\ \hline
Model (activations) & PRESERVE $\rightarrow$ ATTENUATE & Inference (activation) & FairSteer activation steering & \cite{li2025fairsteer} & Low \\ \hline
All output boundaries & ATTENUATE (post-hoc filter) & Inference (output gate) & Guardrails & \cite{dong2024building} & Low \\ \hline
Tool / API Selection (\S{}3.1) & ATTENUATE (source selection) & Retrieval routing & Bias-aware retrieval agent & \cite{singh2025bias}, \cite{singh2025biasb} & Low \\ \hline
Memory \& Retrieval (\S{}3.2) & GENERATE $\rightarrow$ ATTENUATE (embedder) & Retrieval embedding & RAG fairness / embedder control & \cite{kim2024fair}, \cite{kim2025mitigating} & Low \\ \hline
Memory \& Retrieval (\S{}3.2) & AMPLIFY $\rightarrow$ ATTENUATE (memory audit) & Long-term memory & Memory hygiene / selective forgetting & \cite{ma2026implicit} & Very Low \\ \hline
Multi-agent Delegation (\S{}3.3) & AMPLIFY $\rightarrow$ ATTENUATE (ensemble) & Agent coordination & Multi-LLM debiasing ensemble & \cite{owens2024multi}, \cite{xu2024mitigating} & Low \\ \hline
Multi-agent Delegation (\S{}3.3) & GENERATE $\rightarrow$ ATTENUATE (identity) & Agent debate protocol & Identity anonymization & \cite{choi2025identity} & Low \\ \hline
Multi-agent Delegation (\S{}3.3) & AMPLIFY $\rightarrow$ ATTENUATE (debate) & Agent debate protocol & Cultural debate debiasing & \cite{ki2025multiple} & Low \\ \hline
Multi-agent Delegation (\S{}3.3) & AMPLIFY $\rightarrow$ ATTENUATE (recurrence) & Interaction history & Implicit bias detection pipeline & \cite{borah2024implicit} & Low \\ \hline
Planning / Decomposition (\S{}3.4) & PRESERVE $\rightarrow$ ATTENUATE (plan scoring) & Planning trace & Fairness reward model on plan & \cite{hall2025guiding} & Very Low \\ \hline
Planning / Decomposition (\S{}3.4) & PRESERVE $\rightarrow$ ATTENUATE (reasoning) & Fine-tuning trace & Reasoning-guided SFT & \cite{kabra2025reasoning} & Very Low \\ \hline
Tool / API Selection: AMPLIFY & n/a & n/a & \textbf{(No published intervention)} & n/a & \textbf{Empty} \\ \hline
Tool / API Selection: GENERATE & n/a & n/a & \textbf{(No published intervention)} & n/a & \textbf{Empty} \\ \hline
Memory \& Retrieval: AMPLIFY (poisoning) & n/a & n/a & \textbf{(No published intervention)} & n/a & \textbf{Empty} \\ \hline
User Modeling / Personalization (\S{}3.5) & n/a & n/a & \textbf{(No published intervention at locus)} & n/a & \textbf{Empty} \\ \hline
Long-horizon Drift (\S{}3.6): AMPLIFY (feedback loop) & n/a & n/a & \textbf{(No published intervention at locus)} & n/a & \textbf{Empty} \\ \hline
\end{longtable}\end{center}


\textbf{Table 4: Challenges $\rightarrow$ Solutions $\rightarrow$ Gaps: Closing the \S{}3 $\rightarrow$ \S{}5 $\rightarrow$ \S{}8 Loop}

\emph{Takeaway: Five of the six agent components have partial or no mitigation coverage; the most consequential operator types (AMPLIFY and GENERATE at agent loci) are the least addressed, and every component's open gap maps directly onto a \S{}8 research agenda item.}

\begin{center}\footnotesize
\begin{longtable}{|p{0.18\textwidth}|p{0.18\textwidth}|p{0.18\textwidth}|p{0.18\textwidth}|p{0.18\textwidth}|}
\hline \textbf{Component (\S{}3.x)} & \textbf{Primary Harm (BCF operator)} & \textbf{Mitigation Evidence} & \textbf{Coverage} & \textbf{Critical Open Gap ($\rightarrow$ \S{}8)} \\ \hline\endhead
\textbf{Tool / API Selection} (\S{}3.1) & Discriminatory routing via option-order sensitivity, metadata-driven selection skew; PRESERVE or AMPLIFY on zero user-level input \cite{blankenstein2025biasbusters,sneh2025tooltweak} & Source-selection bias mitigation agents \cite{singh2025bias,singh2025biasb}; bias-aware retrieval routing & \textbf{Partial} (selection fairness addressed at source level; no intervention on discrete tool-routing policy during execution) & No action-level metric or training signal for tool-invocation disparity across demographic counterfactuals; AMPLIFY/GENERATE cells empty ($\rightarrow$ \S{}8.1, \S{}8.6) \\ \hline
\textbf{Memory \& Retrieval} (\S{}3.2) & RAG injects demographic skew even for fairness-vigilant users; GENERATE at embedding and AMPLIFY over multi-turn history \cite{hu2024free,fan2024fairmt,wang2025bias} & Embedder-level bias control \cite{kim2025mitigating}; fair RAG ranking \cite{kim2024fair}; memory audit \cite{ma2026implicit} & \textbf{Partial} (embedder and ranking addressed; knowledge-store poisoning and multi-turn accumulation unaddressed) & No joint intervention on embedding-level GENERATE and history-level AMPLIFY within a unified RAG architecture; no longitudinal retrieval fairness benchmark ($\rightarrow$ \S{}8.1, \S{}8.2) \\ \hline
\textbf{Multi-agent Delegation} (\S{}3.3) & Emergent and amplified bias from role assignment and agent interaction; AMPLIFY at interaction edges, GENERATE at identity-conditioned role seeding \cite{nguyen2025social,madigan2025emergent} & Ensemble debiasing \cite{owens2024multi,xu2024mitigating}; identity anonymization \cite{choi2025identity}; cultural debate \cite{ki2025multiple}; implicit bias monitoring \cite{borah2024implicit} & \textbf{Partial} (several targeted interventions; none covers arbitrary topologies or quantifies the conduction factor $\varphi$ at delegation edges) & No metric for "assignment disparity" across demographic counterfactuals; amplification factor unbounded; no intervention on delegation-graph topology as a fairness variable ($\rightarrow$ \S{}8.3) \\ \hline
\textbf{Planning / Decomposition} (\S{}3.4) & Reasoning-stage demographic bias; plan steps, resource allocation, and task assignments skewed; PRESERVE or AMPLIFY within chain-of-thought \cite{li2025actions,parziale2026once,wu2025reasoning} & Fairness reward on planning trace \cite{hall2025guiding}; reasoning-guided SFT \cite{kabra2025reasoning} & \textbf{Partial} (planning-trace scoring and reasoning-supervised training demonstrated; neither tested in tool-using agentic loop) & No benchmark that scores the \emph{plan itself} for demographic skew independently of the final decision; plan-level disparity metric absent ($\rightarrow$ \S{}8.1) \\ \hline
\textbf{User Modeling / Personalization} (\S{}3.5) & Stereotype activation via inferred identity (names, dialect, usage history); differential quality across demographic groups; GENERATE from inferred attribute \cite{kantharuban2024stereotype,pawar2025presumed,mire2025rejected} & Reward-model dialect-bias mitigation \cite{mire2025rejected} (indirect); no direct user-modeling-locus interventions & \textbf{None} (no corpus paper intervenes at the agent's user-modeling or personalization component) & No mechanism to detect when an agent's inferred user model encodes a protected attribute; no intervention on persistent cross-session user representation; cue-channel inconsistency \cite{tonneau2026different} unaddressed in any debiasing design ($\rightarrow$ \S{}8.1, \S{}8.4) \\ \hline
\textbf{Long-horizon Drift} (\S{}3.6) & Bias accumulation over trajectories via sycophancy and feedback loops; AMPLIFY in recurrence equation; $\varphi$\textbackslash\{\}\_fb > 1 drives compounding \cite{ma2026implicit,wyllie2024fairness,liu2025truth} & Memory hygiene \cite{ma2026implicit} (partial, not validated at scale); sequential fairness theory \cite{xu2024adapting,hu2022achieving} (formal, not yet applied to LLM agents) & \textbf{None} (no deployed mitigation that enforces a long-term fairness criterion over an LLM agent trajectory) & No LLM agent evaluated against a formal long-term fairness criterion (equal long-term benefit rate, etc.); sycophancy-driven drift measurable \cite{hong2025measuring} but not mitigated in agentic settings ($\rightarrow$ \S{}8.2) \\ \hline
\end{longtable}\end{center}



\subsection*{7. Coverage Map and Cross-Analysis}

This section projects the 201-paper corpus onto its two organizing axes simultaneously: the six agent-pipeline components and the four fairness dimensions (group, individual, counterfactual, and unspecified). The projection is Figure 2 (the coverage matrix) and its companion, Figure 3 (the locus $\times$ BCF-operator grid). Read together, they convert the survey's central claim (that the field is well-instrumented where measurement is easy and under-instrumented where agentic deployment makes fairness consequential) from an assertion into a structured, falsifiable coordinate.

\textbf{Tagging convention.} The corpus comprises \textbf{201 unique papers}. Of these, \textbf{57 papers} are tagged to at least one of the six agent-pipeline component branches in \texttt{corpus\_master.md}. Because a single paper may address more than one component (e.g., a study of RAG-induced bias in a multi-agent pipeline is tagged to both memory-retrieval and multi-agent delegation), the 57 unique component-tagged papers generate \textbf{68 component-branch entries} across the coverage matrix, more entries than papers. All counts in the narrative below follow this multi-tag convention. The remaining 144 papers (foundations, llm-agents-general, evaluation-metrics-benchmarks, mitigation, high-stakes-domains, and prior-art-additions branches) are covered in \S{}\S{}2-5 as contextual and cross-cutting evidence; they inform the taxonomy but do not appear in the six-component rows of Figure 2.

\subsubsection*{7.1 The Counterfactual Column is Essentially Empty}

The headline of Figure 2 is a near-vacant column. Across all six component rows, formal counterfactual fairness is the named measurement framing in \textbf{exactly one paper}: \cite{she2025fairsense}, under long-horizon drift, which simulates demographic-conditioned ML-system trajectories over time. The remaining five components (tool/API selection, memory and retrieval, multi-agent delegation, planning/decomposition, and user modeling/personalization) each register \textbf{zero} counterfactual entries.

This is a structural inversion of need. The counterfactual flip (re-running the same agent on a profile altered only in a protected attribute and asking whether the trajectory changed) is the natural operationalization of the agentic fairness question: "would the agent have \emph{acted} differently for another demographic?" Group-level metrics, computed over output distributions, cannot directly answer this. Counterfactual flip can. In BCF terms, the counterfactual swap is the procedure that produces D2's $\Delta$\_c, the component counterfactual disparity that enters the conduction equation. The dimension that BCF-based auditing most requires is precisely the dimension the field has least applied.

The gap is not one of missing methods but of missing application. The counterfactual tooling exists: CFR/MASD in \cite{mayilvaghanan2026counterfactual}; the CAFFE framework for systematic counterfactual evaluation \cite{parziale2025systematic}; the SCOPE stereotyped-prompt dataset \cite{parziale2026scope}; the equal-access counterfactual audit \cite{amirimargavi2026equal}; and intervention-consistency testing \cite{basu2026names}. All five sit in the evaluation-metrics-benchmarks branch; none has been wired through any of the six agent-pipeline components as their primary measurement instrument. The gap is one of application, not invention.

\emph{Takeaway (Figure 2): Across all six agent components, counterfactual fairness (the dimension BCF-based auditing requires for $\Delta$\_c estimation) appears in exactly one of the 68 component-branch entries; GENERATE and AMPLIFY conduction signatures remain empirically uncharted.}

\subsubsection*{7.2 Group Fairness Clusters Only Where Decisions are Explicit}

Where a formal fairness dimension is invoked, it is overwhelmingly group fairness, and it clusters in the two components where the agent emits an explicit, attributable decision: \textbf{multi-agent delegation} (5 group-framed papers, 1 individual) and \textbf{planning/decomposition} (4 group, 1 individual). These are the components whose outputs most resemble the binary decisions that group-parity metrics were designed for: a role assignment, a debate verdict, a task allocation, a discrete approval or denial. Multi-agent work measures disparity across agent personas and emergent group outcomes \cite{madigan2025emergent,nguyen2025social,li2025from,mirza2025malibu}; planning work measures decision and allocation disparity \cite{li2025actions,hall2025guiding,parziale2026once}. The pattern is structural: \textbf{formalization tracks decision explicitness, not harm severity}. Components that expose a clean decision variable attract group metrics; components whose harm is diffuse (retrieval ranking, memory accumulation, user-model inference) do not, regardless of how consequential those harms may be.

\subsubsection*{7.3 Tool, Memory, and Personalization Harms are Documented but Rarely Formalized}

The complementary pattern is a heavily weighted unspecified column. Three components are dominated by it almost entirely: \textbf{memory and retrieval} (0 formal, 12 unspecified), \textbf{tool/API selection} (1 group, 11 unspecified), and \textbf{user modeling/personalization} (1 group, 11 unspecified). Memory and retrieval is the starkest case: not a single paper in this branch names a formal fairness dimension, despite a substantial body documenting real harm: RAG degrading fairness even for vigilant users \cite{hu2024free}, multi-turn fairness erosion \cite{fan2024fairmt}, belief drift from accumulating context \cite{geng2025accumulating,simhi2026old}, and adversarial retrieval-store poisoning \cite{wang2025bias}. These papers demonstrate disparity convincingly; they do not measure it against group, individual, or counterfactual criteria, instead reporting bespoke or qualitative gaps. The same holds for tool selection \cite{blankenstein2025biasbusters,sneh2025tooltweak} and personalization \cite{kantharuban2024stereotype,pawar2025presumed,mire2025rejected,tonneau2026different}.

Harm that lives outside the formal-fairness vocabulary cannot be benchmarked, tracked across model versions, or compared across systems. Documentation without formalization is the field's modal state for half the agent pipeline, and it is the direct driver of the P5 Measurement-adequacy gap the BCF identifies.

\subsubsection*{7.4 Figure 3: Locus $\times$ BCF-Operator Grid}

Figure 2 cross-tabulates component with fairness \emph{dimension}; Figure 3 cross-tabulates component-as-locus with BCF conduction \emph{operator} (ATTENUATE / PRESERVE / AMPLIFY / GENERATE). The two visualizations are complementary: Figure 2 reveals what formal vocabulary is used; Figure 3 reveals what \emph{transformation} of disparity is being empirically studied versus merely assumed.

\textbf{Construction.} Each corpus paper with a component tag is further assigned a conduction operator based on whether it documents disparity attenuation (a debiasing intervention that reduces $\Delta$\_c at that component), preservation (disparity passes through unchanged), amplification (disparity increases), or generation (disparity arises at that component from a nominally zero input, the agentic phenomenon with no single-model analogue). Papers are tagged by the operator their primary finding instantiates; a paper can carry multiple tags.

\begin{center}\footnotesize
\begin{longtable}{|p{0.15\textwidth}|p{0.15\textwidth}|p{0.15\textwidth}|p{0.15\textwidth}|p{0.15\textwidth}|p{0.15\textwidth}|}
\hline \textbf{Component \textbackslash\{\} Operator} & \textbf{ATTENUATE $\varphi$<1} & \textbf{PRESERVE $\varphi$$\approx$1} & \textbf{AMPLIFY $\varphi$>1} & \textbf{GENERATE $\varphi$ (zero$\rightarrow$nonzero)} & \textbf{Uncharacterized} \\ \hline\endhead
Tool/API selection & 1 & 3 & 1 & 1 & 6 \\ \hline
Memory \& retrieval & 2 & 2 & 3 & 1 & 4 \\ \hline
Multi-agent delegation & 3 & 1 & 4 & 3 & 1 \\ \hline
Planning/decomposition & 2 & 2 & 2 & 0 & 4 \\ \hline
User modeling/personalization & 1 & 4 & 3 & 1 & 3 \\ \hline
Long-horizon drift & 2 & 1 & 4 & 1 & 2 \\ \hline
\end{longtable}\end{center}

\emph{Takeaway (Figure 3): GENERATE and AMPLIFY cells (the operators that produce distinctly agentic harm, where pipeline edges produce disparity that was absent at the input or that grows across the trajectory) are the least instrumented; the AMPLIFY/GENERATE column for planning/decomposition and the GENERATE column for tool selection are nearly uncharacterized, pointing to the highest-priority measurement gaps for BCF-guided empirical work.}

The most consequential observation is that the \textbf{GENERATE} operator (disparity arising at a pipeline edge from nominally unbiased inputs, the formal signature of the agentic multiplier in P3) has been documented in only a handful of papers. For multi-agent delegation, \cite{nguyen2025social} and \cite{madigan2025emergent} provide the clearest evidence of generation (bias emerging in interaction among individually measured-to-be-unbiased agents); for memory, \cite{ma2026implicit} documents generation in the accumulation process; for tool selection, \cite{blankenstein2025biasbusters} documents generation from tool-description surface features alone. Outside these, the GENERATE column is largely empty. This is the empirical frontier the audit of \S{}6 is designed to probe.

\subsubsection*{7.5 Synthesis and Differentiation}

Four regularities fall out of the combined Figure 2/Figure 3 analysis.

First, \textbf{formalization tracks decision explicitness, not harm severity}: components that emit a clean decision variable attract formal metrics; diffuse-harm components do not.

Second, \textbf{the formal vocabulary thins toward the counterfactual and toward GENERATE/AMPLIFY} exactly as the evaluation question shifts from "what did the agent output?" toward the distinctively agentic questions: "would it have acted differently?" and "did the pipeline produce disparity it did not inherit?"

Third, \textbf{the modal corpus paper documents disparity without naming a fairness dimension or a conduction operator}: 52 of the 68 component-branch entries sit in the unspecified column of Figure 2, and most of the locus $\times$ operator cells in Figure 3 are under-populated. The field is evidence-rich and formalization-poor.

Fourth, the BCF coordinate system makes the gaps precise and actionable. The empty cells of Figure 3 are specific (component, operator) pairs that can be populated by targeted experiments with the counterfactual-perturbation machinery already available \cite{parziale2025systematic,mayilvaghanan2026counterfactual,basu2026names}.

This analysis also sharpens the differentiation from neighboring work (Table 0 in \S{}1.4). Where \cite{chu2024fairness} organizes single-turn LLM fairness by metric and \cite{gallegos2024bias} by datasets and mitigation, neither resolves \emph{where in an acting pipeline} a disparity originates. Where \cite{mohammadi2025evaluation} treats fairness as one sub-topic of agent evaluation and \cite{vatsal2026agentic} as one of seven clinical dimensions, neither makes component coverage its unit of analysis. Where \cite{ranjan2025fairness} proposes a framework for fairness in classical multi-agent systems, it does not operate at the LLM-agent component level or develop the conduction-operator analysis. By projecting the corpus onto the component $\times$ operator grid, this survey turns the measurement gap from an assertion into a structured coordinate (the empty counterfactual column and the under-populated GENERATE/AMPLIFY cells) and hands \S{}8 a concrete, BCF-keyed research target.


\subsection*{8. Open Problems and Research Agenda}

The coverage analysis of \S{}7 converts the measurement gap from an observation into a set of precise coordinates: (component, operator) cells in Figure 3 that are empirically uncharacterized, P5 Measurement-adequacy cells that no existing benchmark populates, and BCF propositions whose empirical grounding is weak. Each problem below opens with its BCF coordinate, states why it matters for consequential deployment, and specifies a concrete first step (a metric definition, a dataset structure, or a system design) that an immediate follow-on could execute.

A governance thread runs through all six problems. Agents are now deployed in hiring, lending, healthcare, and public-benefits adjudication, domains where disparate treatment across race, gender, socioeconomic status, dialect, and their intersections carries direct legal and regulatory weight. The EU AI Act (Regulation 2024/1689) classifies high-risk AI systems in employment, credit, and essential-service sectors as subject to bias-testing and transparency requirements \cite{parliament2024regulation}; the NIST AI Risk Management Framework specifies bias evaluation as a component of AI governance \cite{nist2023artificial}; New York City Local Law 144 requires automated employment decision tools to undergo annual bias audits before deployment; and EEOC guidance applies Title VII disparate-impact standards to algorithmic hiring tools. None of these frameworks currently specifies what a fairness audit of an \emph{acting agent} must cover, and the measurement infrastructure described in \S{}4.4 is not adequate to discharge those requirements. Every open problem below has a governance correlate: the missing metric, dataset, or system design is also the missing regulatory instrument.

\subsubsection*{8.1 Unpopulated P5 Cells: Benchmarks for Acting Agents}

\textbf{BCF coordinate.} P5 Measurement-adequacy states that a metric is adequate for component c only if it can detect non-zero $\Delta$\_c at that component independently of all other components. The benchmark inventory in \S{}4.4 shows that no existing benchmark satisfies P5 for any of the six agent-pipeline components at the action level. Figure 3's uncharacterized cells are exactly the P5 cells that no existing instrument populates.

\textbf{The problem.} The dominant fairness benchmarks score answers, not actions. Clinical fairness sets define the ceiling of current practice: FairMedQA evaluates demographic bias in medical QA \cite{xiao2025fairmedqa}; mFARM assembles multi-faceted harm assessment for clinical decision support on static vignettes \cite{adappanavar2025mfarm}; MedEqualQA probes consistency through counterfactual demographic edits to QA questions \cite{ghosh2025medequalqa}; and the health equity toolbox surfaces harms at the level of generated text \cite{pfohl2024toolbox}. Each is methodologically careful; each stops at the QA boundary. None measures what a clinical acting agent does when it selects which guideline tool to invoke, which prior note to retrieve, whether to escalate to a human, or how it sequences a multi-step workup. The few studies that score actions confirm the divergence: agent decisions reveal biases absent from token-level probes \cite{li2025actions}, and decomposing unfairness in a tool-using chest X-ray agent attributes $\Delta$\_c to specific pipeline stages invisible to answer-level scoring \cite{xu2026ducx}. Existing agent evaluation frameworks (AgentBench \cite{liu2024agentbench}, WebArena \cite{zhou2024webarena}, SWE-bench \cite{jimenez2024swe}) are not demographically instrumented; they evaluate capability (task completion) rather than fairness across demographic counterfactuals, so they cannot be repurposed directly without structural additions.

\textbf{Why it matters.} A system can be near-parity on FairMedQA and still route, escalate, or order differently across demographic counterfactuals once it acts, because $\Delta$\_c enters through tool routing and retrieval rather than the final string. Certifying agents with QA benchmarks licenses deployment on evidence that does not bind the deployed behavior, precisely the accountability gap EU AI Act Article 10 (data quality) and NIST AI RMF Measure 2.5 (bias testing) are meant to close.

\textbf{Concrete first step.} A minimal acting-agent fairness benchmark requires: (i) a set of consequential-task environments (e.g., clinical triage, hiring screening, loan underwriting) in which an agent can take multiple discrete actions; (ii) counterfactual profile pairs that differ only in a demographic signal, following the CAFFE/SCOPE design \cite{parziale2025systematic,parziale2026scope}; (iii) logging of all intermediate actions (tool invocations, retrieval queries, routing decisions) per counterfactual pair; and (iv) reporting of per-edge $\Delta$\_c alongside the final-verdict CFR/MASD. FairMedAgent, described in \S{}6.4, represents one instantiation for the clinical domain; comparable instruments are needed for hiring and lending, where legal audit requirements are already in force.

\subsubsection*{8.2 Unmeasured $\varphi$\_fb Recurrence: Long-Horizon and Compounding Fairness}

\textbf{BCF coordinate.} The feedback-recurrence term $\Delta$\textasciicircum\{\}(t+1) = $\varphi$\_fb $\cdot$ $\Delta$\textasciicircum\{\}(t) + $\Delta$\_entry formalizes how per-step disparity compounds across time when an agent conditions on its own prior outputs. P3 Super-additivity, applied longitudinally, predicts that $\varphi$\_fb > 1 produces exponential divergence of group outcomes across deployment cycles, a prediction no existing evaluation captures.

\textbf{The problem.} Fairness is evaluated as a one-shot property, but agents persist. Implicit bias in long-term memory accumulates and propagates, reinforcing small initial skews across interactions \cite{ma2026implicit}. Sequential decision-making theory establishes the formal stakes: static per-step fairness constraints fail to deliver equitable long-term outcomes, and equal long-term benefit requires reasoning about how present decisions reshape future eligibility \cite{hu2022achieving,xu2024adapting}. Feedback loops amplify this further: training on a model's own synthetic outputs amplifies bias across generations \cite{wyllie2024fairness}, and self-consuming performative loops compound the effect at deployment \cite{wang2026observations}. The theory and the LLM-drift empirics \cite{fan2024fairmt,ma2026implicit,liu2025truth} have not been joined: no study evaluates an LLM agent against a formal long-term fairness criterion over a realistic multi-session deployment.

\textbf{Concrete first step.} Define a \textbf{longitudinal disparity index} (LDI) as the ratio of group-outcome divergence at turn T to turn 1, estimated by simulation over a synthetic trajectory of T interactions using the recurrence $\Delta$\textasciicircum\{\}(t+1) = $\varphi$\_fb $\cdot$ $\Delta$\textasciicircum\{\}(t) + $\Delta$\_entry. Parameterize $\varphi$\_fb empirically from an existing deployed agent using the FairSense simulation tooling \cite{she2025fairsense}. A single-domain pilot (e.g., the loan-triage replication arm of \S{}6) with T = 20 turns and five demographic group pairs would constitute a first binding measurement.

\subsubsection*{8.3 Unbounded GENERATE: Multi-Agent Fairness Dynamics}

\textbf{BCF coordinate.} The GENERATE operator ($\varphi$ applied to nominally zero-$\Delta$\_c input) is the signature failure mode of multi-agent systems: disparity arises in the interaction among agents that are individually measured as unbiased. Figure 3 shows the GENERATE column is the least populated in the multi-agent delegation row, despite the strongest theoretical and empirical motivation.

\textbf{The problem.} Multi-agent orchestration introduces failure modes with no single-agent analogue: bias can \emph{emerge} from interaction even among agents whose individual outputs are measured as near-parity. \cite{madigan2025emergent} formalizes emergent bias in multi-agent decision systems as a property of the collective rather than any constituent; \cite{nguyen2025social} documents bias emerging, propagating, and amplifying through agent-to-agent communication. Disparity in these settings is a function of communication topology, role-assignment structure, and message-passing dynamics. None of these is observable by a per-agent audit, and none is instrumented by existing benchmarks. The MALIBU benchmark \cite{mirza2025malibu} surfaces implicit bias in multi-agent assignments but does not measure the GENERATE operator specifically (bias arising at zero input to each constituent).

\textbf{Concrete first step.} Design a \textbf{GENERATE-detection probe}: a multi-agent task in which each constituent agent is individually certified as CFR = 0 (i.e., no flip under demographic swap) on a single-turn benchmark, then assembled into a two-agent delegation. Measure the CFR of the delegation. Any non-zero CFR at the delegation level constitutes GENERATE evidence; the difference between delegation-level and max(constituent) CFR is an empirical estimate of the GENERATE amplitude for that topology. The probe requires only the counterfactual evaluation machinery of \cite{mayilvaghanan2026counterfactual} and any existing two-agent framework \cite{wu2023autogen}.

\subsubsection*{8.4 Intersectional and Sociotechnical Harms in Deployment}

\textbf{BCF coordinate.} BCF D2 defines $\Delta$\_c as a counterfactual disparity over a single protected attribute; real deployment harms are intersectional (jointly over multiple attributes) and sociotechnical (conditioned on deployment context, not just model parameters). This is a D2 extension problem: the current formalization is a first-order approximation.

\textbf{The problem.} Bias research frequently isolates one protected attribute on a clean benchmark, but real deployments harm people at intersections. The largest-scale evidence is unambiguous: sociodemographic biases in LLM medical decision-making produce differential management recommendations at clinical scale, with race, gender, and socioeconomic status interacting \cite{omar2025sociodemographic}; a mortgage underwriting experiment documents racial disparities in lending \cite{bowen2024measuring}; and judicial fairness evaluation finds systematic differential treatment across defendant demographics \cite{hu2025llms}. These harms are jointly demographic and contextual, depending on the deployment's stakes, decision structure, and the combination of attributes a real person carries. Intersectional treatment is specifically flagged in algorithmic hiring regulation \cite{raghavan2020mitigating}.

\textbf{Concrete first step.} Extend D2 to an intersectional counterfactual disparity $\textbackslash\{\}Delta\_c\textasciicircum\{\}\{A \textbackslash\{\}times B\}$ defined as the $\Delta$\_c estimated under joint swap of attributes A and B minus the sum of the marginal $\Delta$\_c values, measuring the interaction term. Apply this to the \S{}6 audit's hiring task with a 2$\times$2 race-by-gender counterfactual design; this requires only doubling the profile set and is immediately feasible.

\subsubsection*{8.5 Empty Operator Rows of Figure 3: Governance and Standardized Auditing}

\textbf{BCF coordinate.} The ATTENUATE rows of Figure 3 are partially populated (mitigation papers), but the formal connection between a mitigation's claimed ATTENUATE effect and its BCF-level $\varphi$ value is never made: no paper reports how much $\varphi$\_e decreases at a specific pipeline edge following a given intervention. This is a P4 Mitigation-matching gap.

\textbf{The problem.} There is no shared standard for what constitutes a fairness audit of an LLM agent, nor a required reporting format for action-level disparity. The ingredients exist: the CAFFE counterfactual framework \cite{parziale2025systematic}, the SCOPE dataset \cite{parziale2026scope}, the LangFair evaluation toolkit \cite{bouchard2024bring}, the intervention-consistency test \cite{basu2026names}, and the dataset survey \cite{zhang2025datasets}. No governance instrument standardizes them into a repeatable, comparable audit of an \emph{acting} system. The regulatory field has outpaced the measurement infrastructure: the EU AI Act \cite{parliament2024regulation} requires bias testing of high-risk employment and credit systems, the NIST AI RMF \cite{nist2023artificial} specifies bias evaluation as a governance requirement, NYC Local Law 144 mandates annual bias audits for automated employment decision tools, and EEOC disparate-impact analysis applies to algorithmic hiring. None specifies what those audits must cover for an acting agent as opposed to a static classifier.

\textbf{Concrete first step.} Draft a \textbf{minimum audit specification} for acting agents with four components: (i) a protocol for constructing counterfactual profile sets (building on \cite{parziale2025systematic}); (ii) a required metric set covering, at minimum, per-component CFR, tool-invocation disparity, and escalation-rate disparity, aligned with the BCF's D2 definition; (iii) a standardized reporting template with per-component $\Delta$\_c values and the conduction-operator classification (ATTENUATE/PRESERVE/AMPLIFY/GENERATE) for each active edge; and (iv) a recourse specification covering what information must be logged to enable post-hoc attribution and contestation. This specification can be drafted as a technical standard proposal and submitted for comment to the NIST AI 100-1 framework revision process.

\subsubsection*{8.6 Mitigation at the Agent Level, Not the Token Level}

\textbf{BCF coordinate.} P4 Mitigation-matching states that an intervention is matched if it operates at the same pipeline locus as the disparity it is meant to correct. The mitigation map in \S{}5 shows that the ATTENUATE column of Figure 3 is populated almost entirely by token-level interventions (prompting, alignment, output filtering), which are mismatched to the six agent-pipeline loci where the AMPLIFY and GENERATE operators are documented.

\textbf{The problem.} The mitigation literature overwhelmingly intervenes on model outputs: debiasing prompts \cite{furniturewala2024thinking,li2024prompting}, preference optimization \cite{allam2024biasdpo,bai2022constitutional}, inference-time activation steering \cite{li2025fairsteer}, and output guardrails \cite{dong2024building}. These operate on tokens. But \S{}3 locates agentic harm in actions and architecture: which tool is chosen, what is retrieved into context, how roles are assigned, how a plan is decomposed. Token-level debiasing leaves these channels untouched. The few process-aware interventions point the way: balanced-retrieval agents \cite{singh2025biasb}, bias-aware retrieval agents \cite{singh2025bias}, identity anonymization to neutralize role-induced debate bias \cite{choi2025identity}, and multi-agent debiasing frameworks \cite{owens2024multi,xu2024mitigating}. They are scattered relative to the dominant approach, and no unified framework simultaneously addresses tool-selection, memory-accumulation, and delegation loci.

\textbf{Concrete first step.} Implement a \textbf{component-keyed guardrail layer} as a proof of concept: a wrapper that intercepts the agent's action proposal at each of the three highest-priority AMPLIFY/GENERATE loci identified in Figure 3 (tool selection, multi-agent delegation, long-horizon accumulation) and applies a targeted $\Delta$\_c reduction intervention before action execution. The three interventions are, respectively: a demographic-blind re-ranking of tool candidates, a role-assignment audit against an assignment-disparity threshold, and a memory-hygiene filter. Evaluate the wrapper using the \S{}6 audit's C1 to C3 configurations, reporting per-edge $\varphi$\_e before and after to directly test P4 Mitigation-matching.

\subsubsection*{Synthesis}

The autonomy that makes agents valuable (acting, persisting, delegating, self-conditioning) is precisely what current QA-level measurement, single-shot certification, and token-level mitigation cannot observe. Closing the gap requires instruments, protocols, and interventions defined over trajectories and actions, governed at the system level, and reported in a standardized format auditors and regulators can use. The BCF provides the coordinate system; the six open problems above specify the cells to fill.


\subsection*{9. Conclusion}

LLM-based agents have crossed a consequential threshold: they select tools, accumulate memories across turns, delegate subtasks across role-assigned subagents, decompose goals into multi-step plans, adapt to inferred user identities, and persist behavior over long horizons \cite{wang2024survey,xi2023rise}. Each of these pipeline stages is a distinct locus at which demographic disparity can enter, propagate, and compound into consequential harm. The Bias Conduction Framework (BCF) introduced in \S{}3.0 formalizes this structure: D1 defines the agent loop $\langle$$\pi$, T, M, R$\rangle$; D2 defines component counterfactual disparity $\Delta$\_c via trace-replay; D3 defines the conduction operators (ATTENUATE $\varphi$<1, PRESERVE $\varphi$$\approx$1, AMPLIFY $\varphi$>1, GENERATE $\varphi$ on zero input); the conduction equation $\Delta$\_$\tau$ $\approx$ $\Sigma$\_c $\Delta$\_c $\cdot$ $\Pi$ $\varphi$\_e with feedback recurrence $\Delta$\textasciicircum\{\}(t+1) = $\varphi$\_fb $\cdot$ $\Delta$\textasciicircum\{\}(t) + $\Delta$\_entry derives the trajectory disparity from its per-component contributions; and five propositions (P1 Locality, P2 Masking, P3 Super-additivity, P4 Mitigation-matching, P5 Measurement-adequacy) characterize the conditions under which component-level auditing is necessary, sufficient, and policy-relevant.

This survey has presented the first dedicated taxonomy of fairness in LLM-based agents organized by this component axis, spanning tool and API selection (\S{}3.1), memory and retrieval (\S{}3.2), multi-agent delegation and role assignment (\S{}3.3), planning and decomposition (\S{}3.4), user modeling and personalization (\S{}3.5), and long-horizon trajectory drift (\S{}3.6), with a cross-cutting account of how component-level harms compound into the agentic multiplier (\S{}3.7, P3). Against this architecture the survey maps the evaluation field (\S{}4), characterizing group, individual, and counterfactual metric families lifted to agent trajectories and generalizing the CFR and MASD statistics of \cite{mayilvaghanan2026counterfactual} across all pipeline components, while Table 2 documents that nearly every existing benchmark (clinical fairness sets \cite{xiao2025fairmedqa,adappanavar2025mfarm,ghosh2025medequalqa}, hiring and lending audits \cite{an2025measuring,azime2025accept}, compositional evaluation suites \cite{wang2024ceb}) scores model outputs rather than agent actions. The mitigation survey (\S{}5) follows the P4 Mitigation-matching principle, distinguishing prompt-level interventions \cite{furniturewala2024thinking,li2024prompting}, alignment-stage approaches \cite{allam2024biasdpo,bai2022constitutional}, architectural guardrails \cite{dong2024building}, and the emerging class of process-level and multi-agent debiasing strategies \cite{owens2024multi,choi2025identity,ki2025multiple}. The coverage analysis (\S{}7) quantifies the gap: 52 of 68 component-branch entries sit in the unspecified column of Figure 2; the counterfactual column (the dimension required to compute D2's $\Delta$\_c) contains exactly one entry across all six components; and the GENERATE/AMPLIFY cells of Figure 3 are the least populated despite being the loci of distinctively agentic harm. The cross-framework audit (\S{}6) is registered as a protocol to test P2 and P3 empirically; results are pending.

The measurement gap documented by this survey is a gap in accountability infrastructure, not in tooling alone. Agents are deployed in hiring \cite{veldanda2023emily,wilson2024gender,nghiem2024you}, clinical decision support \cite{poulain2024bias,omar2025sociodemographic}, lending \cite{bowen2024measuring,feng2023empowering}, and education \cite{weissburg2024llms}, subject to regulatory requirements (EU AI Act \cite{parliament2024regulation}, NIST AI RMF \cite{nist2023artificial}, NYC Local Law 144, EEOC disparate-impact guidance) that no current acting-agent fairness benchmark is equipped to satisfy. The long-term fairness risks are compounding: feedback recurrence documented by \cite{wyllie2024fairness} and \cite{wang2026observations}, memory accumulation characterized by \cite{ma2026implicit}, and the sequential fairness theory of \cite{hu2022achieving} and \cite{xu2024adapting} collectively predict that systems fair at any single step can drive divergent cumulative outcomes across groups over deployment cycles.

The immediate priority is benchmark construction for acting agents: evaluation infrastructure that scores trajectories and decisions rather than tokens and outputs, applies counterfactual demographic perturbations at the action level \cite{parziale2025systematic,amirimargavi2026equal,basu2026names}, and covers the per-component, per-edge $\Delta$\_c granularity the BCF specifies. That infrastructure is the precondition for every subsequent advance in mitigation (P4), governance (\S{}8.5), and longitudinal auditing (\S{}8.2). The field has the conceptual vocabulary, the BCF framework, and the first components of the counterfactual tooling. The gap is instrumentation.


\subsection*{9.1 Ethical Considerations, Societal Impact, and Limitations of this Survey}

Surveys of bias and fairness in AI systems carry their own ethical charge: the taxonomy and evaluation framework this work advances can be used both to \emph{audit and constrain} unfair agent behavior and to \emph{reverse-engineer and evade} such constraints. We address this dual-use concern and several additional limitations that bear on how the survey's findings should be interpreted and acted upon.

\subsubsection*{9.1.1 Dual-Use of the Taxonomy and BCF Framework}

The BCF's conduction-operator vocabulary (identifying which pipeline edges are AMPLIFY or GENERATE loci, and which mitigation interventions reduce $\varphi$\_e at each edge) is designed as an auditor's instrument. The same information can be used adversarially: knowing that tool-description metadata is a GENERATE locus \cite{sneh2025tooltweak} could guide an actor who wishes to steer an agent's routing by manipulating tool descriptions rather than inputs. By the same logic, the fact that memory-accumulation is an AMPLIFY locus could inform the design of retrieval-store poisoning attacks \cite{wang2025bias} calibrated to exploit that amplification. The survey discloses these mechanisms because responsible auditing requires understanding them; the disclosure is not an operational guide for exploitation. We recommend that practitioners who use the BCF framework as an audit template also consult it as an adversarial-risk checklist, evaluating deployed agents not only for inadvertent bias but also for susceptibility to targeted disparity injection at the loci this taxonomy identifies.

\subsubsection*{9.1.2 Academic-vs-Legal Fairness Gap}

The fairness criteria formalized in \S{}2.1 and operationalized throughout the survey (demographic parity, equalized odds, counterfactual fairness) are academic constructs. Their relationship to the legal standards that govern high-stakes AI deployment is partial and contested. Disparate-impact doctrine under Title VII (applied to hiring), the Equal Credit Opportunity Act (lending), and the EU AI Act's non-discrimination provisions each impose distinct evidentiary standards, burden-of-proof structures, and remediation requirements that academic fairness metrics do not automatically satisfy. In particular: (i) a system achieving statistical demographic parity may still constitute unlawful disparate treatment under intent-based theories; (ii) a system failing the counterfactual flip test may nonetheless satisfy legal requirements if the disparity is justified by a legitimate business necessity; and (iii) the granular per-component $\Delta$\_c measurements the BCF calls for are not currently required by any regulatory framework, which means deployers may satisfy legal compliance while exhibiting large action-level disparities invisible to regulators. This gap reflects a genuine normative distance between what the research community can measure and what policymakers have chosen to require; it is not a flaw in academic fairness research. We encourage legal scholars and regulators to engage with the BCF's measurement proposals as inputs to the standard-setting processes now underway under the EU AI Act and NIST AI RMF frameworks.

\subsubsection*{9.1.3 English- and US/EU-Centric Corpus Positionality}

The 201-paper corpus was assembled from English-language academic venues (arXiv, ACL Anthology, ACM DL, NeurIPS/ICML/ICLR, FAccT, Nature), all of which have geographic and linguistic concentrations in the United States and Western Europe. Several structural biases follow from this:

\begin{itemize}
\item \textbf{Language bias}: nearly all benchmark datasets and evaluation prompts are in English; dialect bias findings (e.g., \cite{mire2025rejected} on African American Language; \cite{lin2024assessing} on non-standard dialects) are primarily US-context findings. Bias patterns in Arabic, Hindi, Swahili, and other high-speaker-count languages are substantially underrepresented.
\item \textbf{Protected-attribute scope}: the protected attributes most frequently studied (race (primarily Black/white US framing), binary gender, and US-relevant socioeconomic status) do not map cleanly onto the protected categories, group identities, and fairness norms of other legal and cultural contexts. India's caste system, Brazil's color classification, the EU's national-origin protections, and indigenous identity frameworks are effectively absent from the corpus.
\item \textbf{High-stakes-domain coverage}: the corpus's consequential-domain evidence clusters in US hiring (name-based resume audits), US lending (mortgage underwriting), and US clinical settings (medical QA). Non-US regulatory contexts, public-benefit adjudication systems, and welfare-state service delivery are not covered.
\end{itemize}

These positionality constraints mean the survey's findings are most directly applicable to LLM-based agents deployed in English-language, US- or EU-regulated, employment/credit/healthcare contexts, and should be extrapolated to other contexts with caution. Geographic and multilingual expansion of the corpus is the highest-priority scope extension for a follow-up survey.

\subsubsection*{9.1.4 Limitations Arising from the Unrun Audit}

The cross-framework audit of \S{}6 is pre-registered but not yet executed. As a result:

\begin{itemize}
\item The BCF propositions P2 and P3 are supported by theoretical derivation and by the existing multi-paper evidence cited throughout \S{}\S{}3-4, but are not yet supported by the within-survey experimental results the audit would provide. Readers should treat P2 and P3 as well-motivated hypotheses with strong indirect support, not as empirically confirmed claims.
\item The conduction-operator tagging in Figure 3 is based on the framing and findings of the cited papers as reported; it reflects the survey authors' interpretation of whether a given paper's finding instantiates ATTENUATE, PRESERVE, AMPLIFY, or GENERATE at the relevant locus, not a standardized classification applied by the original authors. Some cells may be reclassified as the audit methodology is refined.
\item The audit's action-level metrics (tool-invocation disparity, escalation disparity, plan-allocation disparity) are specified in \S{}6.2 but have not been computed on any dataset. Their empirical behavior (distributional properties, sensitivity, test-retest reliability) is unknown until the run is complete.
\end{itemize}

These limitations are intentional: the survey's contribution is the framework and the protocol, not the numerical results. The protocol is released as a replication target; results will be incorporated in the camera-ready version.

\subsubsection*{9.1.5 Fast-Moving Field and Recency Constraints}

The corpus cutoff is mid-2026. The LLM-agent field's publication rate has accelerated sharply since 2023 \cite{wang2024survey}, and the fairness subfield is following the same trajectory. Several developments that may alter the picture drawn here are already visible at the corpus boundary:

\begin{itemize}
\item \textbf{Frontier model capability shifts}: improvements in instruction-following fidelity and tool-calling reliability may change the empirical profile of which pipeline edges are AMPLIFY versus PRESERVE loci, as models become more or less sensitive to demographic cues at each stage.
\item \textbf{Regulatory enactments}: the EU AI Act's implementing acts specifying audit requirements for high-risk AI systems are expected to issue in 2025 to 2026; these may operationalize bias-testing requirements in ways that alter the alignment between academic fairness metrics and legal compliance requirements.
\item \textbf{Benchmark releases}: AgentBench \cite{liu2024agentbench}, WebArena \cite{zhou2024webarena}, and $\tau$-bench are actively adding evaluation dimensions; it is possible that demographically instrumented agentic fairness evaluation will be released in these frameworks within the survey's publication window, which would partially populate the empty cells of Figure 2.
\end{itemize}

Readers should treat the coverage matrix as a snapshot of the field as of mid-2026, not as a permanent characterization. The BCF framework and the open-problem agenda of \S{}8 are designed to remain relevant as the empirical field changes; the specific cell counts in Figure 2 and Figure 3 should be updated as new work appears.

\begin{thebibliography}{999}
\bibitem{wang2024survey} Wang et al.. (2024). A Survey on Large Language Model based Autonomous Agents. Frontiers of Computer Science. arXiv:2308.11432. doi:10.1007/s11704-024-40231-1
\bibitem{xi2023rise} Xi et al.. (2023). The Rise and Potential of Large Language Model Based Agents: A Survey. arXiv preprint. arXiv:2309.07864
\bibitem{yao2023react} Yao et al.. (2023). ReAct: Synergizing Reasoning and Acting in Language Models. ICLR 2023. arXiv:2210.03629
\bibitem{sumers2024cognitive} Sumers et al.. (2024). Cognitive Architectures for Language Agents. Transactions on Machine Learning Research (TMLR). arXiv:2309.02427
\bibitem{wu2023autogen} Wu et al.. (2023). AutoGen: Enabling Next-Gen LLM Applications via Multi-Agent Conversation. arXiv preprint (ICLR 2024 workshop). arXiv:2308.08155
\bibitem{park2023generative} Park et al.. (2023). Generative Agents: Interactive Simulacra of Human Behavior. UIST 2023. arXiv:2304.03442. doi:10.1145/3586183.3606763
\bibitem{ma2026implicit} Ma et al.. (2026). How Implicit Bias Accumulates and Propagates in LLM Long-term Memory. arXiv preprint. arXiv:2602.01558
\bibitem{nguyen2025social} Nguyen et al.. (2025). The Social Cost of Intelligence: Emergence, Propagation, and Amplification of Stereotypical Bias in Multi-Agent Systems. arXiv preprint. arXiv:2510.10943
\bibitem{an2025measuring} An et al.. (2025). Measuring gender and racial biases in large language models: Intersectional evidence from automated resume evaluation. PNAS Nexus. arXiv:2403.15281. doi:10.1093/pnasnexus/pgaf089
\bibitem{veldanda2023emily} Veldanda et al.. (2023). Are Emily and Greg Still More Employable than Lakisha and Jamal? Investigating Algorithmic Hiring Bias in the Era of ChatGPT. EMNLP 2024 (arXiv preprint 2023). arXiv:2310.05135
\bibitem{nghiem2024you} Nghiem et al.. (2024). "You Gotta be a Doctor, Lin": An Investigation of Name-Based Bias of Large Language Models in Employment Recommendations. EMNLP 2024. arXiv:2406.12232
\bibitem{bowen2024measuring} Bowen et al.. (2024). Measuring and Mitigating Racial Disparities in LLMs: Evidence from a Mortgage Underwriting Experiment. SSRN Working Paper. doi:10.2139/ssrn.4812158
\bibitem{feng2023empowering} Feng et al.. (2023). Empowering Many, Biasing a Few: Generalist Credit Scoring through Large Language Models. arXiv. arXiv:2310.00566
\bibitem{azime2025accept} Azime et al.. (2025). Accept or Deny? Evaluating LLM Fairness and Performance in Loan Approval across Table-to-Text Serialization Approaches. arXiv preprint. arXiv:2508.21512
\bibitem{omar2025sociodemographic} Omar et al.. (2025). Sociodemographic biases in medical decision making by large language models. Nature Medicine. doi:10.1038/s41591-025-03626-6
\bibitem{zack2024coding} Zack et al.. (2024). Coding Inequity: Assessing GPT-4's Potential for Perpetuating Racial and Gender Biases in Healthcare. The Lancet Digital Health. doi:10.1016/S2589-7500(23)00225-X
\bibitem{poulain2024bias} Poulain et al.. (2024). Bias patterns in the application of LLMs for clinical decision support: A comprehensive study. arXiv. arXiv:2404.15149
\bibitem{raghavan2020mitigating} Raghavan et al.. (2020). Mitigating Bias in Algorithmic Hiring: Evaluating Claims and Practices. ACM FAccT. doi:10.1145/3351095.3372828
\bibitem{basu2026names} Basu et al.. (2026). When Names Change Verdicts: Intervention Consistency Reveals Systematic Bias in LLM Decision-Making. arXiv preprint. arXiv:2603.18530
\bibitem{hu2025llms} Hu et al.. (2025). LLMs on Trial: Evaluating Judicial Fairness for Large Language Models. arXiv. arXiv:2507.10852
\bibitem{blankenstein2025biasbusters} Blankenstein et al.. (2025). BiasBusters: Uncovering and Mitigating Tool Selection Bias in Large Language Models. ArXiv preprint. arXiv:2510.00307
\bibitem{wu2024rag} Wu et al.. (2024). Does RAG Introduce Unfairness in LLMs? Evaluating Fairness in Retrieval-Augmented Generation Systems. ArXiv preprint. arXiv:2409.19804
\bibitem{hu2024free} Hu et al.. (2024). No Free Lunch: Retrieval-Augmented Generation Undermines Fairness in LLMs, Even for Vigilant Users. arXiv preprint. arXiv:2410.07589
\bibitem{gupta2023bias} Gupta et al.. (2023). Bias Runs Deep: Implicit Reasoning Biases in Persona-Assigned LLMs. ICLR 2024 (submitted Nov 2023). arXiv:2311.04892
\bibitem{cao2026from} Cao et al.. (2026). From Biased Chatbots to Biased Agents: Examining Role Assignment Effects on LLM Agent Robustness. AAAI 2026 TrustAgent Workshop. arXiv:2602.12285
\bibitem{liu2025truth} Liu et al.. (2025). Truth Decay: Quantifying Multi-Turn Sycophancy in Language Models. arXiv preprint. arXiv:2503.11656
\bibitem{li2025actions} Li et al.. (2025). Actions Speak Louder than Words: Agent Decisions Reveal Implicit Biases in Language Models. arXiv preprint. arXiv:2501.17420. doi:10.48550/arXiv.2501.17420
\bibitem{selbst2019fairness} Selbst et al.. (2019). Fairness and Abstraction in Sociotechnical Systems. ACM FAccT. doi:10.1145/3287560.3287598
\bibitem{blodgett2020language} Blodgett et al.. (2020). Language (Technology) is Power: A Critical Survey of Bias in NLP. ACL. arXiv:2005.14050
\bibitem{kleinberg2017inherent} Kleinberg et al.. (2017). Inherent Trade-Offs in the Fair Determination of Risk Scores. ITCS. arXiv:1609.05807
\bibitem{chouldechova2017fair} Chouldechova. (2017). Fair Prediction with Disparate Impact: A Study of Bias in Recidivism Prediction Instruments. Big Data. arXiv:1703.00056. doi:10.1089/big.2016.0047
\bibitem{mayilvaghanan2026counterfactual} Mayilvaghanan et al.. (2026). Counterfactual Fairness Evaluation of LLM-Based Contact Center Agent Quality Assurance System. arXiv preprint. arXiv:2602.14970
\bibitem{chu2024fairness} Chu et al.. (2024). Fairness in Large Language Models: A Taxonomic Survey. ACM SIGKDD Explorations Newsletter. arXiv:2404.01349. doi:10.1145/3682112.3682117
\bibitem{dwork2012fairness} Dwork et al.. (2012). Fairness Through Awareness. Proceedings of the 3rd Innovations in Theoretical Computer Science Conference (ITCS). arXiv:1104.3913. doi:10.1145/2090236.2090255
\bibitem{hardt2016equality} Hardt et al.. (2016). Equality of Opportunity in Supervised Learning. Advances in Neural Information Processing Systems (NeurIPS). arXiv:1610.02413
\bibitem{kusner2017counterfactual} Kusner et al.. (2017). Counterfactual Fairness. Advances in Neural Information Processing Systems (NeurIPS). arXiv:1703.06856
\bibitem{gallegos2024bias} Gallegos et al.. (2024). Bias and Fairness in Large Language Models: A Survey. Computational Linguistics. arXiv:2309.00770
\bibitem{mehrabi2021survey} Mehrabi et al.. (2021). A Survey on Bias and Fairness in Machine Learning. ACM Computing Surveys. arXiv:1908.09635. doi:10.1145/3457607
\bibitem{pessach2022review} Dana Pessach, Eytan Shmueli. (2022). A Review on Fairness in Machine Learning. ACM Computing Surveys, vol. 55, no. 3, Article 51. doi:10.1145/3494672
\bibitem{caton2024fairness} Simon Caton, Christian Haas. (2024). Fairness in Machine Learning: A Survey. ACM Computing Surveys, vol. 56, no. 7, Article 144. arXiv:2010.04053. doi:10.1145/3616865
\bibitem{mohammadi2025evaluation} Mohammadi et al.. (2025). Evaluation and Benchmarking of LLM Agents: A Survey. arXiv preprint. arXiv:2507.21504
\bibitem{vatsal2026agentic} Vatsal et al.. (2026). Agentic AI in Healthcare \textbackslash\{\}\& Medicine: A Seven-Dimensional Taxonomy for Empirical Evaluation of LLM-based Agents. IEEE Access. arXiv:2602.04813
\bibitem{ranjan2025fairness} Rajesh Ranjan, Shailja Gupta, Surya Narayan Singh. (2025). Fairness in Agentic AI: A Unified Framework for Ethical and Equitable Multi-Agent System. arXiv preprint arXiv:2502.07254. arXiv:2502.07254
\bibitem{ebrahimi2025survey} Sana Ebrahimi, Abolfazl Asudeh. (2025). A Survey on Reliability, Transparency, Accountability, and Fairness in LLM-based Multi-Agent Systems through the Responsibility Lens. Preprint (ResearchGate/unpublished), University of Illinois Chicago
\bibitem{binkyte2025interactional} Ruta Binkyte. (2025). Interactional Fairness in LLM Multi-Agent Systems: An Evaluation Framework. AAAI/ACM AIES 2025 (Proceedings, Vol. 8, pp. 457-468). arXiv:2505.12001. doi:10.1609/aies.v8i1.36563
\bibitem{yu2025survey} Yu et al.. (2025). A Survey on Trustworthy LLM Agents: Threats and Countermeasures. ACM SIGKDD 2025. arXiv:2503.09648
\bibitem{feldman2015certifying} Feldman et al.. (2015). Certifying and Removing Disparate Impact. ACM KDD. arXiv:1412.3756. doi:10.1145/2783258.2783311
\bibitem{barocas2023fairness} Barocas, Hardt, and Narayanan. (2023). Fairness and Machine Learning: Limitations and Opportunities. MIT Press
\bibitem{tang2023what} Zeyu Tang, Jiji Zhang, Kun Zhang. (2023). What-is and How-to for Fairness in Machine Learning: A Survey, Reflection, and Perspective. ACM Computing Surveys, vol. 55, no. 13s, Article 299. arXiv:2206.04101. doi:10.1145/3597199
\bibitem{navigli2023biases} Navigli et al.. (2023). Biases in Large Language Models: Origins, Inventory, and Discussion. ACM JDIQ. doi:10.1145/3597307
\bibitem{buolamwini2018gender} Buolamwini and Gebru. (2018). Gender Shades: Intersectional Accuracy Disparities in Commercial Gender Classification. FAccT (PMLR v81)
\bibitem{caliskan2017semantics} Caliskan et al.. (2017). Semantics Derived Automatically from Language Corpora Contain Human-Like Biases. Science. arXiv:1608.07187. doi:10.1126/science.aal4230
\bibitem{zhao2018gender} Zhao et al.. (2018). Gender Bias in Coreference Resolution: Evaluation and Debiasing Methods (WinoBias). NAACL. arXiv:1804.06876
\bibitem{rudinger2018gender} Rudinger et al.. (2018). Gender Bias in Coreference Resolution (WinoGender). NAACL. arXiv:1804.09301
\bibitem{parrish2022bbq} Parrish et al.. (2022). BBQ: A Hand-Built Bias Benchmark for Question Answering. Findings of ACL 2022. arXiv:2110.08193
\bibitem{dhamala2021bold} Dhamala et al.. (2021). BOLD: Dataset and Metrics for Measuring Biases in Open-Ended Language Generation. Proceedings of the 2021 ACM Conference on Fairness, Accountability, and Transparency (FAccT). arXiv:2101.11718. doi:10.1145/3442188.3445924
\bibitem{nadeem2021stereoset} Nadeem et al.. (2021). StereoSet: Measuring Stereotypical Bias in Pretrained Language Models. Proceedings of ACL-IJCNLP 2021. arXiv:2004.09456
\bibitem{nangia2020crows} Nangia et al.. (2020). CrowS-Pairs: A Challenge Dataset for Measuring Social Biases in Masked Language Models. Proceedings of EMNLP 2020. arXiv:2010.00133
\bibitem{smith2022sorry} Smith et al.. (2022). "I'm sorry to hear that": Finding New Biases in Language Models with a Holistic Descriptor Dataset. Proceedings of EMNLP 2022. arXiv:2205.09209
\bibitem{gehman2020realtoxicityprompts} Gehman et al.. (2020). RealToxicityPrompts: Evaluating Neural Toxic Degeneration in Language Models. EMNLP Findings. arXiv:2009.11462
\bibitem{hartvigsen2022toxigen} Hartvigsen et al.. (2022). ToxiGen: A Large-Scale Machine-Generated Dataset for Adversarial and Implicit Hate Speech Detection. ACL. arXiv:2203.09509
\bibitem{wang2023decodingtrust} Wang et al.. (2023). DecodingTrust: A Comprehensive Assessment of Trustworthiness in GPT Models. NeurIPS. arXiv:2306.11698
\bibitem{liang2023holistic} Liang et al.. (2023). Holistic Evaluation of Language Models. TMLR. arXiv:2211.09110
\bibitem{huang2024trustllm} Huang et al.. (2024). TrustLLM: Trustworthiness in Large Language Models. ICML. arXiv:2401.05561
\bibitem{blodgett2021stereotyping} Blodgett et al.. (2021). Stereotyping Norwegian Salmon: An Inventory of Pitfalls in Fairness Benchmark Datasets. ACL-IJCNLP. doi:10.18653/v1/2021.acl-long.81
\bibitem{guo2024large} Guo et al.. (2024). Large Language Model based Multi-Agents: A Survey of Progress and Challenges. IJCAI. arXiv:2402.01680
\bibitem{durante2024agent} Durante et al.. (2024). Agent AI: Surveying the Horizons of Multimodal Interaction. arXiv preprint. arXiv:2401.03568
\bibitem{schick2023toolformer} Schick et al.. (2023). Toolformer: Language Models Can Teach Themselves to Use Tools. NeurIPS 2023. arXiv:2302.04761
\bibitem{lewis2020retrieval} Lewis et al.. (2020). Retrieval-Augmented Generation for Knowledge-Intensive NLP Tasks. NeurIPS 2020. arXiv:2005.11401
\bibitem{hong2024metagpt} Hong et al.. (2024). MetaGPT: Meta Programming for A Multi-Agent Collaborative Framework. ICLR. arXiv:2308.00352
\bibitem{shinn2023reflexion} Shinn et al.. (2023). Reflexion: Language Agents with Verbal Reinforcement Learning. NeurIPS 2023. arXiv:2303.11366
\bibitem{fan2024fairmt} Fan et al.. (2024). FairMT-Bench: Benchmarking Fairness for Multi-turn Dialogue in Conversational LLMs. ICLR 2025 (spotlight). arXiv:2410.19317
\bibitem{wu2025reasoning} Wu et al.. (2025). Does Reasoning Introduce Bias? A Study of Social Bias Evaluation and Mitigation in LLM Reasoning. arXiv preprint. arXiv:2502.15361
\bibitem{kantharuban2024stereotype} Kantharuban et al.. (2024). Stereotype or Personalization? User Identity Biases Chatbot Recommendations. arXiv preprint. arXiv:2410.05613
\bibitem{wyllie2024fairness} Wyllie et al.. (2024). Fairness Feedback Loops: Training on Synthetic Data Amplifies Bias. ACM FAccT 2024. arXiv:2403.07857. doi:10.1145/3630106.3659029
\bibitem{zhou2024webarena} Zhou et al.. (2024). WebArena: A Realistic Web Environment for Building Autonomous Agents. ICLR. arXiv:2307.13854
\bibitem{liu2024agentbench} Liu et al.. (2024). AgentBench: Evaluating LLMs as Agents. ICLR. arXiv:2308.03688
\bibitem{jimenez2024swe} Jimenez et al.. (2024). SWE-bench: Can Language Models Resolve Real-World GitHub Issues?. ICLR. arXiv:2310.06770
\bibitem{deng2023mind2web} Deng et al.. (2023). Mind2Web: Towards a Generalist Agent for the Web. NeurIPS. arXiv:2306.06070
\bibitem{hu2022achieving} Hu and Zhang. (2022). Achieving Long-Term Fairness in Sequential Decision Making. AAAI 2022. arXiv:2204.01819
\bibitem{xu2024adapting} Xu et al.. (2024). Adapting Static Fairness to Sequential Decision-Making: Bias Mitigation Strategies towards Equal Long-term Benefit Rate. ICML 2024. arXiv:2309.03426
\bibitem{zehlike2022fairness} Zehlike et al.. (2022). Fairness in Ranking: A Survey. ACM Computing Surveys, vol. 55, no. 6, Article 118. arXiv:2103.14000. doi:10.1145/3533379
\bibitem{wang2023survey} Yifan Wang, Weizhi Ma, Min Zhang, Yiqun Liu, Shaoping Ma. (2023). A Survey on the Fairness of Recommender Systems. ACM Transactions on Information Systems, vol. 41, no. 2, Article 38. arXiv:2206.03761. doi:10.1145/3547333
\bibitem{moher2009preferred} Moher et al.. (2009). Preferred Reporting Items for Systematic Reviews and Meta-Analyses: The PRISMA Statement. PLoS Medicine. doi:10.1371/journal.pmed.1000097
\bibitem{wang2026observations} Wang et al.. (2026). Observations and Remedies for Large Language Model Bias in Self-Consuming Performative Loop. arXiv preprint. arXiv:2601.05184
\bibitem{obermeyer2019dissecting} Obermeyer et al.. (2019). Dissecting Racial Bias in an Algorithm Used to Manage the Health of Populations. Science. doi:10.1126/science.aax2342
\bibitem{qu2024tool} Qu et al.. (2024). Tool Learning with Large Language Models: A Survey. Frontiers of Computer Science. arXiv:2405.17935
\bibitem{xu2026ducx} Xu et al.. (2026). DUCX: Decomposing Unfairness in Tool-Using Chest X-ray Agents. ArXiv preprint. arXiv:2603.00777
\bibitem{choi2025identity} Choi et al.. (2025). When Identity Skews Debate: Anonymization for Bias-Reduced Multi-Agent Reasoning. ACL 2026. arXiv:2510.07517
\bibitem{kim2025mitigating} Kim et al.. (2025). Mitigating Bias in RAG: Controlling the Embedder. arXiv preprint. arXiv:2502.17390
\bibitem{owens2024multi} Owens et al.. (2024). A Multi-LLM Debiasing Framework. arXiv preprint. arXiv:2409.13884
\bibitem{mire2025rejected} Mire et al.. (2025). Rejected Dialects: Biases Against African American Language in Reward Models. NAACL Findings 2025. arXiv:2502.12858
\bibitem{mansoury2020feedback} Mansoury et al.. (2020). Feedback Loop and Bias Amplification in Recommender Systems. CIKM 2020. arXiv:2007.13019
\bibitem{xu2025quantifying} Xu et al.. (2025). Quantifying Fairness in LLMs Beyond Tokens: A Semantic and Statistical Perspective. arXiv preprint. arXiv:2506.19028
\bibitem{turpin2023language} Turpin et al.. (2023). Language Models Don't Always Say What They Think: Unfaithful Explanations in Chain-of-Thought Prompting. NeurIPS. arXiv:2305.04388
\bibitem{madigan2025emergent} Madigan et al.. (2025). Emergent Bias and Fairness in Multi-Agent Decision Systems. arXiv preprint. arXiv:2512.16433
\bibitem{tonneau2026different} Tonneau et al.. (2026). Different Demographic Cues Yield Inconsistent Conclusions About LLM Personalization and Bias. arXiv preprint. arXiv:2601.18486
\bibitem{singh2025biasb} Singh et al.. (2025). Bias Mitigation Agent: Optimizing Source Selection for Fair and Balanced Knowledge Retrieval. KDD 2025 Agent4IR Workshop. arXiv:2508.18724
\bibitem{dai2024bias} Dai et al.. (2024). Bias and Unfairness in Information Retrieval Systems: New Challenges in the LLM Era. Proceedings of the 30th ACM SIGKDD Conference on Knowledge Discovery and Data Mining. arXiv:2404.11457. doi:10.1145/3637528.3671458
\bibitem{dai2025trustworthy} Sunhao Dai, Changle Qu, Sirui Chen, Xiao Zhang, Jun Xu. (2025). Trustworthy Information Retrieval in the LLM Era: Bias, Unfairness, and Hallucination. Proceedings of the 2025 Annual International ACM SIGIR Conference (SIGIR-AP '25). doi:10.1145/3767695.3769670
\bibitem{parziale2026once} Parziale et al.. (2026). Once Upon a Team: Investigating Bias in LLM-Driven Software Team Composition and Task Allocation. arXiv preprint. arXiv:2601.03857
\bibitem{staab2023memorization} Staab et al.. (2023). Beyond Memorization: Violating Privacy Via Inference with Large Language Models. arXiv preprint (ICLR 2024 submission). arXiv:2310.07298
\bibitem{u2018algorithms} Noble S. U.. (2018). Algorithms of Oppression: How Search Engines Reinforce Racism. NYU Press. doi:10.2307/j.ctt1pwt9w5
\bibitem{makhortykh2021detecting} Makhortykh et al.. (2021). Detecting race and gender bias in visual representation of AI on web search engines. Advances in Bias-Free Computing (Springer, LNCS 12817). arXiv:2106.14072. doi:10.1007/978-3-030-78818-6\_5
\bibitem{pezeshkpour2023large} Pezeshkpour et al.. (2023). Large Language Models Sensitivity to The Order of Options in Multiple-Choice Questions. ArXiv preprint. arXiv:2308.11483
\bibitem{sneh2025tooltweak} Sneh et al.. (2025). ToolTweak: An Attack on Tool Selection in LLM-based Agents. ArXiv preprint. arXiv:2510.02554
\bibitem{tang2026your} Tang et al.. (2026). Is Your LLM-as-a-Recommender Agent Trustable? LLMs' Recommendation is Easily Hacked by Biases (Preferences). ArXiv preprint. arXiv:2603.17417
\bibitem{singh2025bias} Singh et al.. (2025). Bias-Aware Agent: Enhancing Fairness in AI-Driven Knowledge Retrieval. Companion Proceedings of the ACM Web Conference 2025 (WWW '25). arXiv:2503.21237
\bibitem{zhang2025evaluating} Zhang et al.. (2025). Evaluating the Effect of Retrieval Augmentation on Social Biases. arXiv preprint. arXiv:2502.17611
\bibitem{wang2025bias} Wang et al.. (2025). Bias Amplification in RAG: Poisoning Knowledge Retrieval to Steer LLMs. ArXiv preprint. arXiv:2506.11415
\bibitem{geng2025accumulating} Geng et al.. (2025). Accumulating Context Changes the Beliefs of Language Models. arXiv preprint. arXiv:2511.01805
\bibitem{simhi2026old} Simhi et al.. (2026). Old Habits Die Hard: How Conversational History Geometrically Traps LLMs. arXiv preprint. arXiv:2603.03308
\bibitem{liu2023lost} Liu et al.. (2023). Lost in the Middle: How Language Models Use Long Contexts. Transactions of the Association for Computational Linguistics (TACL). arXiv:2307.03172
\bibitem{coppolillo2025unmasking} Coppolillo et al.. (2025). Unmasking Conversational Bias in AI Multiagent Systems. arXiv preprint. arXiv:2501.14844
\bibitem{kim2024fair} Kim and Diaz. (2024). Towards Fair RAG: On the Impact of Fair Ranking in Retrieval-Augmented Generation. arXiv preprint (ACM SIGIR ICTIR 2025). arXiv:2409.11598
\bibitem{qian2024chatdev} Qian et al.. (2024). ChatDev: Communicative Agents for Software Development. ACL. arXiv:2307.07924
\bibitem{deshpande2023toxicity} Deshpande et al.. (2023). Toxicity in ChatGPT: Analyzing Persona-assigned Language Models. EMNLP Findings. arXiv:2304.05335
\bibitem{mirza2025malibu} Mirza et al.. (2025). MALIBU Benchmark: Multi-Agent LLM Implicit Bias Uncovered. arXiv preprint. arXiv:2507.01019
\bibitem{li2025from} Li et al.. (2025). From Single to Societal: Analyzing Persona-Induced Bias in Multi-Agent Interactions. arXiv preprint. arXiv:2511.11789
\bibitem{borah2024implicit} Borah and Mihalcea. (2024). Towards Implicit Bias Detection and Mitigation in Multi-Agent LLM Interactions. arXiv preprint. arXiv:2410.02584
\bibitem{campedelli2024want} Campedelli et al.. (2024). I Want to Break Free! Persuasion and Anti-Social Behavior of LLMs in Multi-Agent Settings with Social Hierarchy. arXiv preprint (EACL 2024 Workshop). arXiv:2410.07109
\bibitem{ki2025multiple} Ki et al.. (2025). Multiple LLM Agents Debate for Equitable Cultural Alignment. ACL 2025 (Oral). arXiv:2505.24671
\bibitem{kambhampati2024large} Kambhampati. (2024). Can Large Language Models Reason and Plan?. Annals of the New York Academy of Sciences. arXiv:2403.04121. doi:10.1111/nyas.15125
\bibitem{zhou2025veracity} Zhou et al.. (2025). Veracity Bias and Beyond: Uncovering LLMs' Hidden Beliefs in Problem-Solving Reasoning. ACL 2025 (Main). arXiv:2505.16128
\bibitem{kabra2025reasoning} Kabra et al.. (2025). Reasoning Towards Fairness: Mitigating Bias in Language Models through Reasoning-Guided Fine-Tuning. arXiv preprint. arXiv:2504.05632
\bibitem{hall2025guiding} Hall et al.. (2025). Guiding LLM Decision-Making with Fairness Reward Models. arXiv preprint. arXiv:2507.11344
\bibitem{santurkar2023whose} Santurkar et al.. (2023). Whose Opinions Do Language Models Reflect?. ICML. arXiv:2303.17548
\bibitem{pawar2025presumed} Pawar et al.. (2025). Presumed Cultural Identity: How Names Shape LLM Responses. arXiv preprint. arXiv:2502.11995
\bibitem{weissburg2024llms} Weissburg et al.. (2024). LLMs are Biased Teachers: Evaluating LLM Bias in Personalized Education. NAACL Findings 2025. arXiv:2410.14012
\bibitem{deldjoo2024cfairllm} Deldjoo et al.. (2024). CFaiRLLM: Consumer Fairness Evaluation in Large-Language Model Recommender System. arXiv preprint. arXiv:2403.05668
\bibitem{deldjoo2024normative} Deldjoo et al.. (2024). A Normative Framework for Benchmarking Consumer Fairness in Large Language Model Recommender Systems. arXiv preprint. arXiv:2405.02219
\bibitem{arzaghi2024understanding} Arzaghi et al.. (2024). Understanding Intrinsic Socioeconomic Biases in Large Language Models. arXiv preprint. arXiv:2405.18662
\bibitem{lin2024assessing} Lin et al.. (2024). Assessing Dialect Fairness and Robustness of Large Language Models in Reasoning Tasks. ACL 2025. arXiv:2410.11005
\bibitem{vijjini2024exploring} Vijjini et al.. (2024). Exploring Safety-Utility Trade-Offs in Personalized Language Models. NAACL 2025. arXiv:2406.11107
\bibitem{sharma2024understanding} Sharma et al.. (2024). Towards Understanding Sycophancy in Language Models. ICLR. arXiv:2310.13548
\bibitem{hong2025measuring} Hong et al.. (2025). Measuring Sycophancy of Language Models in Multi-turn Dialogues. EMNLP 2025 Findings. arXiv:2505.23840
\bibitem{she2025fairsense} She et al.. (2025). FairSense: Long-Term Fairness Analysis of ML-Enabled Systems. ICSE 2025. arXiv:2501.01665. doi:10.1109/ICSE55347.2025.00159
\bibitem{wilson2024gender} Wilson et al.. (2024). Gender, Race, and Intersectional Bias in Resume Screening via Language Model Retrieval. arXiv. arXiv:2407.20371
\bibitem{adappanavar2025mfarm} Adappanavar et al.. (2025). mFARM: Towards Multi-Faceted Fairness Assessment based on HARMs in Clinical Decision Support. arXiv preprint. arXiv:2509.02007
\bibitem{song2025large} Song et al.. (2025). Towards Large Language Models that Benefit for All: Benchmarking Group Fairness in Reward Models. arXiv preprint. arXiv:2503.07806
\bibitem{parziale2025systematic} Parziale et al.. (2025). Toward Systematic Counterfactual Fairness Evaluation of Large Language Models: The CAFFE Framework. ICSE 2026. arXiv:2512.16816
\bibitem{parziale2026scope} Parziale et al.. (2026). SCOPE: A Dataset of Stereotyped Prompts for Counterfactual Fairness Assessment of LLMs. IEEE (ICSE 2026 proceedings). arXiv:2604.05555
\bibitem{amirimargavi2026equal} Amiri-Margavi et al.. (2026). Equal Access, Unequal Interaction: A Counterfactual Audit of LLM Fairness. arXiv preprint. arXiv:2602.02932
\bibitem{huang2023trustgpt} Huang et al.. (2023). TrustGPT: A Benchmark for Trustworthy and Responsible Large Language Models. arXiv preprint. arXiv:2306.11507
\bibitem{xiao2025fairmedqa} Xiao et al.. (2025). FairMedQA: Benchmarking Bias in Large Language Models for Medical Question Answering. arXiv preprint. arXiv:2505.19562
\bibitem{ghosh2025medequalqa} Ghosh et al.. (2025). MEDEQUALQA: Evaluating Biases in LLMs with Counterfactual Reasoning. arXiv preprint. arXiv:2510.12818
\bibitem{wu2024fmbench} Wu et al.. (2024). FMBench: Benchmarking Fairness in Multimodal Large Language Models on Medical Tasks. arXiv preprint (ICLR 2025 workshop). arXiv:2410.01089
\bibitem{pfohl2024toolbox} Pfohl et al.. (2024). A Toolbox for Surfacing Health Equity Harms and Biases in Large Language Models. Nature Medicine. arXiv:2403.12025. doi:10.1038/s41591-024-03258-2
\bibitem{wang2024ceb} Wang et al.. (2024). CEB: Compositional Evaluation Benchmark for Fairness in Large Language Models. ICLR 2025 Spotlight. arXiv:2407.02408
\bibitem{bouchard2024bring} Bouchard, D.. (2024). Bring Your Own Prompts: Use-Case-Specific Bias and Fairness Evaluation for LLMs (LangFair). arXiv preprint. arXiv:2407.10853
\bibitem{zhang2025datasets} Zhang et al.. (2025). Datasets for Fairness in Language Models: An In-Depth Survey. arXiv preprint. arXiv:2506.23411
\bibitem{yehudai2025survey} Asaf Yehudai, Lilach Eden, Alan Li, Guy Uziel, Yilun Zhao, Roy Bar-Haim, Arman Cohan, Michal Shmueli-Scheuer. (2025). A Survey on Evaluation of LLM-based Agents. arXiv preprint arXiv:2503.16416 (revised April 2026). arXiv:2503.16416
\bibitem{chang2024survey} Yupeng Chang, Xu Wang, Jindong Wang, Yuan Wu, Linyi Yang, Kaijie Zhu, Hao Chen, Xiaoyuan Yi, Cunxiang Wang, Yidong Wang, Wei Ye, Yue Zhang, Yi Chang, Philip S. Yu, Qiang Yang, Xing Xie. (2024). A Survey on Evaluation of Large Language Models. ACM Transactions on Intelligent Systems and Technology (TIST), vol. 15, no. 3, Article 39. arXiv:2307.03109. doi:10.1145/3641289
\bibitem{li2024prompting} Li et al.. (2024). Prompting Fairness: Integrating Causality to Debias Large Language Models. ICLR 2025. arXiv:2403.08743
\bibitem{zhang2024causal} Zhang et al.. (2024). Causal Prompting: Debiasing Large Language Model Prompting based on Front-Door Adjustment. AAAI 2025. arXiv:2403.02738
\bibitem{furniturewala2024thinking} Furniturewala et al.. (2024). Thinking Fair and Slow: On the Efficacy of Structured Prompts for Debiasing Language Models. arXiv preprint. arXiv:2405.10431
\bibitem{du2024causal} Du et al.. (2024). Causal-Guided Active Learning for Debiasing Large Language Models. ACL 2024 (Outstanding Paper). arXiv:2408.12942
\bibitem{ouyang2022training} Ouyang et al.. (2022). Training language models to follow instructions with human feedback. NeurIPS 2022. arXiv:2203.02155
\bibitem{xiao2024algorithmic} Xiao et al.. (2024). On the Algorithmic Bias of Aligning Large Language Models with RLHF: Preference Collapse and Matching Regularization. Journal of the American Statistical Association 2025. arXiv:2405.16455
\bibitem{chakraborty2024maxmin} Chakraborty et al.. (2024). MaxMin-RLHF: Alignment with Diverse Human Preferences. ICML 2024. arXiv:2402.08925
\bibitem{allam2024biasdpo} Allam et al.. (2024). BiasDPO: Mitigating Bias in Language Models through Direct Preference Optimization. arXiv preprint. arXiv:2407.13928
\bibitem{bai2022constitutional} Bai et al.. (2022). Constitutional AI: Harmlessness from AI Feedback. arXiv preprint. arXiv:2212.08073
\bibitem{huang2024collective} Huang et al.. (2024). Collective Constitutional AI: Aligning a Language Model with Public Input. FAccT 2024. arXiv:2406.07814
\bibitem{dong2024building} Dong et al.. (2024). Building Guardrails for Large Language Models. arXiv preprint. arXiv:2402.01822
\bibitem{xu2024mitigating} Xu et al.. (2024). Mitigating Social Bias in Large Language Models: A Multi-Objective Approach within a Multi-Agent Framework. AAAI 2025. arXiv:2412.15504
\bibitem{li2025fairsteer} Li et al.. (2025). FairSteer: Inference Time Debiasing for LLMs with Dynamic Activation Steering. ACL 2025. arXiv:2504.14492
\bibitem{parliament2024regulation} European Parliament and Council. (2024). Regulation (EU) 2024/1689 (Artificial Intelligence Act). Official Journal of the EU
\bibitem{nist2023artificial} NIST. (2023). Artificial Intelligence Risk Management Framework (AI RMF 1.0). NIST AI 100-1. doi:10.6028/NIST.AI.100-1
\end{thebibliography}
\end{document}
